\documentclass[12pt, a4paper]{report}

% --- PACKAGES ---
\usepackage{geometry}
\geometry{a4paper, margin=2cm}
\usepackage{fontspec}
\usepackage{xeCJK}
\usepackage{xcolor}
\usepackage{tcolorbox}
\tcbuselibrary{skins, breakable}
\usepackage{parskip}
\usepackage{mwe}
\usepackage[zerostyle=b,scaled=.75]{newtxtt}
\usepackage{booktabs}
\usepackage{listings}

% --- FONT CONFIGURATION ---
% Pastikan font Noto Sans JP sudah terinstall di Mac kamu.
\setmainfont{Arial} % Font utama untuk teks alfabet
\setCJKmainfont[AutoFakeBold=true, AutoFakeSlant=true]{Noto Sans JP}

% --- CUSTOM COLORS ---
\definecolor{kanjicolor}{HTML}{1E3A8A}    % Biru gelap untuk kanji
\definecolor{furicolor}{HTML}{D946EF}     % Magenta untuk furigana
\definecolor{romacolor}{HTML}{EA580C}     % Oranye gelap untuk romaji
\definecolor{articolor}{HTML}{047857}     % Hijau untuk bahasa indonesia
\definecolor{boxbg}{HTML}{FEF9C3}         % Kuning pastel untuk box percakapan
\definecolor{boxframe}{HTML}{F59E0B}      % Oranye untuk border percakapan
\definecolor{lessonbg}{HTML}{F0FDF4}      % Hijau pastel untuk lesson
\definecolor{lessonframe}{HTML}{10B981}   % Hijau terang untuk border lesson

% --- CUSTOM MACROS UNTUK KANJI, FURIGANA, & ROMAJI ---
% \jpword{furigana}{kanji}{romaji}
\newcommand{\jpword}[3]{%
  \begin{tabular}[b]{@{}c@{}}
    {\fontsize{10}{12}\selectfont\color{furicolor}\sffamily #1} \\[-0.2ex]
    {\fontsize{24}{28}\selectfont\bfseries\color{kanjicolor} #2} \\[-0.2ex]
    {\fontsize{12}{14}\selectfont\color{romacolor}\sffamily #3}
  \end{tabular}%
  \hspace{0.2em}%
}

% \jpkana{kana}{romaji} -> Untuk hiragana/katakana yang tidak butuh furigana
\newcommand{\jpkana}[2]{%
  \begin{tabular}[b]{@{}c@{}}
    {\fontsize{10}{12}\selectfont\color{furicolor}\vphantom{あ}} \\[-0.2ex]
    {\fontsize{24}{28}\selectfont\bfseries\color{kanjicolor} #1} \\[-0.2ex]
    {\fontsize{12}{14}\selectfont\color{romacolor}\sffamily #2}
  \end{tabular}%
  \hspace{0.2em}%
}

% --- TCOLORBOX STYLES ---
\newtcolorbox{convobox}[1]{
    colback=boxbg,
    colframe=boxframe,
    boxrule=2pt,
    arc=5mm,
    title={\Large\bfseries \sffamily Kalimat #1},
    fonttitle=\color{white},
    attach boxed title to top left={xshift=5mm, yshift=-3mm},
    boxed title style={colback=boxframe, rounded corners},
    left=5mm, right=5mm, top=7mm, bottom=5mm,
    breakable
}

\newtcolorbox{lessonbox}[1]{
    colback=lessonbg,
    colframe=lessonframe,
    boxrule=2pt,
    arc=5mm,
    title={\Large\bfseries \sffamily 💡 #1},
    fonttitle=\color{white},
    attach boxed title to top center={yshift=-3mm},
    boxed title style={colback=lessonframe, rounded corners},
    left=5mm, right=5mm, top=7mm, bottom=5mm,
    breakable
}

\begin{document}
\tableofcontents
\chapter{Chatting with a Stranger on Your Trip in Japan}

% ================= PERCAKAPAN 1-4 =================
\begin{convobox}{1}
    \noindent
    \jpkana{こんにちは。}{Konnichiwa.} 
    \jpword{きょう}{今日}{Kyou} \jpkana{は}{wa} 
    \jpword{りょこう}{ご旅行}{go-ryokou} \jpkana{ですか？}{desu ka?}

    \vspace{1.5em}
    \noindent{\Large \color{articolor} \textbf{Arti:} Halo. Apakah hari ini kamu sedang jalan-jalan?}
\end{convobox}

\begin{convobox}{2}
    \noindent
    \jpkana{はい、}{Hai,} 
    \jpword{りょこう}{旅行}{ryokou} \jpkana{です。}{desu.}

    \vspace{1.5em}
    \noindent{\Large \color{articolor} \textbf{Arti:} Ya, (sedang) jalan-jalan.}
\end{convobox}

\begin{convobox}{3}
    \noindent
    \jpkana{どちら}{Dochira} \jpkana{の}{no} 
    \jpword{くに}{国}{kuni} \jpkana{から}{kara} 
    \jpkana{いらっしゃった}{irasshatta} \jpkana{ん}{n} \jpkana{ですか？}{desu ka?}

    \vspace{1.5em}
    \noindent{\Large \color{articolor} \textbf{Arti:} Dari negara mana kamu berasal?}
\end{convobox}

\begin{convobox}{4}
    \noindent
    \jpkana{インドネシア}{Indoneshia} \jpkana{から}{kara} 
    \jpword{き}{来}{ki}\jpkana{ました。}{mashita.}

    \vspace{1.5em}
    \noindent{\Large \color{articolor} \textbf{Arti:} Saya datang dari Indonesia.}
\end{convobox}

% --- LESSON BUFFER 1 ---
\begin{lessonbox}{Catatan Belajar: Asal Negara!}
    Halo teman-teman! Hebat banget kan Andi berani ngobrol sama orang Jepang asli? Di kalimat 3 dan 4, kita belajar cara kasih tahu dari mana kita berasal.\\
    
    Kalau ada yang tanya pakai kata \textbf{"\dots kara kimashita"} (datang dari\dots), kamu tinggal sebut negaramu lalu ditambah kata itu. Contohnya: \textit{Indoneshia kara kimashita!} Gampang banget, kan? Yuk lanjut lihat mereka mau ke mana!
\end{lessonbox}

% ================= PERCAKAPAN 5-9 =================
\begin{convobox}{5}
    \noindent
    \jpword{にほん}{日本}{Nihon} \jpkana{は}{wa} 
    \jpword{はじ}{初}{haji}\jpkana{めて}{mete} \jpkana{ですか？}{desu ka?}

    \vspace{1.5em}
    \noindent{\Large \color{articolor} \textbf{Arti:} Apakah ini pertama kalinya ke Jepang?}
\end{convobox}

\begin{convobox}{6}
    \noindent
    \jpkana{はい、}{Hai,} 
    \jpword{はじ}{初}{haji}\jpkana{めて}{mete} \jpkana{です。}{desu.}

    \vspace{1.5em}
    \noindent{\Large \color{articolor} \textbf{Arti:} Ya, pertama kali.}
\end{convobox}

\begin{convobox}{7}
    \noindent
    \jpword{きょう}{今日}{Kyou} \jpkana{は}{wa} 
    \jpkana{どちら}{dochira} \jpkana{まで}{made} 
    \jpword{い}{行}{i}\jpkana{かれる}{kareru} \jpkana{ん}{n} \jpkana{ですか？}{desu ka?}

    \vspace{1.5em}
    \noindent{\Large \color{articolor} \textbf{Arti:} Hari ini mau pergi sampai ke mana?}
\end{convobox}

\begin{convobox}{8}
    \noindent
    \jpword{きょうと}{京都}{Kyouto} \jpkana{まで}{made} 
    \jpword{い}{行}{i}\jpkana{きます。}{kimasu.}

    \vspace{1.5em}
    \noindent{\Large \color{articolor} \textbf{Arti:} Pergi sampai ke Kyoto.}
\end{convobox}

\begin{convobox}{9}
    \noindent
    \jpkana{そうですか。}{Sou desu ka.} 
    \jpword{わたし}{私}{Watashi} \jpkana{は}{wa} 
    \jpword{なごや}{名古屋}{Nagoya} \jpkana{まで}{made} 
    \jpword{い}{行}{i}\jpkana{きます。}{kimasu.}

    \vspace{1.5em}
    \noindent{\Large \color{articolor} \textbf{Arti:} Oh begitu. Saya pergi sampai ke Nagoya.}
\end{convobox}

% --- LESSON BUFFER 2 ---
\begin{lessonbox}{Catatan Belajar: Mau ke mana kita?}
    Perhatikan kata \textbf{まで (made)} di atas! Kata ini artinya "sampai". Kalau kamu mau bilang pergi sampai ke suatu tempat, sebut nama kotanya ditambah \textit{made ikimasu}. \\
    
    Orang Jepang suka nanya \textit{"dochira made?"} (mau sampai ke arah mana?). Nah, sekarang siapkan tiketmu, karena Andi mau ke Kyoto dan teman barunya mau ke Nagoya!
\end{lessonbox}

% ================= PERCAKAPAN 10-17 =================
\begin{convobox}{10}
    \noindent
    \jpword{とうちゃく}{到着}{Touchaku} \jpkana{したら、}{shitara,} 
    \jpkana{どこ}{doko} \jpkana{に}{ni} 
    \jpword{い}{行}{i}\jpkana{かれる}{kareru} 
    \jpword{よてい}{予定}{yotei} \jpkana{ですか？}{desu ka?}

    \vspace{1.5em}
    \noindent{\Large \color{articolor} \textbf{Arti:} Kalau sudah tiba, rencana mau pergi ke mana?}
\end{convobox}

\begin{convobox}{11}
    \noindent
    \jpword{さくら}{桜}{Sakura} \jpkana{が}{ga} 
    \jpword{ゆうめい}{有名}{yuumei} \jpkana{な}{na} 
    \jpword{てら}{お寺}{o-tera} \jpkana{に}{ni} 
    \jpword{い}{行}{i}\jpkana{く}{ku} \jpword{よてい}{予定}{yotei} \jpkana{です。}{desu.}

    \vspace{1.5em}
    \noindent{\Large \color{articolor} \textbf{Arti:} Rencananya pergi ke kuil yang terkenal dengan bunga sakuranya.}
\end{convobox}

\begin{convobox}{12}
    \noindent
    \jpword{たの}{楽}{Tano}\jpkana{しそう}{shisou} \jpkana{ですね。}{desu ne.}

    \vspace{1.5em}
    \noindent{\Large \color{articolor} \textbf{Arti:} Sepertinya menyenangkan ya.}
\end{convobox}

\begin{convobox}{13}
    \noindent
    \jpkana{それにしても、}{Sore ni shitemo,} 
    \jpword{にほんご}{日本語}{Nihongo} \jpkana{が}{ga} 
    \jpkana{とても}{totemo} \jpword{じょうず}{上手}{jouzu} \jpkana{ですね。}{desu ne.}

    \vspace{1.5em}
    \noindent{\Large \color{articolor} \textbf{Arti:} Ngomong-ngomong, bahasa Jepangnya sangat pintar ya.}
\end{convobox}

\begin{convobox}{14}
    \noindent
    \jpkana{いつ}{Itsu} \jpkana{から}{kara} 
    \jpword{にほんご}{日本語}{Nihongo} \jpkana{を}{o} 
    \jpword{べんきょう}{勉強}{benkyou} \jpkana{している}{shite iru} \jpkana{ん}{n} \jpkana{ですか？}{desu ka?}

    \vspace{1.5em}
    \noindent{\Large \color{articolor} \textbf{Arti:} Sejak kapan belajar bahasa Jepang?}
\end{convobox}

\begin{convobox}{15}
    \noindent
    \jpword{いちねんまえ}{1年前}{Ichi-nen mae} \jpkana{から}{kara} 
    \jpword{べんきょう}{勉強}{benkyou} \jpkana{しています。}{shite imasu.}

    \vspace{1.5em}
    \noindent{\Large \color{articolor} \textbf{Arti:} Belajar sejak 1 tahun yang lalu.}
\end{convobox}

\begin{convobox}{16}
    \noindent
    \jpkana{なんで}{Nande} \jpword{にほんご}{日本語}{Nihongo} \jpkana{の}{no} 
    \jpword{べんきょう}{勉強}{benkyou} \jpkana{を}{o} 
    \jpword{はじ}{始}{haji}\jpkana{めた}{meta} \jpkana{ん}{n} \jpkana{ですか？}{desu ka?}

    \vspace{1.5em}
    \noindent{\Large \color{articolor} \textbf{Arti:} Kenapa mulai belajar bahasa Jepang?}
\end{convobox}

\begin{convobox}{17}
    \noindent
    \jpword{にほん}{日本}{Nihon} \jpkana{の}{no} 
    \jpword{ぶんか}{文化}{bunka} \jpkana{と}{to} 
    \jpword{れきし}{歴史}{rekishi} \jpkana{に}{ni} 
    \jpword{きょうみ}{興味}{kyoumi} \jpkana{が}{ga} \jpkana{ある}{aru} \jpkana{から}{kara} \jpkana{です。}{desu.}

    \vspace{1.5em}
    \noindent{\Large \color{articolor} \textbf{Arti:} Karena saya tertarik dengan budaya dan sejarah Jepang.}
\end{convobox}

% --- LESSON BUFFER 3 ---
\begin{lessonbox}{Catatan Belajar: Pujian dan Hobi}
    Wah, bahasa Jepang Andi dipuji \textbf{上手 (jouzu)} yang artinya pintar/jago! Kalau kamu dipuji seperti ini, kamu bisa senyum bahagia lho. \\
    
    Di sini Andi juga ngomongin alasannya belajar bahasa Jepang karena tertarik dengan \textit{bunka} (budaya) dan \textit{rekishi} (sejarah). Orang Jepang sangat senang loh kalau kita tertarik dengan negara mereka!
\end{lessonbox}

% ================= PERCAKAPAN 18-27 =================
\begin{convobox}{18}
    \noindent
    \jpword{にほん}{日本}{Nihon} \jpkana{の}{no} 
    \jpword{た}{食}{ta}\jpkana{べ}{be}\jpword{もの}{物}{mono} \jpkana{は}{wa} 
    \jpword{す}{好}{su}\jpkana{き}{ki} \jpkana{ですか？}{desu ka?}

    \vspace{1.5em}
    \noindent{\Large \color{articolor} \textbf{Arti:} Apakah kamu suka makanan Jepang?}
\end{convobox}

\begin{convobox}{19}
    \noindent
    \jpkana{はい、}{Hai,} \jpword{だいす}{大好}{daisu}\jpkana{き}{ki} \jpkana{です。}{desu.}

    \vspace{1.5em}
    \noindent{\Large \color{articolor} \textbf{Arti:} Ya, sangat suka.}
\end{convobox}

\begin{convobox}{20}
    \noindent
    \jpword{いま}{今}{Ima} \jpkana{まで}{made} 
    \jpword{にほん}{日本}{Nihon} \jpkana{で}{de} 
    \jpword{た}{食}{ta}\jpkana{べた}{beta} \jpword{もの}{物}{mono} \jpkana{の}{no} 
    \jpword{なか}{中}{naka} \jpkana{で、}{de,}

    \vspace{1.5em}
    \noindent{\Large \color{articolor} \textbf{Arti:} Di antara makanan yang pernah dimakan di Jepang selama ini,}
\end{convobox}

\begin{convobox}{21}
    \noindent
    \jpword{なに}{何}{Nani} \jpkana{が}{ga} 
    \jpword{いちばん}{一番}{ichiban} \jpkana{おいしかった}{oishikatta} \jpkana{ですか？}{desu ka?}

    \vspace{1.5em}
    \noindent{\Large \color{articolor} \textbf{Arti:} Apa yang paling lezat?}
\end{convobox}

\begin{convobox}{22}
    \noindent
    \jpword{とうきょう}{東京}{Toukyou} \jpkana{で}{de} 
    \jpword{た}{食}{ta}\jpkana{べた}{beta} \jpword{すし}{寿司}{sushi} \jpkana{が}{ga} 
    \jpword{いちばん}{一番}{ichiban} \jpkana{おいしかった}{oishikatta} \jpkana{です。}{desu.}

    \vspace{1.5em}
    \noindent{\Large \color{articolor} \textbf{Arti:} Sushi yang saya makan di Tokyo adalah yang paling lezat.}
\end{convobox}

\begin{convobox}{23}
    \noindent
    \jpkana{（まもなく、}{Mamonaku,} 
    \jpword{なごや}{名古屋}{Nagoya}\jpkana{、}{,} \jpword{なごや}{名古屋}{Nagoya} \jpkana{です。）}{desu.}

    \vspace{1.5em}
    \noindent{\Large \color{articolor} \textbf{Arti:} (Sebentar lagi, tiba di Nagoya, Nagoya.)}
\end{convobox}

\begin{convobox}{24}
    \noindent
    \jpkana{あっ、}{A,,} \jpkana{もう}{mou} 
    \jpword{つ}{着}{tsu}\jpkana{きました。}{kimashita.}

    \vspace{1.5em}
    \noindent{\Large \color{articolor} \textbf{Arti:} Ah, sudah sampai.}
\end{convobox}

\begin{convobox}{25}
    \noindent
    \jpword{しんかんせん}{新幹線}{Shinkansen} \jpkana{は}{wa} 
    \jpword{はや}{速}{haya}\jpkana{い}{i} \jpkana{ですね。}{desu ne.}

    \vspace{1.5em}
    \noindent{\Large \color{articolor} \textbf{Arti:} Shinkansen itu cepat ya.}
\end{convobox}

\begin{convobox}{26}
    \noindent
    \jpword{にほん}{日本}{Nihon} \jpword{りょこう}{旅行}{ryokou}\jpkana{、}{,} 
    \jpword{たの}{楽}{tano}\jpkana{しんで}{shinde} \jpkana{ください。}{kudasai.}

    \vspace{1.5em}
    \noindent{\Large \color{articolor} \textbf{Arti:} Tolong nikmati perjalananmu di Jepang ya.}
\end{convobox}

\begin{convobox}{27}
    \noindent
    \jpkana{それでは、}{Sore dewa,} 
    \jpword{しつれい}{失礼}{shitsurei} \jpkana{します。}{shimasu.}

    \vspace{1.5em}
    \noindent{\Large \color{articolor} \textbf{Arti:} Kalau begitu, saya permisi dulu.}
\end{convobox}

% --- SUMMARY BOX ---
\vspace{2em}
\begin{tcolorbox}[
    colback=blue!5!white,
    colframe=blue!80!black,
    boxrule=3pt,
    arc=5mm,
    title={\huge\bfseries \sffamily 🌟 Rangkuman Bab 1 🌟},
    fonttitle=\color{white},
    attach boxed title to top center={yshift=-4mm},
    boxed title style={colback=blue!80!black, rounded corners},
    left=5mm, right=5mm, top=7mm, bottom=5mm
]
\Large \textbf{Kerja bagus teman-teman!} Di bab ini kita sudah belajar banyak hal keren saat ngobrol di dalam kereta Jepang (Shinkansen):
\begin{itemize}
    \item \textbf{Menjawab asal negara:} "\textit{\dots kara kimashita}" (Datang dari\dots)
    \item \textbf{Bilang ke mana kita pergi:} "\textit{\dots made ikimasu}" (Pergi sampai ke\dots)
    \item \textbf{Bicara soal durasi:} "\textit{Ichi-nen mae kara\dots}" (Sejak 1 tahun lalu\dots)
    \item \textbf{Membicarakan kesukaan (makanan):} "\textit{Daisuki desu!}" (Sangat suka!)
    \item \textbf{Berpamitan sopan:} "\textit{Shitsurei shimasu}" (Permisi / Saya pamit)
\end{itemize}
Jangan lupa diulang-ulang ya bacanya biar makin lancar seperti Andi! Sampai jumpa di perjalanan berikutnya!
\end{tcolorbox}

\newpage

\chapter{Help a Lost Tourist in Your Town}

% ================= PERCAKAPAN 1-7 =================
\begin{convobox}{1}
    \noindent
    \jpkana{あ、}{A,} \jpkana{すみません。}{sumimasen.} 

    \vspace{1.5em}
    \noindent{\Large \color{articolor} \textbf{Arti:} Ah, permisi.}
\end{convobox}

\begin{convobox}{2}
    \noindent
    \jpword{いま}{今}{Ima}\jpkana{、}{,} 
    \jpkana{ちょっと}{chotto} 
    \jpkana{いいですか？}{ii desu ka?}

    \vspace{1.5em}
    \noindent{\Large \color{articolor} \textbf{Arti:} Apakah kamu ada waktu sebentar sekarang?}
\end{convobox}

\begin{convobox}{3}
    \noindent
    \jpword{びじゅつかん}{美術館}{Bijutsukan} \jpkana{に}{ni} 
    \jpword{い}{行}{i}\jpkana{きたいんですが、}{kitai n desu ga,}

    \vspace{1.5em}
    \noindent{\Large \color{articolor} \textbf{Arti:} Saya ingin pergi ke museum seni,}
\end{convobox}

\begin{convobox}{4}
    \noindent
    \jpkana{この}{Kono} 
    \jpword{ちか}{近}{chika}\jpkana{く}{ku} \jpkana{に}{ni} 
    \jpkana{ありますか？}{arimasu ka?}

    \vspace{1.5em}
    \noindent{\Large \color{articolor} \textbf{Arti:} Apakah ada di dekat sini?}
\end{convobox}

\begin{convobox}{5}
    \noindent
    \jpkana{はい。}{Hai.} \jpkana{ありますよ。}{Arimasu yo.}

    \vspace{1.5em}
    \noindent{\Large \color{articolor} \textbf{Arti:} Ya. Ada kok.}
\end{convobox}

\begin{convobox}{6}
    \noindent
    \jpword{ある}{歩}{Aru}\jpkana{いて}{ite} 
    \jpword{なんぷん}{何分}{nanpun} 
    \jpkana{ぐらいですか？}{gurai desu ka?}

    \vspace{1.5em}
    \noindent{\Large \color{articolor} \textbf{Arti:} Kira-kira berapa menit kalau jalan kaki?}
\end{convobox}

\begin{convobox}{7}
    \noindent
    \jpword{ある}{歩}{Aru}\jpkana{いて}{ite} 
    \jpword{にじゅっぷん}{20分}{ni-juppun} 
    \jpkana{ぐらいです。}{gurai desu.}

    \vspace{1.5em}
    \noindent{\Large \color{articolor} \textbf{Arti:} Kira-kira 20 menit jalan kaki.}
\end{convobox}

% % --- LESSON BUFFER 1 ---
\begin{lessonbox}{Catatan Belajar: Nanya Tempat Waktu}
    Wah, ada turis yang nyasar nih! Kalau kamu mau bertanya apakah ada suatu tempat di dekat sini, kamu bisa bilang: \textbf{"Kono chikaku ni arimasu ka?"} (Apakah ada di dekat sini?). \\
    
    Terus, kalau kamu mau tahu berapa lama perjalanannya kalau jalan kaki, gunakan \textbf{"Aruite nanpun gurai desu ka?"}. Yuk kita lihat apakah turisnya mau jalan kaki selama 20 menit!
\end{lessonbox}

% ================= PERCAKAPAN 8-16 =================
\begin{convobox}{8}
    \noindent
    \jpword{けっこう}{結構}{Kekkou} 
    \jpword{とお}{遠}{too}\jpkana{いですね。}{i desu ne.}

    \vspace{1.5em}
    \noindent{\Large \color{articolor} \textbf{Arti:} Lumayan jauh ya.}
\end{convobox}

\begin{convobox}{9}
    \noindent
    \jpkana{タクシー}{Takushii} \jpkana{は}{wa} 
    \jpword{たか}{高}{taka}\jpkana{いですか？}{i desu ka?}

    \vspace{1.5em}
    \noindent{\Large \color{articolor} \textbf{Arti:} Apakah taksi mahal?}
\end{convobox}

\begin{convobox}{10}
    \noindent
    \jpkana{タクシー}{Takushii} \jpkana{は}{wa} 
    \jpkana{ちょっと}{chotto} 
    \jpword{たか}{高}{taka}\jpkana{いと}{i to} 
    \jpword{おも}{思}{omo}\jpkana{います。}{imasu.}

    \vspace{1.5em}
    \noindent{\Large \color{articolor} \textbf{Arti:} Menurutku taksi sedikit mahal.}
\end{convobox}

\begin{convobox}{11}
    \noindent
    \jpword{でんしゃ}{電車}{Densha} \jpkana{とか}{toka} 
    \jpkana{バス}{basu} \jpkana{は}{wa} 
    \jpkana{ありますか？}{arimasu ka?}

    \vspace{1.5em}
    \noindent{\Large \color{articolor} \textbf{Arti:} Apakah ada kereta atau bus?}
\end{convobox}

\begin{convobox}{12}
    \noindent
    \jpkana{バス}{Basu} \jpkana{は}{wa} 
    \jpkana{ありませんが、}{arimasen ga,} 
    \jpword{でんしゃ}{電車}{densha} \jpkana{は}{wa} 
    \jpkana{ありますよ。}{arimasu yo.}

    \vspace{1.5em}
    \noindent{\Large \color{articolor} \textbf{Arti:} Bus tidak ada, tapi kereta ada lho.}
\end{convobox}

\begin{convobox}{13}
    \noindent
    \jpword{えき}{駅}{Eki} \jpkana{は}{wa} 
    \jpkana{どこ}{doko} \jpkana{に}{ni} 
    \jpkana{ありますか？}{arimasu ka?}

    \vspace{1.5em}
    \noindent{\Large \color{articolor} \textbf{Arti:} Stasiunnya ada di mana?}
\end{convobox}

\begin{convobox}{14}
    \noindent
    \jpkana{この}{Kono} 
    \jpword{みち}{道}{michi} \jpkana{を}{o} 
    \jpkana{まっすぐ}{massugu} 
    \jpword{い}{行}{i}\jpkana{って、}{tte,}

    \vspace{1.5em}
    \noindent{\Large \color{articolor} \textbf{Arti:} Ikuti jalan ini lurus terus,}
\end{convobox}

\begin{convobox}{15}
    \noindent
    \jpword{かど}{角}{Kado} \jpkana{を}{o} 
    \jpword{みぎ}{右}{migi} \jpkana{に}{ni} 
    \jpword{ま}{曲}{ma}\jpkana{がると、}{garu to,} 
    \jpword{えき}{駅}{eki} \jpkana{が}{ga} 
    \jpkana{あります。}{arimasu.}

    \vspace{1.5em}
    \noindent{\Large \color{articolor} \textbf{Arti:} Kalau belok kanan di tikungan, stasiunnya ada di sana.}
\end{convobox}

\begin{convobox}{16}
    \noindent
    \jpkana{ありがとうございます。}{Arigatou gozaimasu.}

    \vspace{1.5em}
    \noindent{\Large \color{articolor} \textbf{Arti:} Terima kasih banyak.}
\end{convobox}

% --- LESSON BUFFER 2 ---
\begin{lessonbox}{Catatan Belajar: Memberi Petunjuk Jalan!}
    Ternyata turisnya mau naik kereta! Saat kamu mau memberi tahu arah jalan ke seseorang, kamu bisa memakai kosakata ini:
    \begin{itemize}
        \item \textbf{Massugu} = Lurus terus
        \item \textbf{Migi} = Kanan (kalau Kiri = \textit{Hidari})
        \item \textbf{Magaru} = Belok
    \end{itemize}
    Jadi, \textit{"Massugu itte, migi ni magaru"} artinya "Lurus terus, lalu belok kanan". Keren banget kan bisa bantu turis nyasar!
\end{lessonbox}

% ================= PERCAKAPAN 17-22 =================
\begin{convobox}{17}
    \noindent
    \jpkana{この}{Kono} 
    \jpword{まち}{街}{machi} \jpkana{で}{de} 
    \jpword{ゆうめい}{有名}{yuumei}\jpkana{な}{na} 
    \jpword{かんこうち}{観光地}{kankouchi} \jpkana{や、}{ya,}

    \vspace{1.5em}
    \noindent{\Large \color{articolor} \textbf{Arti:} Tempat wisata yang terkenal di kota ini, atau}
\end{convobox}

\begin{convobox}{18}
    \noindent
    \jpkana{おすすめ}{Osusume} \jpkana{の}{no} 
    \jpword{ばしょ}{場所}{basho} \jpkana{は}{wa} 
    \jpkana{ありますか？}{arimasu ka?}

    \vspace{1.5em}
    \noindent{\Large \color{articolor} \textbf{Arti:} apakah ada tempat yang kamu rekomendasikan?}
\end{convobox}

\begin{convobox}{19}
    \noindent
    \jpword{おお}{大}{Oo}\jpkana{きな}{kina} 
    \jpword{ひろば}{広場}{hiroba} \jpkana{が}{ga} 
    \jpword{ゆうめい}{有名}{yuumei} \jpkana{です。}{desu.}

    \vspace{1.5em}
    \noindent{\Large \color{articolor} \textbf{Arti:} Alun-alun yang besar sangat terkenal lho.}
\end{convobox}

\begin{convobox}{20}
    \noindent
    \jpkana{いろいろな}{Iroiro na} 
    \jpword{みせ}{お店}{o-mise} \jpkana{が}{ga} \jpkana{あって}{atte} 
    \jpword{おも}{面}{omo}\jpword{しろ}{白}{shiro}\jpkana{いですよ。}{i desu yo.}

    \vspace{1.5em}
    \noindent{\Large \color{articolor} \textbf{Arti:} Ada bermacam-macam toko dan tempatnya sangat menarik.}
\end{convobox}

\begin{convobox}{21}
    \noindent
    \jpkana{ときどき}{Tokidoki} 
    \jpkana{イベント}{ibento} \jpkana{も}{mo} 
    \jpkana{あります。}{arimasu.}

    \vspace{1.5em}
    \noindent{\Large \color{articolor} \textbf{Arti:} Kadang-kadang juga ada acara (event).}
\end{convobox}

\begin{convobox}{22}
    \noindent
    \jpkana{わかりました。}{Wakarimashita.} 
    \jpkana{あとで}{Ato de} 
    \jpword{い}{行}{i}\jpkana{ってみます。}{tte mimasu.}

    \vspace{1.5em}
    \noindent{\Large \color{articolor} \textbf{Arti:} Saya mengerti. Nanti saya akan coba pergi ke sana.}
\end{convobox}

% --- LESSON BUFFER 3 ---
\begin{lessonbox}{Catatan Belajar: Rekomendasi Tempat!}
    Kalau kamu jalan-jalan dan bingung mau ke mana, kata ajaibnya adalah \textbf{おすすめ (Osusume)} yang berarti "Rekomendasi". \\
    
    Orang Jepang biasanya akan menjawab dengan antusias dan memberitahu tempat yang \textbf{有名 (Yuumei)} alias terkenal, atau yang \textbf{面白い (Omoshiroi)} alias menarik/seru!
\end{lessonbox}

% ================= PERCAKAPAN 23-34 =================
\begin{convobox}{23}
    \noindent
    \jpkana{そういえば、}{Sou ieba,} 
    \jpkana{のどが}{nodo ga} 
    \jpkana{かわいた}{kawaita} 
    \jpkana{ので}{node}

    \vspace{1.5em}
    \noindent{\Large \color{articolor} \textbf{Arti:} Ngomong-ngomong, karena saya merasa haus,}
\end{convobox}

\begin{convobox}{24}
    \noindent
    \jpword{みず}{水}{Mizu} \jpkana{を}{o} 
    \jpword{か}{買}{ka}\jpkana{いたいんですが、}{itai n desu ga,}

    \vspace{1.5em}
    \noindent{\Large \color{articolor} \textbf{Arti:} saya ingin membeli air,}
\end{convobox}

\begin{convobox}{25}
    \noindent
    \jpkana{この}{Kono} 
    \jpword{ちか}{近}{chika}\jpkana{く}{ku} \jpkana{に}{ni} 
    \jpkana{スーパー}{suupaa} \jpkana{とか}{toka} 
    \jpkana{コンビニ}{konbini} \jpkana{は}{wa} 
    \jpkana{ありますか？}{arimasu ka?}

    \vspace{1.5em}
    \noindent{\Large \color{articolor} \textbf{Arti:} apakah di dekat sini ada supermarket atau minimarket?}
\end{convobox}

\begin{convobox}{26}
    \noindent
    \jpkana{はい。}{Hai.} 
    \jpkana{あの}{ano} 
    \jpword{きいろ}{黄色}{kiiro}\jpkana{い}{i} 
    \jpword{かんばん}{看板}{kanban}\jpkana{、}{,} 
    \jpword{み}{見}{mi}\jpkana{えますか？}{emasu ka?}

    \vspace{1.5em}
    \noindent{\Large \color{articolor} \textbf{Arti:} Ya. Apakah kamu bisa melihat papan reklame kuning itu?}
\end{convobox}

\begin{convobox}{27}
    \noindent
    \jpkana{あれが}{Are ga} 
    \jpkana{コンビニ}{konbini} \jpkana{です。}{desu.}

    \vspace{1.5em}
    \noindent{\Large \color{articolor} \textbf{Arti:} Itu adalah minimarket.}
\end{convobox}

\begin{convobox}{28}
    \noindent
    \jpkana{わかりました。}{Wakarimashita.} 
    \jpkana{ありがとうございます。}{Arigatou gozaimasu.}

    \vspace{1.5em}
    \noindent{\Large \color{articolor} \textbf{Arti:} Saya mengerti. Terima kasih banyak.}
\end{convobox}

\begin{convobox}{29}
    \noindent
    \jpkana{あ、}{A,} 
    \jpkana{もうすぐ}{mou sugu} 
    \jpword{ひる}{お昼}{o-hiru} \jpkana{ですね。}{desu ne.}

    \vspace{1.5em}
    \noindent{\Large \color{articolor} \textbf{Arti:} Ah, sebentar lagi makan siang ya.}
\end{convobox}

\begin{convobox}{30}
    \noindent
    \jpkana{この}{Kono} 
    \jpword{ちか}{近}{chika}\jpkana{く}{ku} \jpkana{に}{ni} 
    \jpkana{おすすめの}{osusume no} 
    \jpkana{レストラン}{resutoran} \jpkana{や}{ya}

    \vspace{1.5em}
    \noindent{\Large \color{articolor} \textbf{Arti:} Di dekat sini, apakah ada restoran yang direkomendasikan, atau}
\end{convobox}

\begin{convobox}{31}
    \noindent
    \jpkana{カフェ}{kafe} \jpkana{は}{wa} 
    \jpkana{ありますか？}{arimasu ka?}

    \vspace{1.5em}
    \noindent{\Large \color{articolor} \textbf{Arti:} kafe?}
\end{convobox}

\begin{convobox}{32}
    \noindent
    \jpword{えき}{駅}{Eki} \jpkana{の}{no} 
    \jpword{ちか}{近}{chika}\jpkana{く}{ku} \jpkana{に}{ni} 
    \jpkana{『Pomodoro』}{Pomodoro} \jpkana{という}{to iu}

    \vspace{1.5em}
    \noindent{\Large \color{articolor} \textbf{Arti:} Di dekat stasiun ada tempat bernama 'Pomodoro', yaitu}
\end{convobox}

\begin{convobox}{33}
    \noindent
    \jpkana{イタリアンレストラン}{Itarian resutoran} \jpkana{が}{ga} 
    \jpkana{あります。}{arimasu.}

    \vspace{1.5em}
    \noindent{\Large \color{articolor} \textbf{Arti:} restoran Italia.}
\end{convobox}

\begin{convobox}{34}
    \noindent
    \jpkana{ピザ}{Piza} \jpkana{が}{ga} 
    \jpkana{とても}{totemo} 
    \jpkana{おいしい}{oishii} \jpkana{ので}{node} 
    \jpkana{おすすめ}{osusume} \jpkana{です。}{desu.}

    \vspace{1.5em}
    \noindent{\Large \color{articolor} \textbf{Arti:} Karena pizzanya sangat enak, saya merekomendasikannya.}
\end{convobox}

% --- LESSON BUFFER 4 ---
\begin{lessonbox}{Catatan Belajar: Menyatakan Alasan!}
    Di bagian ini, turisnya memakai kata \textbf{ので (Node)} yang artinya "Karena".
    \begin{itemize}
        \item \textit{Nodo ga kawaita \textbf{node}} = Karena saya haus
        \item \textit{Piza ga oishii \textbf{node}} = Karena pizzanya enak
    \end{itemize}
    Jadi, kalau kamu mau memberi tahu alasan kamu melakukan sesuatu, kamu bisa menempelkan \textbf{"node"} di tengah kalimat. Coba deh bikin kalimat alasanmu sendiri!
\end{lessonbox}

% ================= PERCAKAPAN 35-42 =================
\begin{convobox}{35}
    \noindent
    \jpkana{あっ、}{A,,} 
    \jpkana{どうしよう…}{dou shiyou...}

    \vspace{1.5em}
    \noindent{\Large \color{articolor} \textbf{Arti:} Ah, bagaimana ini...}
\end{convobox}

\begin{convobox}{36}
    \noindent
    \jpkana{スマホ}{Sumaho} \jpkana{の}{no} 
    \jpkana{バッテリー}{batterii} \jpkana{が}{ga} 
    \jpkana{もうすぐ}{mou sugu} 
    \jpkana{なくなります。}{nakunarimasu.}

    \vspace{1.5em}
    \noindent{\Large \color{articolor} \textbf{Arti:} Baterai HP saya sebentar lagi mau habis.}
\end{convobox}

\begin{convobox}{37}
    \noindent
    \jpkana{この}{Kono} 
    \jpword{ちか}{近}{chika}\jpkana{く}{ku} \jpkana{に}{ni} 
    \jpkana{スマホ}{sumaho} \jpkana{を}{o} 
    \jpword{じゅうでん}{充電}{juuden} \jpkana{できる}{dekiru} 
    \jpword{ばしょ}{場所}{basho} \jpkana{は}{wa} 
    \jpkana{ありますか？}{arimasu ka?}

    \vspace{1.5em}
    \noindent{\Large \color{articolor} \textbf{Arti:} Apakah ada tempat untuk bisa mengecas HP di dekat sini?}
\end{convobox}

\begin{convobox}{38}
    \noindent
    \jpkana{コンビニ}{Konbini} \jpkana{で}{de} 
    \jpkana{モバイルバッテリー}{mobairu batterii} \jpkana{が}{ga}

    \vspace{1.5em}
    \noindent{\Large \color{articolor} \textbf{Arti:} Powerbank di minimarket,}
\end{convobox}

\begin{convobox}{39}
    \noindent
    \jpword{か}{買}{ka}\jpkana{えると}{eru to} 
    \jpword{おも}{思}{omo}\jpkana{います。}{imasu.}

    \vspace{1.5em}
    \noindent{\Large \color{articolor} \textbf{Arti:} menurutku bisa dibeli.}
\end{convobox}

\begin{convobox}{40}
    \noindent
    \jpkana{わかりました。}{Wakarimashita.}

    \vspace{1.5em}
    \noindent{\Large \color{articolor} \textbf{Arti:} Saya mengerti.}
\end{convobox}

\begin{convobox}{41}
    \noindent
    \jpkana{いろいろ}{Iroiro} 
    \jpword{おし}{教}{oshi}\jpkana{えてくださって}{ete kudasatte} 
    \jpkana{ありがとうございました。}{arigatou gozaimashita.}

    \vspace{1.5em}
    \noindent{\Large \color{articolor} \textbf{Arti:} Terima kasih banyak sudah mengajariku (memberitahuku) banyak hal.}
\end{convobox}

\begin{convobox}{42}
    \noindent
    \jpkana{それでは、}{Sore dewa,} 
    \jpword{しつれい}{失礼}{shitsurei} \jpkana{します。}{shimasu.}

    \vspace{1.5em}
    \noindent{\Large \color{articolor} \textbf{Arti:} Kalau begitu, saya permisi dulu.}
\end{convobox}

% --- SUMMARY BOX ---
\vspace{2em}
\begin{tcolorbox}[
    colback=blue!5!white,
    colframe=blue!80!black,
    boxrule=3pt,
    arc=5mm,
    title={\huge\bfseries \sffamily 🌟 Rangkuman Bab 2 🌟},
    fonttitle=\color{white},
    attach boxed title to top center={yshift=-4mm},
    boxed title style={colback=blue!80!black, rounded corners},
    left=5mm, right=5mm, top=7mm, bottom=5mm
]
\Large \textbf{Yey, kamu berhasil bantu turis!} Di bab ini kita belajar banyak kosakata super penting kalau jalan-jalan:
\begin{itemize}
    \item \textbf{Arah Jalan:} \textit{Massugu} (Lurus), \textit{Migi / Hidari} (Kanan / Kiri), \textit{Magaru} (Belok).
    \item \textbf{Bertanya Lokasi:} \textit{\dots arimasu ka?} (Apakah ada\dots?)
    \item \textbf{Rekomendasi:} \textit{Osusume} (Rekomendasi / Saran yang bagus).
    \item \textbf{Memberi Alasan:} Gunakan partikel \textit{Node} (Karena).
    \item \textbf{Darurat HP Mati:} \textit{Juuden dekiru} (Bisa mengecas baterai).
\end{itemize}
Bagus sekali! Sekarang kamu sudah siap jadi pemandu wisata dadakan di kotamu. Semangat terus belajarnya ya!
\end{tcolorbox}

\newpage
\chapter{Catching Up with Friends at a Cozy Autumn Café}

% ================= PERCAKAPAN 1-5 =================
\begin{convobox}{1}
    \noindent
    \jpword{ひさ}{久}{Hisa}\jpkana{しぶり！}{shiburi!} 
    \jpword{げんき}{元気}{Genki}\jpkana{だった？}{datta?}

    \vspace{1.5em}
    \noindent{\Large \color{articolor} \textbf{Arti:} Lama tak jumpa! Apa kabar?}
\end{convobox}

\begin{convobox}{2}
    \noindent
    \jpkana{うん、}{Un,} 
    \jpword{げんき}{元気}{genki}\jpkana{だったよ。}{datta yo.} 
    \jpkana{たなか}{Tanaka}\jpkana{さんは？}{san wa?}

    \vspace{1.5em}
    \noindent{\Large \color{articolor} \textbf{Arti:} Ya, aku baik-baik saja. Kalau Tanaka-san?}
\end{convobox}

\begin{convobox}{3}
    \noindent
    \jpword{さいご}{最後}{Saigo} \jpkana{に}{ni} 
    \jpword{あ}{会}{a}\jpkana{ったのは、}{tta no wa,} 
    \jpkana{いつだったっけ？}{itsu datta kke?}

    \vspace{1.5em}
    \noindent{\Large \color{articolor} \textbf{Arti:} Terakhir kali kita bertemu, kapan ya?}
\end{convobox}

\begin{convobox}{4}
    \noindent
    \jpkana{えーっと…}{Eetto...} 
    \jpword{きょねん}{去年}{kyonen} \jpkana{の}{no} 
    \jpword{ふゆ}{冬}{fuyu} \jpkana{じゃない？}{ja nai?}

    \vspace{1.5em}
    \noindent{\Large \color{articolor} \textbf{Arti:} Emm... bukannya musim dingin tahun lalu?}
\end{convobox}

\begin{convobox}{5}
    \noindent
    \jpkana{いっしょに}{Issho ni} 
    \jpword{おんせん}{温泉}{onsen} \jpkana{に}{ni} 
    \jpword{い}{行}{i}\jpkana{ったよね。}{tta yo ne.}

    \vspace{1.5em}
    \noindent{\Large \color{articolor} \textbf{Arti:} Kita pergi ke pemandian air panas (onsen) bareng, kan.}
\end{convobox}

% --- LESSON BUFFER 1 ---
\begin{lessonbox}{Catatan Belajar: Selamat Datang di Bahasa Kasual!}
    Teman-teman sadar nggak? Di bab ini nggak ada kata \textbf{"desu"} atau \textbf{"masu"}! Kalau kita ngobrol dengan teman dekat, kita pakai bentuk kasual (Biasa).\\
    
    Misalnya, nanya kabar nggak lagi pakai \textit{"Genki desu ka?"}, tapi cukup \textbf{"Genki datta?"}. Trus, "Ya" bukan lagi \textit{"Hai"}, melainkan \textbf{"Un"}. Lebih santai dan akrab banget, kan?
\end{lessonbox}

% ================= PERCAKAPAN 6-11 =================
\begin{convobox}{6}
    \noindent
    \jpword{さいきん}{最近}{Saikin} \jpkana{は}{wa} 
    \jpkana{どう？}{dou?} 
    \jpword{いそが}{忙}{isoga}\jpkana{しい？}{shii?}

    \vspace{1.5em}
    \noindent{\Large \color{articolor} \textbf{Arti:} Akhir-akhir ini gimana? Sibuk?}
\end{convobox}

\begin{convobox}{7}
    \noindent
    \jpkana{うん、}{Un,} 
    \jpword{しごと}{仕事}{shigoto} \jpkana{が}{ga} 
    \jpkana{けっこう}{kekkou} 
    \jpword{いそが}{忙}{isoga}\jpkana{しい。}{shii.}

    \vspace{1.5em}
    \noindent{\Large \color{articolor} \textbf{Arti:} Ya, kerjaan lumayan sibuk.}
\end{convobox}

\begin{convobox}{8}
    \noindent
    \jpkana{でも}{Demo} 
    \jpword{たの}{楽}{tano}\jpkana{しいよ。}{shii yo.}

    \vspace{1.5em}
    \noindent{\Large \color{articolor} \textbf{Arti:} Tapi menyenangkan kok.}
\end{convobox}

\begin{convobox}{9}
    \noindent
    \jpkana{そっか。}{Sokka.} 
    \jpword{やす}{休}{yasu}\jpkana{みの}{mi no} 
    \jpword{ひ}{日}{hi} \jpkana{は}{wa} 
    \jpkana{なにしてる？}{nani shiteru?}

    \vspace{1.5em}
    \noindent{\Large \color{articolor} \textbf{Arti:} Begitu ya. Pas hari libur ngapain aja?}
\end{convobox}

\begin{convobox}{10}
    \noindent
    \jpword{りょうり}{料理}{Ryouri} \jpkana{を}{o} 
    \jpkana{したり、}{shitari,} 
    \jpkana{カフェ}{kafe} \jpkana{に}{ni} 
    \jpword{い}{行}{i}\jpkana{ったり、}{ttari,}

    \vspace{1.5em}
    \noindent{\Large \color{articolor} \textbf{Arti:} Memasak, pergi ke kafe,}
\end{convobox}

\begin{convobox}{11}
    \noindent
    \jpword{いえ}{家}{Ie} \jpkana{で}{de} 
    \jpkana{ごろごろしたり}{gorogoro shitari} 
    \jpkana{してるよ。}{shiteru yo.}

    \vspace{1.5em}
    \noindent{\Large \color{articolor} \textbf{Arti:} bermalas-malasan di rumah, gitu deh.}
\end{convobox}

% --- LESSON BUFFER 2 ---
\begin{lessonbox}{Catatan Belajar: Ngapain Aja? Pakai \textit{~tari}!}
    Kalau kamu mau cerita beberapa aktivitas yang kamu lakuin (nggak harus berurutan), kamu bisa pakai pola \textbf{~tari, ~tari shimasu}. \\
    
    Contoh di atas: \textit{Ryouri o shi\textbf{tari}} (masak), \textit{Kafe ni i\textbf{ttari}} (pergi ke kafe). Ini asik banget dipakai kalau ditanya "Liburan ngapain aja?"
\end{lessonbox}

% ================= PERCAKAPAN 12-19 =================
\begin{convobox}{12}
    \noindent
    \jpword{さいきん}{最近}{Saikin}\jpkana{、}{,} 
    \jpkana{けっこう}{kekkou} 
    \jpword{さむ}{寒}{samu}\jpkana{くなってきたよね。}{ku natte kita yo ne.}

    \vspace{1.5em}
    \noindent{\Large \color{articolor} \textbf{Arti:} Akhir-akhir ini, (cuacanya) sudah lumayan dingin ya.}
\end{convobox}

\begin{convobox}{13}
    \noindent
    \jpword{さむ}{寒}{Samu}\jpkana{いのは}{i no wa} 
    \jpword{にがて}{苦手}{nigate}\jpkana{？}{?}

    \vspace{1.5em}
    \noindent{\Large \color{articolor} \textbf{Arti:} Kamu nggak tahan sama dingin?}
\end{convobox}

\begin{convobox}{14}
    \noindent
    \jpword{さむ}{寒}{Samu}\jpkana{いのは}{i no wa} 
    \jpword{だいじょうぶ}{大丈夫}{daijoubu}\jpkana{。}{.}

    \vspace{1.5em}
    \noindent{\Large \color{articolor} \textbf{Arti:} Kalau dingin nggak apa-apa (tahan).}
\end{convobox}

\begin{convobox}{15}
    \noindent
    \jpword{あつ}{暑}{Atsu}\jpkana{い}{i} 
    \jpword{ほう}{方}{hou} \jpkana{が}{ga} 
    \jpword{にがて}{苦手}{nigate} \jpkana{かな。}{kana.}

    \vspace{1.5em}
    \noindent{\Large \color{articolor} \textbf{Arti:} Kayaknya aku lebih nggak tahan panas deh.}
\end{convobox}

\begin{convobox}{16}
    \noindent
    \jpword{さいきん}{最近}{Saikin}\jpkana{、}{,} 
    \jpkana{どこか}{doko ka} 
    \jpword{りょこう}{旅行}{ryokou} \jpkana{に}{ni} 
    \jpword{い}{行}{i}\jpkana{った？}{tta?}

    \vspace{1.5em}
    \noindent{\Large \color{articolor} \textbf{Arti:} Akhir-akhir ini, ada pergi liburan ke suatu tempat?}
\end{convobox}

\begin{convobox}{17}
    \noindent
    \jpword{せんげつ}{先月}{Sengetsu}\jpkana{、}{,} 
    \jpword{なごや}{名古屋}{Nagoya} \jpkana{に}{ni} 
    \jpword{い}{行}{i}\jpkana{ったよ。}{tta yo.}

    \vspace{1.5em}
    \noindent{\Large \color{articolor} \textbf{Arti:} Bulan lalu, (aku) pergi ke Nagoya lho.}
\end{convobox}

\begin{convobox}{18}
    \noindent
    \jpkana{ジブリパーク}{Jiburi paaku} \jpkana{と}{to} 
    \jpword{なごやじょう}{名古屋城}{Nagoya-jou} \jpkana{に}{ni} 
    \jpword{い}{行}{i}\jpkana{った。}{tta.}

    \vspace{1.5em}
    \noindent{\Large \color{articolor} \textbf{Arti:} Pergi ke Ghibli Park dan Kastil Nagoya.}
\end{convobox}

\begin{convobox}{19}
    \noindent
    \jpkana{おいしい}{Oishii} 
    \jpword{もの}{物}{mono} \jpkana{も}{mo} 
    \jpkana{たくさん}{takusan} 
    \jpword{た}{食}{ta}\jpkana{べて、}{bete,} 
    \jpkana{すごく}{sugoku} 
    \jpword{たの}{楽}{tano}\jpkana{しかった。}{shikatta.}

    \vspace{1.5em}
    \noindent{\Large \color{articolor} \textbf{Arti:} Makan banyak makanan enak, (pokoknya) sangat menyenangkan.}
\end{convobox}

% --- LESSON BUFFER 3 ---
\begin{lessonbox}{Catatan Belajar: Tahan atau Nggak Tahan?}
    Orang Jepang suka nanya apakah kita bisa tahan sesuatu (makanan pedas, cuaca panas/dingin).
    \begin{itemize}
        \item \textbf{苦手 (Nigate)} = Lemah terhadap sesuatu / Nggak tahan / Kurang suka.
        \item \textbf{大丈夫 (Daijoubu)} = Nggak apa-apa / Aman / Tahan.
    \end{itemize}
    Jadi kalau ditanya "Samui no wa nigate?" artinya "Kamu nggak kuat dingin ya?".
\end{lessonbox}

% ================= PERCAKAPAN 20-25 =================
\begin{convobox}{20}
    \noindent
    \jpkana{あっ、}{A,,} 
    \jpword{きょう}{今日}{kyou} \jpkana{は}{wa} 
    \jpkana{そろそろ}{sorosoro} 
    \jpword{かえ}{帰}{kae}\jpkana{らなくちゃ。}{ranakucha.}

    \vspace{1.5em}
    \noindent{\Large \color{articolor} \textbf{Arti:} Ah, hari ini harus segera pulang nih.}
\end{convobox}

\begin{convobox}{21}
    \noindent
    \jpkana{よかったら、}{Yokattara,} 
    \jpword{らいげつ}{来月}{raigetsu} \jpkana{も}{mo} 
    \jpkana{また}{mata} 
    \jpkana{ランチ}{ranchi} \jpkana{か}{ka} 
    \jpword{おちゃ}{お茶}{o-cha} \jpkana{でも}{demo} 
    \jpkana{しない？}{shinai?}

    \vspace{1.5em}
    \noindent{\Large \color{articolor} \textbf{Arti:} Kalau berkenan, bulan depan mau makan siang atau minum teh lagi nggak?}
\end{convobox}

\begin{convobox}{22}
    \noindent
    \jpkana{いいね！}{Ii ne!}

    \vspace{1.5em}
    \noindent{\Large \color{articolor} \textbf{Arti:} Ide bagus!}
\end{convobox}

\begin{convobox}{23}
    \noindent
    \jpword{へいじつ}{平日}{Heijitsu} \jpkana{は}{wa} 
    \jpword{しごと}{仕事}{shigoto} \jpkana{だけど、}{dakedo,} 
    \jpword{どにち}{土日}{donichi} \jpkana{は}{wa} 
    \jpword{だいたい}{大体}{daitai} 
    \jpword{あ}{空}{a}\jpkana{いてるよ。}{iteru yo.}

    \vspace{1.5em}
    \noindent{\Large \color{articolor} \textbf{Arti:} Kalau hari biasa sih kerja, tapi akhir pekan biasanya kosong (jadwalnya) kok.}
\end{convobox}

\begin{convobox}{24}
    \noindent
    \jpkana{じゃあ、}{Jaa,} \jpkana{また}{mata} 
    \jpword{れんらく}{連絡}{renraku} \jpkana{するね。}{suru ne.}

    \vspace{1.5em}
    \noindent{\Large \color{articolor} \textbf{Arti:} Kalau gitu, nanti aku hubungi lagi ya.}
\end{convobox}

\begin{convobox}{25}
    \noindent
    \jpword{きょう}{今日}{Kyou} \jpkana{は}{wa} \jpkana{ありがとう。}{arigatou.}

    \vspace{1.5em}
    \noindent{\Large \color{articolor} \textbf{Arti:} Terima kasih ya untuk hari ini.}
\end{convobox}

% --- SUMMARY BOX ---
\vspace{2em}
\begin{tcolorbox}[
    colback=blue!5!white,
    colframe=blue!80!black,
    boxrule=3pt,
    arc=5mm,
    title={\huge\bfseries \sffamily 🌟 Rangkuman Bab 3 🌟},
    fonttitle=\color{white},
    attach boxed title to top center={yshift=-4mm},
    boxed title style={colback=blue!80!black, rounded corners},
    left=5mm, right=5mm, top=7mm, bottom=5mm
]
\Large \textbf{Keren banget! Sekarang kamu bisa ngobrol santai!} Di bab ini kita belajar:
\begin{itemize}
    \item \textbf{Sapaan Akrab:} \textit{Hisashiburi!} (Lama tak jumpa!)
    \item \textbf{Menyebut Berbagai Aktivitas:} \textit{\dots shitari, \dots ittari} (Bisa berarti: melakukan X, pergi ke Y, dll).
    \item \textbf{Menyatakan Kesukaan/Kelemahan:} \textit{Nigate} (Kurang suka / Nggak tahan) dan \textit{Daijoubu} (Tahan / Aman).
    \item \textbf{Harus melakukan sesuatu:} \textit{Kaeranakucha!} (Harus pulang!).
    \item \textbf{Berpisah dengan teman:} \textit{Mata renraku suru ne} (Nanti aku kabari / hubungi lagi ya).
\end{itemize}
Coba deh praktekkan ngobrol dengan gaya kasual ini ke teman belajarmu. Pasti asik banget!
\end{tcolorbox}

\newpage
\chapter{Meeting a Japanese Teacher Online}

% ================= PERCAKAPAN 1-6 =================
\begin{convobox}{1}
    \noindent
    \jpkana{あっ、}{A,,} 
    \jpkana{こんにちは。}{konnichiwa.} 

    \vspace{1.5em}
    \noindent{\Large \color{articolor} \textbf{Arti:} Ah, halo (selamat siang).}
\end{convobox}

\begin{convobox}{2}
    \noindent
    \jpword{わたし}{私}{Watashi} \jpkana{の}{no} 
    \jpword{こえ}{声}{koe}\jpkana{、}{,} 
    \jpword{き}{聞}{ki}\jpkana{こえますか？}{koemasu ka?}

    \vspace{1.5em}
    \noindent{\Large \color{articolor} \textbf{Arti:} Apakah suaraku terdengar?}
\end{convobox}

\begin{convobox}{3}
    \noindent
    \jpkana{はい、}{Hai,} 
    \jpword{き}{聞}{ki}\jpkana{こえます。}{koemasu.}

    \vspace{1.5em}
    \noindent{\Large \color{articolor} \textbf{Arti:} Ya, terdengar.}
\end{convobox}

\begin{convobox}{4}
    \noindent
    \jpkana{はじめまして、}{Hajimemashite,} 
    \jpkana{Tanaka}{Tanaka} 
    \jpkana{と}{to} 
    \jpword{もう}{申}{mou}\jpkana{します。}{shimasu.}

    \vspace{1.5em}
    \noindent{\Large \color{articolor} \textbf{Arti:} Salam kenal, namaku Tanaka.}
\end{convobox}

\begin{convobox}{5}
    \noindent
    \jpkana{えっと、}{Eetto,} 
    \jpkana{お}{o}\jpword{なまえ}{名前}{namae} \jpkana{は}{wa} 
    \jpword{なん}{何}{nan} \jpkana{ですか？}{desu ka?}

    \vspace{1.5em}
    \noindent{\Large \color{articolor} \textbf{Arti:} Emm, nama kamu siapa?}
\end{convobox}

\begin{convobox}{6}
    \noindent
    \jpkana{クロワッサン}{Kurowassan} 
    \jpkana{と}{to} 
    \jpword{もう}{申}{mou}\jpkana{します。}{shimasu.} 
    \jpkana{よろしくお}{Yoroshiku o}\jpword{ねが}{願}{nega}\jpkana{いします。}{ishimasu.}

    \vspace{1.5em}
    \noindent{\Large \color{articolor} \textbf{Arti:} Namaku Croissant. Mohon bantuannya (salam kenal).}
\end{convobox}

% --- LESSON BUFFER 1 ---
\begin{lessonbox}{Catatan Belajar: Halo, Kenalan Yuk!}
    Halo adik-adik yang pintar! Kalau kita baru pertama kali bertemu seseorang, kita mengucapkan \textbf{はじめまして (Hajimemashite)} yang artinya "Salam kenal". \\
    
    Nah, untuk menyebutkan nama, biasanya kita pakai "Namaku... desu". Tapi, guru Tanaka menggunakan kata \textbf{と申します (to moushimasu)}. Wah, apa itu? Itu adalah cara yang saaaangat sopan untuk menyebut nama kita! Jadi, kalau mau terdengar sangat keren dan sopan, kamu bisa bilang: \textit{"(Namamu) to moushimasu!"}. Jangan lupa akhiri perkenalanmu dengan \textit{Yoroshiku onegaishimasu} ya!
\end{lessonbox}

% ================= PERCAKAPAN 7-13 =================
\begin{convobox}{7}
    \noindent
    \jpkana{よろしくお}{Yoroshiku o}\jpword{ねが}{願}{nega}\jpkana{いします。}{ishimasu.} 

    \vspace{1.5em}
    \noindent{\Large \color{articolor} \textbf{Arti:} Salam kenal juga.}
\end{convobox}

\begin{convobox}{8}
    \noindent
    \jpkana{オンライン}{Onrain} \jpkana{の}{no} 
    \jpword{にほんご}{日本語}{Nihongo} 
    \jpkana{レッスン}{ressun} \jpkana{は、}{wa,} 
    \jpword{はじ}{初}{haji}\jpkana{めてですか？}{mete desu ka?}

    \vspace{1.5em}
    \noindent{\Large \color{articolor} \textbf{Arti:} Apakah ini pertama kalinya ikut pelajaran bahasa Jepang online?}
\end{convobox}

\begin{convobox}{9}
    \noindent
    \jpkana{はい、}{Hai,} 
    \jpword{はじ}{初}{haji}\jpkana{めてです。}{mete desu.}

    \vspace{1.5em}
    \noindent{\Large \color{articolor} \textbf{Arti:} Ya, ini pertama kalinya.}
\end{convobox}

\begin{convobox}{10}
    \noindent
    \jpkana{そうですか。}{Sou desu ka.} 
    \jpkana{お}{o}\jpword{あ}{会}{a}\jpkana{いできて}{i dekite} 
    \jpkana{うれしいです。}{ureshii desu.}

    \vspace{1.5em}
    \noindent{\Large \color{articolor} \textbf{Arti:} Oh begitu. Aku senang bisa bertemu denganmu.}
\end{convobox}

\begin{convobox}{11}
    \noindent
    \jpkana{いつから}{Itsu kara} 
    \jpword{にほんご}{日本語}{Nihongo} \jpkana{の}{no} 
    \jpword{べんきょう}{勉強}{benkyou} \jpkana{を}{o} 
    \jpkana{していますか？}{shite imasu ka?}

    \vspace{1.5em}
    \noindent{\Large \color{articolor} \textbf{Arti:} Sejak kapan kamu belajar bahasa Jepang?}
\end{convobox}

\begin{convobox}{12}
    \noindent
    \jpword{にねんまえ}{2年前}{Ni-nen mae} \jpkana{から}{kara} 
    \jpword{べんきょう}{勉強}{benkyou} \jpkana{しています。}{shite imasu.}

    \vspace{1.5em}
    \noindent{\Large \color{articolor} \textbf{Arti:} Aku belajar sejak 2 tahun yang lalu.}
\end{convobox}

\begin{convobox}{13}
    \noindent
    \jpkana{どうして}{Doushite} 
    \jpword{にほんご}{日本語}{Nihongo} \jpkana{の}{no} 
    \jpword{べんきょう}{勉強}{benkyou} \jpkana{を}{o} 
    \jpword{はじ}{始}{haji}\jpkana{めたんですか？}{meta n desu ka?}

    \vspace{1.5em}
    \noindent{\Large \color{articolor} \textbf{Arti:} Kenapa kamu mulai belajar bahasa Jepang?}
\end{convobox}

% --- LESSON BUFFER 2 ---
\begin{lessonbox}{Catatan Belajar: Bertanya "Kapan?" dan "Kenapa?"}
    Saat kita penasaran, kita bisa bertanya loh! Di percakapan ini, Guru Tanaka bertanya:
    \begin{itemize}
        \item \textbf{いつ (Itsu)} = Kapan?
        \item \textbf{どうして (Doushite)} = Kenapa?
    \end{itemize}
    Coba bayangkan bahasa Jepang itu seperti main tebak-tebakan. Kalau gurumu tanya \textit{"Doushite?"} (Kenapa), nanti kamu harus jawab dengan alasanmu! Yuk kita lihat apa alasannya di kotak selanjutnya!
\end{lessonbox}

% ================= PERCAKAPAN 14-20 =================
\begin{convobox}{14}
    \noindent
    \jpword{にほん}{日本}{Nihon} \jpkana{の}{no} 
    \jpword{まんが}{漫画}{manga} \jpkana{が}{ga} 
    \jpword{す}{好}{su}\jpkana{きだからです。}{ki da kara desu.}

    \vspace{1.5em}
    \noindent{\Large \color{articolor} \textbf{Arti:} Karena aku suka komik Jepang.}
\end{convobox}

\begin{convobox}{15}
    \noindent
    \jpword{にほん}{日本}{Nihon} \jpkana{に}{ni} 
    \jpword{き}{来}{ki}\jpkana{たことは}{ta koto wa} 
    \jpkana{ありますか？}{arimasu ka?}

    \vspace{1.5em}
    \noindent{\Large \color{articolor} \textbf{Arti:} Apakah kamu pernah datang ke Jepang?}
\end{convobox}

\begin{convobox}{16}
    \noindent
    \jpkana{はい、}{Hai,} 
    \jpword{いっかい}{1回}{ikkai} \jpkana{あります。}{arimasu.} 
    \jpword{とうきょう}{東京}{Toukyou} \jpkana{と}{to} 
    \jpword{おおさか}{大阪}{Oosaka} \jpkana{と}{to} 
    \jpword{なら}{奈良}{Nara} \jpkana{に}{ni} 
    \jpword{い}{行}{i}\jpkana{きました。}{kimashita.}

    \vspace{1.5em}
    \noindent{\Large \color{articolor} \textbf{Arti:} Ya, pernah 1 kali. Aku pergi ke Tokyo, Osaka, dan Nara.}
\end{convobox}

\begin{convobox}{17}
    \noindent
    \jpword{ふだん}{普段}{Fudan} 
    \jpkana{どうやって}{dou yatte} 
    \jpword{にほんご}{日本語}{Nihongo} \jpkana{を}{o} 
    \jpword{べんきょう}{勉強}{benkyou} \jpkana{していますか？}{shite imasu ka?}

    \vspace{1.5em}
    \noindent{\Large \color{articolor} \textbf{Arti:} Biasanya bagaimana cara kamu belajar bahasa Jepang?}
\end{convobox}

\begin{convobox}{18}
    \noindent
    \jpword{にほんご}{日本語}{Nihongo} \jpkana{の}{no} 
    \jpword{まんが}{漫画}{manga} \jpkana{を}{o} 
    \jpword{よ}{読}{yo}\jpkana{んで}{nde} 
    \jpword{べんきょう}{勉強}{benkyou} \jpkana{しています。}{shite imasu.}

    \vspace{1.5em}
    \noindent{\Large \color{articolor} \textbf{Arti:} Aku belajar dengan membaca komik bahasa Jepang.}
\end{convobox}

\begin{convobox}{19}
    \noindent
    \jpkana{それから、}{Sore kara,} 
    \jpkana{ときどき}{tokidoki} 
    \jpkana{YouTube}{YouTube} \jpkana{を}{o} 
    \jpword{み}{観}{mi}\jpkana{て}{te}

    \vspace{1.5em}
    \noindent{\Large \color{articolor} \textbf{Arti:} Lalu, kadang-kadang aku menonton YouTube}
\end{convobox}

\begin{convobox}{20}
    \noindent
    \jpword{かいわ}{会話}{kaiwa} \jpkana{の}{no} 
    \jpword{れんしゅう}{練習}{renshuu} \jpkana{を}{o} 
    \jpkana{しています。}{shite imasu.}

    \vspace{1.5em}
    \noindent{\Large \color{articolor} \textbf{Arti:} dan latihan percakapan.}
\end{convobox}

% --- LESSON BUFFER 3 ---
\begin{lessonbox}{Catatan Belajar: Menjawab "Karena..."}
    Nah, ini dia jawaban dari \textit{"Doushite?"} (Kenapa?) tadi! Kalau kamu mau bilang "Karena...", kamu tinggal menambahkan \textbf{〜から (kara)} di akhir kalimatmu. \\
    
    Contohnya: \textit{Manga ga suki da \textbf{kara} desu!} (Karena suka komik!). \\
    Gampang banget kan? Huruf-huruf di atas mungkin terlihat banyak, tapi pelan-pelan saja bacanya. Romaji di bawahnya akan menuntunmu!
\end{lessonbox}

% ================= PERCAKAPAN 21-27 =================
\begin{convobox}{21}
    \noindent
    \jpkana{レッスン}{Ressun} \jpkana{で}{de} 
    \jpkana{どんなことを}{donna koto o} 
    \jpkana{したいですか？}{shitai desu ka?}

    \vspace{1.5em}
    \noindent{\Large \color{articolor} \textbf{Arti:} Di pelajaran nanti, apa yang ingin kamu lakukan?}
\end{convobox}

\begin{convobox}{22}
    \noindent
    \jpword{かいわ}{会話}{Kaiwa} \jpkana{の}{no} 
    \jpword{れんしゅう}{練習}{renshuu} \jpkana{を}{o} 
    \jpkana{したいですか？}{shitai desu ka?}

    \vspace{1.5em}
    \noindent{\Large \color{articolor} \textbf{Arti:} Apakah ingin berlatih percakapan?}
\end{convobox}

\begin{convobox}{23}
    \noindent
    \jpkana{それとも、}{Sore tomo,} 
    \jpword{ぶんぽう}{文法}{bunpou} \jpkana{を}{o} 
    \jpword{べんきょう}{勉強}{benkyou} \jpkana{したいですか？}{shitai desu ka?}

    \vspace{1.5em}
    \noindent{\Large \color{articolor} \textbf{Arti:} Atau, apakah ingin belajar tata bahasa?}
\end{convobox}

\begin{convobox}{24}
    \noindent
    \jpword{かいわ}{会話}{Kaiwa} \jpkana{の}{no} 
    \jpword{れんしゅう}{練習}{renshuu} \jpkana{を}{o} 
    \jpkana{たくさん}{takusan} 
    \jpkana{したいです。}{shitai desu.}

    \vspace{1.5em}
    \noindent{\Large \color{articolor} \textbf{Arti:} Aku ingin banyak berlatih percakapan.}
\end{convobox}

\begin{convobox}{25}
    \noindent
    \jpkana{あと、}{Ato,} 
    \jpword{しゅくだい}{宿題}{shukudai} \jpkana{も}{mo} 
    \jpkana{たくさん}{takusan} 
    \jpword{だ}{出}{da}\jpkana{してほしいです。}{shite hoshii desu.}

    \vspace{1.5em}
    \noindent{\Large \color{articolor} \textbf{Arti:} Terus, aku ingin diberikan banyak PR juga.}
\end{convobox}

\begin{convobox}{26}
    \noindent
    \jpword{きょうかしょ}{教科書}{Kyoukasho} \jpkana{は、}{wa,} 
    \jpkana{何か}{nani ka} 
    \jpword{も}{持}{mo}\jpkana{っていますか？}{tte imasu ka?}

    \vspace{1.5em}
    \noindent{\Large \color{articolor} \textbf{Arti:} Apakah kamu punya buku pelajaran?}
\end{convobox}

\begin{convobox}{27}
    \noindent
    \jpkana{いえ、}{Ie,} 
    \jpword{なに}{何}{nani}\jpkana{も}{mo} 
    \jpword{も}{持}{mo}\jpkana{っていません。}{tte imasen.} 
    \jpkana{おすすめの}{Osusume no} 
    \jpword{きょうかしょ}{教科書}{kyoukasho} \jpkana{は}{wa} 
    \jpkana{ありますか？}{arimasu ka?}

    \vspace{1.5em}
    \noindent{\Large \color{articolor} \textbf{Arti:} Tidak, aku tidak punya apa-apa. Apakah ada buku pelajaran yang direkomendasikan?}
\end{convobox}

% --- LESSON BUFFER 4 ---
\begin{lessonbox}{Catatan Belajar: "Aku Ingin..."}
    Tahukah kamu cara bilang "Aku ingin..." dalam bahasa Jepang? Gampang banget! Kita cukup menambahkan kata \textbf{〜たい (〜tai)} di belakang kegiatannya!\\
    
    Contoh:
    \begin{itemize}
        \item Belajar = \textit{Benkyou shimasu} $\rightarrow$ Ingin belajar = \textbf{Benkyou shitai desu}.
    \end{itemize}
    Wow, teman kita ini rajin sekali ya, sampai minta banyak \textbf{宿題 (Shukudai)} alias PR! Adik-adik suka PR juga nggak nih?
\end{lessonbox}

% ================= PERCAKAPAN 28-35 =================
\begin{convobox}{28}
    \noindent
    \jpkana{そろそろ}{Sorosoro} 
    \jpword{じかん}{時間}{jikan} \jpkana{なので、}{na node,} 
    \jpword{きょう}{今日}{kyou} \jpkana{は}{wa} 
    \jpkana{ここまでにしましょう。}{koko made ni shimashou.}

    \vspace{1.5em}
    \noindent{\Large \color{articolor} \textbf{Arti:} Karena sudah hampir waktunya, hari ini kita akhiri sampai di sini ya.}
\end{convobox}

\begin{convobox}{29}
    \noindent
    \jpword{さいご}{最後}{Saigo} \jpkana{に、}{ni,} 
    \jpkana{何か}{nani ka} 
    \jpword{しつもん}{質問}{shitsumon} \jpkana{は}{wa} 
    \jpkana{ありますか？}{arimasu ka?}

    \vspace{1.5em}
    \noindent{\Large \color{articolor} \textbf{Arti:} Terakhir, apakah ada pertanyaan?}
\end{convobox}

\begin{convobox}{30}
    \noindent
    \jpkana{いえ、}{Ie,} 
    \jpword{だいじょうぶ}{大丈夫}{daijoubu} \jpkana{です。}{desu.}

    \vspace{1.5em}
    \noindent{\Large \color{articolor} \textbf{Arti:} Tidak, sudah cukup (tidak ada).}
\end{convobox}

\begin{convobox}{31}
    \noindent
    \jpword{きょう}{今日}{Kyou} \jpkana{の}{no} 
    \jpkana{レッスン}{ressun} \jpkana{は、}{wa,} 
    \jpkana{どうでしたか？}{dou deshita ka?} 
    \jpword{むずか}{難}{muzuka}\jpkana{しかったですか？}{shikatta desu ka?}

    \vspace{1.5em}
    \noindent{\Large \color{articolor} \textbf{Arti:} Pelajaran hari ini bagaimana? Apakah sulit?}
\end{convobox}

\begin{convobox}{32}
    \noindent
    \jpword{にほんご}{日本語}{Nihongo} \jpkana{で}{de} 
    \jpword{はな}{話}{hana}\jpkana{すのは}{su no wa} 
    \jpword{むずか}{難}{muzuka}\jpkana{しかったですが、}{shikatta desu ga,} 
    \jpword{たの}{楽}{tano}\jpkana{しかったです。}{shikatta desu.}

    \vspace{1.5em}
    \noindent{\Large \color{articolor} \textbf{Arti:} Berbicara dengan bahasa Jepang itu sulit, tapi menyenangkan.}
\end{convobox}

\begin{convobox}{33}
    \noindent
    \jpword{じかい}{次回}{Jikai} \jpkana{の}{no} 
    \jpkana{レッスン}{ressun} \jpkana{は、}{wa,} 
    \jpword{なんようび}{何曜日}{nan-youbi} \jpkana{の}{no} 
    \jpword{なんじ}{何時}{nan-ji} \jpkana{ごろが}{goro ga} 
    \jpkana{いいですか？}{ii desu ka?}

    \vspace{1.5em}
    \noindent{\Large \color{articolor} \textbf{Arti:} Pelajaran berikutnya, enaknya hari apa dan kira-kira jam berapa?}
\end{convobox}

\begin{convobox}{34}
    \noindent
    \jpword{へいじつ}{平日}{Heijitsu} \jpkana{は}{wa} 
    \jpword{しごと}{仕事}{shigoto} \jpkana{なので、}{na node,} 
    \jpword{どようび}{土曜日}{doyoubi} \jpkana{の}{no} 
    \jpword{あさ}{朝}{asa} \jpword{じゅうじ}{10時}{juu-ji} \jpkana{でもいいですか？}{demo ii desu ka?}

    \vspace{1.5em}
    \noindent{\Large \color{articolor} \textbf{Arti:} Karena hari biasa aku bekerja, apakah boleh hari Sabtu pagi jam 10?}
\end{convobox}

\begin{convobox}{35}
    \noindent
    \jpkana{わかりました。}{Wakarimashita.} 
    \jpkana{では、}{Dewa,} 
    \jpkana{また}{mata} 
    \jpword{じかい}{次回}{jikai} \jpkana{の}{no} 
    \jpkana{レッスン}{ressun} \jpkana{で}{de} 
    \jpkana{お}{o}\jpword{あ}{会}{a}\jpkana{いしましょう。}{i shimashou.} 
    \jpword{しつれい}{失礼}{shitsurei} \jpkana{します。}{shimasu.}

    \vspace{1.5em}
    \noindent{\Large \color{articolor} \textbf{Arti:} Aku mengerti. Kalau begitu, sampai jumpa lagi di pelajaran berikutnya. Permisi.}
\end{convobox}

% --- SUMMARY BOX ---
\vspace{2em}
\begin{tcolorbox}[
    colback=blue!5!white,
    colframe=blue!80!black,
    boxrule=3pt,
    arc=5mm,
    title={\huge\bfseries \sffamily 🌟 Rangkuman Bab 4 🌟},
    fonttitle=\color{white},
    attach boxed title to top center={yshift=-4mm},
    boxed title style={colback=blue!80!black, rounded corners},
    left=5mm, right=5mm, top=7mm, bottom=5mm
]
\Large \textbf{Hebat sekali! Kamu sudah berani ikut kelas online!} Mari kita ingat-ingat lagi apa saja yang sudah kita pelajari:
\begin{itemize}
    \item \textbf{Menyebutkan nama dengan super sopan:} \textit{\dots to moushimasu} (Namaku adalah\dots)
    \item \textbf{Kata tanya penting:} \textit{Itsu} (Kapan) dan \textit{Doushite} (Kenapa).
    \item \textbf{Menjawab "Karena":} \textit{\dots kara desu}.
    \item \textbf{Bilang "Aku ingin\dots":} Tambahkan akhiran \textit{\dots tai desu}.
    \item \textbf{Menanyakan hari dan jam:} \textit{Nan-youbi?} (Hari apa?) dan \textit{Nan-ji?} (Jam berapa?).
\end{itemize}
Semakin lama belajar, rasanya semakin seru ya! Terus berlatih agar lidahmu semakin terbiasa mengucapkan huruf-huruf bahasa Jepang. Semangat, kamu pasti bisa!
\end{tcolorbox}

\newpage
\chapter{Let's Talk about Your Childhood}

% ================= PERCAKAPAN 1-5 =================
\begin{convobox}{1}
    \noindent
    \jpkana{こんにちは。}{Konnichiwa.}

    \vspace{1.5em}
    \noindent{\Large \color{articolor} \textbf{Arti:} Halo (selamat siang).}
\end{convobox}

\begin{convobox}{2}
    \noindent
    \jpword{ひさ}{久}{Hisa}\jpkana{しぶりに}{shiburi ni} 
    \jpword{こうえん}{公園}{kouen} \jpkana{に}{ni} 
    \jpword{き}{来}{ki}\jpkana{ました。}{mashita.}

    \vspace{1.5em}
    \noindent{\Large \color{articolor} \textbf{Arti:} Sudah lama tidak datang ke taman.}
\end{convobox}

\begin{convobox}{3}
    \noindent
    \jpword{きょう}{今日}{Kyou} \jpkana{は}{wa} 
    \jpkana{ここで、}{koko de,} 
    \jpword{こ}{子}{ko}\jpkana{どものころの}{domo no koro no} 
    \jpword{はなし}{話}{hanashi} \jpkana{をしましょう。}{o shimashou.}

    \vspace{1.5em}
    \noindent{\Large \color{articolor} \textbf{Arti:} Hari ini di sini, mari kita bicarakan tentang masa kecil.}
\end{convobox}

\begin{convobox}{4}
    \noindent
    \jpword{こ}{子}{ko}\jpkana{どものころ、}{domo no koro,} 
    \jpkana{どこに}{doko ni} 
    \jpword{す}{住}{su}\jpkana{んでいましたか？}{nde imashita ka?}

    \vspace{1.5em}
    \noindent{\Large \color{articolor} \textbf{Arti:} Sewaktu kecil, kamu tinggal di mana?}
\end{convobox}

\begin{convobox}{5}
    \noindent
    \jpword{おおさか}{大阪}{Oosaka} \jpkana{に}{ni} 
    \jpword{す}{住}{su}\jpkana{んでいました。}{nde imashita.}

    \vspace{1.5em}
    \noindent{\Large \color{articolor} \textbf{Arti:} Dulu tinggal di Osaka.}
\end{convobox}

% --- LESSON BUFFER 1 ---
\begin{lessonbox}{Catatan Belajar: Tata Bahasa "Masa Lalu" (Bentuk Lampau)}
    Hai adik-adik! Pernahkah kamu bercerita tentang kejadian waktu kamu masih kecil? Di bahasa Jepang, kita pakai kata ajaib \textbf{子どものころ (Kodomo no koro)} yang artinya "Sewaktu masih kecil".\\
    
    Nah, karena ini cerita zaman dulu, perhatikan deh kalimat nomor 4 dan 5. Ada kata \textbf{〜ていました (〜te imashita)}. Ini adalah grammar untuk menyatakan \textbf{keadaan yang berlangsung di masa lalu}.
    \begin{itemize}
        \item \textit{Sunde imasu} (Sekarang sedang tinggal) $\rightarrow$ \textbf{Sunde imashita} (Dulu tinggal).
    \end{itemize}
    Coba ingat-ingat, waktu kecil kamu \textit{sunde imashita} di mana?
\end{lessonbox}

% ================= PERCAKAPAN 6-10 =================
\begin{convobox}{6}
    \noindent
    \jpword{なに}{何}{Nani} \jpkana{を}{o} 
    \jpkana{して}{shite} 
    \jpword{あそ}{遊}{aso}\jpkana{ぶのが}{bu no ga} 
    \jpword{す}{好}{su}\jpkana{きでしたか？}{ki deshita ka?}

    \vspace{1.5em}
    \noindent{\Large \color{articolor} \textbf{Arti:} Dulu suka bermain apa?}
\end{convobox}

\begin{convobox}{7}
    \noindent
    \jpword{こうえん}{公園}{Kouen} \jpkana{で}{de} 
    \jpkana{サッカー}{sakkaa} \jpkana{を}{o} 
    \jpkana{して}{shite} 
    \jpword{あそ}{遊}{aso}\jpkana{ぶのが}{bu no ga} 
    \jpword{す}{好}{su}\jpkana{きでした。}{ki deshita.}

    \vspace{1.5em}
    \noindent{\Large \color{articolor} \textbf{Arti:} Dulu suka bermain sepak bola di taman.}
\end{convobox}

\begin{convobox}{8}
    \noindent
    \jpword{そと}{外}{Soto} \jpkana{で}{de} 
    \jpword{あそ}{遊}{aso}\jpkana{ぶのと、}{bu no to,} 
    \jpword{いえ}{家}{ie} \jpkana{の}{no} 
    \jpword{なか}{中}{naka} \jpkana{で}{de} 
    \jpword{あそ}{遊}{aso}\jpkana{ぶの、}{bu no,}

    \vspace{1.5em}
    \noindent{\Large \color{articolor} \textbf{Arti:} Antara bermain di luar dan bermain di dalam rumah,}
\end{convobox}

\begin{convobox}{9}
    \noindent
    \jpkana{どっちの}{Dotchi no} 
    \jpword{ほう}{方}{hou} \jpkana{が}{ga} 
    \jpword{す}{好}{su}\jpkana{きでしたか？}{ki deshita ka?}

    \vspace{1.5em}
    \noindent{\Large \color{articolor} \textbf{Arti:} kamu lebih suka yang mana?}
\end{convobox}

\begin{convobox}{10}
    \noindent
    \jpword{そと}{外}{Soto} \jpkana{で}{de} 
    \jpword{あそ}{遊}{aso}\jpkana{ぶ}{bu} 
    \jpword{ほう}{方}{hou} \jpkana{が}{ga} 
    \jpword{す}{好}{su}\jpkana{きでした。}{ki deshita.}

    \vspace{1.5em}
    \noindent{\Large \color{articolor} \textbf{Arti:} Dulu aku lebih suka bermain di luar.}
\end{convobox}

% --- LESSON BUFFER 2 ---
\begin{lessonbox}{Catatan Belajar: Tata Bahasa Memilih 2 Pilihan (A vs B)}
    Bagaimana cara bertanya "Kamu lebih suka A atau B?" dalam bahasa Jepang? Kita bisa gunakan rumus grammar yang super keren ini:\\
    
    \textbf{[Pilihan A] to, [Pilihan B] to, dotchi no hou ga [kata sifat] desu ka?}\\
    
    Contoh di atas: Bermain di luar (\textit{Soto de asobu no}) VS Bermain di dalam (\textit{Ie no naka de asobu no}). Kata \textbf{どっち (Dotchi)} artinya "Yang mana". Untuk menjawabnya, sebutkan pilihan yang kamu suka dan tambahkan \textbf{〜ほうが (〜hou ga)}. Wah, kamu semakin jago merangkai kalimat panjang!
\end{lessonbox}

% ================= PERCAKAPAN 11-18 =================
\begin{convobox}{11}
    \noindent
    \jpword{こ}{子}{ko}\jpkana{どものころ、}{domo no koro,} 
    \jpkana{ペット}{petto} \jpkana{を}{o} 
    \jpword{か}{飼}{ka}\jpkana{っていましたか？}{tte imashita ka?}

    \vspace{1.5em}
    \noindent{\Large \color{articolor} \textbf{Arti:} Waktu kecil, apakah kamu memelihara hewan peliharaan?}
\end{convobox}

\begin{convobox}{12}
    \noindent
    \jpkana{はい。}{Hai.} 
    \jpword{きんぎょ}{金魚}{Kingyo} \jpkana{を}{o} 
    \jpword{か}{飼}{ka}\jpkana{っていました。}{tte imashita.}

    \vspace{1.5em}
    \noindent{\Large \color{articolor} \textbf{Arti:} Ya. Aku memelihara ikan mas koki.}
\end{convobox}

\begin{convobox}{13}
    \noindent
    \jpword{こ}{子}{ko}\jpkana{どものころ、}{domo no koro,} 
    \jpword{す}{好}{su}\jpkana{きだった}{ki datta} 
    \jpkana{アニメ}{anime} \jpkana{は}{wa} 
    \jpword{なん}{何}{nan} \jpkana{でしたか？}{deshita ka?}

    \vspace{1.5em}
    \noindent{\Large \color{articolor} \textbf{Arti:} Waktu kecil, anime apa yang kamu sukai?}
\end{convobox}

\begin{convobox}{14}
    \noindent
    \jpkana{「キャプテン}{"Kyaputen} \jpword{つばさ}{翼}{Tsubasa}\jpkana{」という}{" to iu} 
    \jpkana{アニメ}{anime} \jpkana{が}{ga} 
    \jpword{す}{好}{su}\jpkana{きでした。}{ki deshita.}

    \vspace{1.5em}
    \noindent{\Large \color{articolor} \textbf{Arti:} Aku dulu suka anime bernama 'Captain Tsubasa'.}
\end{convobox}

\begin{convobox}{15}
    \noindent
    \jpword{こ}{子}{ko}\jpkana{どものころ、}{domo no koro,} 
    \jpword{にがて}{苦手}{nigate} \jpkana{だった}{datta} 
    \jpword{た}{食}{ta}\jpkana{べ}{be}\jpword{もの}{物}{mono} \jpkana{は}{wa} 
    \jpkana{ありましたか？}{arimashita ka?}

    \vspace{1.5em}
    \noindent{\Large \color{articolor} \textbf{Arti:} Waktu kecil, apakah ada makanan yang tidak kamu sukai?}
\end{convobox}

\begin{convobox}{16}
    \noindent
    \jpkana{ピーマン}{Piiman} \jpkana{と}{to} 
    \jpkana{にんじん}{ninjin} \jpkana{が}{ga} 
    \jpword{にがて}{苦手}{nigate} \jpkana{でした。}{deshita.}

    \vspace{1.5em}
    \noindent{\Large \color{articolor} \textbf{Arti:} Dulu aku tidak suka paprika hijau dan wortel.}
\end{convobox}

\begin{convobox}{17}
    \noindent
    \jpword{こ}{子}{ko}\jpkana{どものころ、}{domo no koro,} 
    \jpword{だれ}{誰}{dare} \jpkana{と}{to} 
    \jpkana{よく}{yoku} 
    \jpword{あそ}{遊}{aso}\jpkana{んでいましたか？}{nde imashita ka?}

    \vspace{1.5em}
    \noindent{\Large \color{articolor} \textbf{Arti:} Waktu kecil, sering bermain dengan siapa?}
\end{convobox}

\begin{convobox}{18}
    \noindent
    \jpword{きんじょ}{近所}{Kinjo} \jpkana{に}{ni} 
    \jpword{す}{住}{su}\jpkana{んでいた、}{nde ita,} 
    \jpkana{たいやきちゃんと}{Taiyaki-chan to} 
    \jpkana{よく}{yoku} 
    \jpword{あそ}{遊}{aso}\jpkana{んでいました。}{nde imashita.}

    \vspace{1.5em}
    \noindent{\Large \color{articolor} \textbf{Arti:} Aku dulu sering bermain dengan Taiyaki-chan yang tinggal di lingkungan sekitar.}
\end{convobox}

% --- LESSON BUFFER 3 ---
\begin{lessonbox}{Catatan Belajar: Mengubah Sifat Menjadi Masa Lalu!}
    Tahukah kamu? Kata sifat dan kata benda dalam bahasa Jepang akan berubah bentuk kalau peristiwanya sudah lewat lho!
    \begin{itemize}
        \item \textbf{Bentuk Sekarang:} \textit{Suki desu} (Suka) $\rightarrow$ \textbf{Bentuk Lampau:} \textit{Suki \textbf{deshita}} (Dulu suka).
        \item \textbf{Bentuk Sekarang:} \textit{Nigate desu} (Tidak suka/lemah) $\rightarrow$ \textbf{Bentuk Lampau:} \textit{Nigate \textbf{deshita}} (Dulu tidak suka).
    \end{itemize}
    Jadi ingat ya, kalau kamu mau cerita kesukaanmu pas masih TK atau SD, pakailah kata \textbf{〜でした (〜deshita)} di belakangnya!
\end{lessonbox}

% ================= PERCAKAPAN 19-26 =================
\begin{convobox}{19}
    \noindent
    \jpword{がっこう}{学校}{Gakkou} \jpkana{に}{ni} 
    \jpword{い}{行}{i}\jpkana{くのは}{ku no wa} 
    \jpword{す}{好}{su}\jpkana{きでしたか？}{ki deshita ka?} 
    \jpkana{きらいでしたか？}{Kirai deshita ka?}

    \vspace{1.5em}
    \noindent{\Large \color{articolor} \textbf{Arti:} Apakah kamu suka pergi ke sekolah? Atau benci?}
\end{convobox}

\begin{convobox}{20}
    \noindent
    \jpword{す}{好}{su}\jpkana{きでしたが、}{ki deshita ga,} 
    \jpword{がっこう}{学校}{gakkou} \jpkana{に}{ni} 
    \jpword{い}{行}{i}\jpkana{きたくない}{kitakunai} 
    \jpword{ひ}{日}{hi} \jpkana{も}{mo} 
    \jpkana{ありました。}{arimashita.}

    \vspace{1.5em}
    \noindent{\Large \color{articolor} \textbf{Arti:} Aku suka, tapi ada juga hari-hari di mana aku tidak ingin pergi ke sekolah.}
\end{convobox}

\begin{convobox}{21}
    \noindent
    \jpword{べんきょう}{勉強}{Benkyou} \jpkana{は}{wa} 
    \jpword{とくい}{得意}{tokui} \jpkana{でしたか？}{deshita ka?} 
    \jpword{にがて}{苦手}{nigate} \jpkana{でしたか？}{deshita ka?}

    \vspace{1.5em}
    \noindent{\Large \color{articolor} \textbf{Arti:} Apakah kamu jago belajar? Atau tidak bisa (lemah)?}
\end{convobox}

\begin{convobox}{22}
    \noindent
    \jpkana{ちょっと}{Chotto} 
    \jpword{にがて}{苦手}{nigate} \jpkana{でした。}{deshita.}

    \vspace{1.5em}
    \noindent{\Large \color{articolor} \textbf{Arti:} Sedikit tidak bisa (kurang pandai).}
\end{convobox}

\begin{convobox}{23}
    \noindent
    \jpword{がっこう}{学校}{Gakkou} \jpkana{には}{ni wa} 
    \jpword{なん}{何}{nan} \jpkana{で}{de} 
    \jpword{かよ}{通}{kayo}\jpkana{っていましたか？}{tte imashita ka?}

    \vspace{1.5em}
    \noindent{\Large \color{articolor} \textbf{Arti:} Pergi ke sekolah menggunakan (kendaraan) apa?}
\end{convobox}

\begin{convobox}{24}
    \noindent
    \jpword{たと}{例}{Tato}\jpkana{えば、}{eba,} 
    \jpword{くるま}{車}{kuruma} \jpkana{とか、}{toka,} 
    \jpkana{バス}{basu} \jpkana{とか。}{toka.}

    \vspace{1.5em}
    \noindent{\Large \color{articolor} \textbf{Arti:} Misalnya, dengan mobil, atau bus.}
\end{convobox}

\begin{convobox}{25}
    \noindent
    \jpword{しょうがっこう}{小学校}{Shougakkou} \jpkana{には、}{ni wa,} 
    \jpword{ある}{歩}{aru}\jpkana{いて}{ite} 
    \jpword{かよ}{通}{kayo}\jpkana{っていました。}{tte imashita.}

    \vspace{1.5em}
    \noindent{\Large \color{articolor} \textbf{Arti:} Ke SD, aku berjalan kaki.}
\end{convobox}

\begin{convobox}{26}
    \noindent
    \jpword{ちゅうがっこう}{中学校}{Chuugakkou} \jpkana{には、}{ni wa,} 
    \jpword{じてんしゃ}{自転車}{jitensha} \jpkana{で}{de} 
    \jpword{かよ}{通}{kayo}\jpkana{っていました。}{tte imashita.}

    \vspace{1.5em}
    \noindent{\Large \color{articolor} \textbf{Arti:} Ke SMP, aku naik sepeda.}
\end{convobox}

% --- LESSON BUFFER 4 ---
\begin{lessonbox}{Catatan Belajar: Naik Kendaraan vs Jalan Kaki}
    Coba lihat kalimat nomor 23. Ada grammar partikel \textbf{で (De)}!\\
    Dalam bahasa Jepang, kalau kita mau bilang pergi menggunakan alat transportasi, kita pakai partikel "De". 
    \begin{itemize}
        \item \textit{Kuruma \textbf{de}} (Menggunakan mobil)
        \item \textit{Jitensha \textbf{de}} (Menggunakan sepeda)
    \end{itemize}
    Tapi hati-hati! Kalau kamu pergi **dengan berjalan kaki**, jangan pakai "Ashi de" (dengan kaki) ya! Cukup bilang \textbf{歩いて (Aruite)} yang artinya "berjalan kaki" tanpa perlu partikel apa pun! Hebat kan?
\end{lessonbox}

% ================= PERCAKAPAN 27-36 =================
\begin{convobox}{27}
    \noindent
    \jpword{こ}{子}{ko}\jpkana{どものころ、}{domo no koro,} 
    \jpword{こわ}{怖}{kowa}\jpkana{かった}{katta} 
    \jpword{もの}{物}{mono} \jpkana{は}{wa} 
    \jpkana{ありましたか？}{arimashita ka?}

    \vspace{1.5em}
    \noindent{\Large \color{articolor} \textbf{Arti:} Waktu kecil, apakah ada hal yang menakutkan (buatmu)?}
\end{convobox}

\begin{convobox}{28}
    \noindent
    \jpkana{おばけ}{Obake} \jpkana{と}{to} 
    \jpword{かみなり}{雷}{kaminari} \jpkana{が}{ga} 
    \jpword{こわ}{怖}{kowa}\jpkana{かったです。}{katta desu.}

    \vspace{1.5em}
    \noindent{\Large \color{articolor} \textbf{Arti:} Dulu aku takut hantu dan petir.}
\end{convobox}

\begin{convobox}{29}
    \noindent
    \jpword{こ}{子}{ko}\jpkana{どものころの}{domo no koro no} 
    \jpword{ゆめ}{夢}{yume} \jpkana{は}{wa} 
    \jpword{なん}{何}{nan} \jpkana{でしたか？}{deshita ka?}

    \vspace{1.5em}
    \noindent{\Large \color{articolor} \textbf{Arti:} Waktu kecil, apa mimpimu (cita-cita)?}
\end{convobox}

\begin{convobox}{30}
    \noindent
    \jpword{うちゅうひこうし}{宇宙飛行士}{Uchuuhikoushi} \jpkana{でした。}{deshita.}

    \vspace{1.5em}
    \noindent{\Large \color{articolor} \textbf{Arti:} Dulu cita-citaku astronot.}
\end{convobox}

\begin{convobox}{31}
    \noindent
    \jpword{こ}{子}{ko}\jpkana{どものころの}{domo no koro no} 
    \jpword{じぶん}{自分}{jibun} \jpkana{に、}{ni,} 
    \jpkana{何か}{nani ka} 
    \jpword{い}{言}{i}\jpkana{いたいことは}{itai koto wa} 
    \jpkana{ありますか？}{arimasu ka?}

    \vspace{1.5em}
    \noindent{\Large \color{articolor} \textbf{Arti:} Apakah ada sesuatu yang ingin dikatakan kepada dirimu di masa kecil?}
\end{convobox}

\begin{convobox}{32}
    \noindent
    \jpword{ふとん}{布団}{Futon} \jpkana{の}{no} 
    \jpword{なか}{中}{naka} \jpkana{で}{de} 
    \jpkana{まんが}{manga} \jpkana{を}{o} 
    \jpword{よ}{読}{yo}\jpkana{むと}{mu to} 
    \jpword{め}{目}{me} \jpkana{が}{ga} 
    \jpword{わる}{悪}{waru}\jpkana{くなるよ、}{ku naru yo,} 
    \jpkana{と}{to} \jpword{い}{言}{i}\jpkana{いたいです。}{itai desu.}

    \vspace{1.5em}
    \noindent{\Large \color{articolor} \textbf{Arti:} Aku ingin bilang, "Kalau baca komik di dalam selimut, matamu akan rusak lho".}
\end{convobox}

\begin{convobox}{33}
    \noindent
    \jpword{きょう}{今日}{Kyou} \jpkana{は}{wa} 
    \jpkana{いろいろ}{iroiro} 
    \jpword{はな}{話}{hana}\jpkana{してくれて}{shite kurete} 
    \jpkana{ありがとうございました。}{arigatou gozaimashita.}

    \vspace{1.5em}
    \noindent{\Large \color{articolor} \textbf{Arti:} Terima kasih sudah menceritakan banyak hal hari ini.}
\end{convobox}

\begin{convobox}{34}
    \noindent
    \jpkana{たまには}{Tama ni wa} 
    \jpword{むかし}{昔}{mukashi} \jpkana{の}{no} 
    \jpword{はなし}{話}{hanashi} \jpkana{をするのも}{o suru no mo} 
    \jpkana{いいですね。}{ii desu ne.}

    \vspace{1.5em}
    \noindent{\Large \color{articolor} \textbf{Arti:} Sesekali membicarakan masa lalu itu menyenangkan ya.}
\end{convobox}

\begin{convobox}{35}
    \noindent
    \jpkana{また}{Mata} 
    \jpword{いっしょ}{一緒}{issho} \jpkana{に}{ni} 
    \jpword{はな}{話}{hana}\jpkana{しましょう。}{shimashou.}

    \vspace{1.5em}
    \noindent{\Large \color{articolor} \textbf{Arti:} Mari mengobrol bersama lagi nanti.}
\end{convobox}

\begin{convobox}{36}
    \noindent
    \jpkana{それでは、}{Sore dewa,} 
    \jpkana{また。}{mata.}

    \vspace{1.5em}
    \noindent{\Large \color{articolor} \textbf{Arti:} Kalau begitu, sampai jumpa.}
\end{convobox}

% --- SUMMARY BOX ---
\vspace{2em}
\begin{tcolorbox}[
    colback=blue!5!white,
    colframe=blue!80!black,
    boxrule=3pt,
    arc=5mm,
    title={\huge\bfseries \sffamily 🌟 Rangkuman Bab 5 🌟},
    fonttitle=\color{white},
    attach boxed title to top center={yshift=-4mm},
    boxed title style={colback=blue!80!black, rounded corners},
    left=5mm, right=5mm, top=7mm, bottom=5mm
]
\Large \textbf{Luar biasa! Kini kamu mahir bercerita masa lalu!} Yuk ingat grammar utamanya:
\begin{itemize}
    \item \textbf{Kata Kunci Masa Lalu:} \textit{Kodomo no koro} (Waktu kecil) / \textit{Mukashi} (Zaman dulu).
    \item \textbf{Akhiran Lampau:} Ubah \textit{Desu} menjadi \textbf{Deshita}. (Cth: \textit{Suki deshita} = Dulu suka).
    \item \textbf{Lebih suka yang mana?} Pakai grammar \textit{Dotchi no hou ga\dots} dan jawab dengan \textit{\dots hou ga suki deshita}.
    \item \textbf{Alat Transportasi:} Pakai partikel \textbf{De} (\textit{Basu de, Jitensha de}), kecuali jalan kaki gunakan \textbf{Aruite}.
\end{itemize}
Jangan lupa nasihat dari cerita di atas ya, jangan baca komik dalam gelap nanti matamu sakit! Sampai jumpa di bab selanjutnya yang tidak kalah seru!
\end{tcolorbox}
\newpage
\chapter{How's Your Japanese Learning Going?}

% ================= PERCAKAPAN 1-5 =================
\begin{convobox}{1}
    \noindent
    \jpkana{こんにちは。}{Konnichiwa.}

    \vspace{1.5em}
    \noindent{\Large \color{articolor} \textbf{Arti:} Halo.}
\end{convobox}

\begin{convobox}{2}
    \noindent
    \jpword{にほんご}{日本語}{Nihongo} \jpkana{の}{no} 
    \jpword{べんきょう}{勉強}{benkyou}\jpkana{、}{,} 
    \jpkana{お}{o}\jpword{つか}{疲}{tsuka}\jpkana{れ}{re}\jpword{さま}{様}{sama} \jpkana{です。}{desu.}

    \vspace{1.5em}
    \noindent{\Large \color{articolor} \textbf{Arti:} Terima kasih atas kerja kerasmu belajar bahasa Jepang (Semangat belajarnya).}
\end{convobox}

\begin{convobox}{3}
    \noindent
    \jpword{すこ}{少}{suko}\jpkana{し}{shi} 
    \jpword{きゅう}{休}{kyuu}\jpkana{けいして、}{kei shite,} 
    \jpword{いっしょ}{一緒}{issho} \jpkana{に}{ni} 
    \jpkana{おしゃべりしましょう。}{oshaberi shimashou.}

    \vspace{1.5em}
    \noindent{\Large \color{articolor} \textbf{Arti:} Mari istirahat sebentar dan mengobrol bersama.}
\end{convobox}

\begin{convobox}{4}
    \noindent
    \jpword{さいきん}{最近}{Saikin}\jpkana{、}{,} 
    \jpword{にほんご}{日本語}{Nihongo} \jpkana{の}{no} 
    \jpword{べんきょう}{勉強}{benkyou} \jpkana{は}{wa} 
    \jpkana{どうですか？}{dou desu ka?}

    \vspace{1.5em}
    \noindent{\Large \color{articolor} \textbf{Arti:} Akhir-akhir ini, bagaimana dengan belajar bahasa Jepangnya?}
\end{convobox}

\begin{convobox}{5}
    \noindent
    \jpword{すこ}{少}{suko}\jpkana{し}{shi} 
    \jpword{たいへん}{大変}{taihen} \jpkana{ですが、}{desu ga,} 
    \jpword{じゅんちょう}{順調}{junchou} \jpkana{です。}{desu.}

    \vspace{1.5em}
    \noindent{\Large \color{articolor} \textbf{Arti:} Sedikit berat/susah, tapi berjalan lancar.}
\end{convobox}

% --- LESSON BUFFER 1 ---
\begin{lessonbox}{Catatan Belajar: Kata Ajaib "Otsukaresama"}
    Adik-adik, tahukah kamu kata \textbf{お疲れ様です (Otsukaresama desu)}? Ini adalah kalimat pujian super ajaib di Jepang! Artinya kurang lebih, "Terima kasih atas kerja kerasmu hari ini." \\
    
    Orang Jepang sering sekali mengucapkannya setelah selesai belajar, selesai bekerja, atau melihat temannya berusaha keras. Coba ucapkan ke teman belajarmu ya, pasti mereka akan senang!
\end{lessonbox}

% ================= PERCAKAPAN 6-12 =================
\begin{convobox}{6}
    \noindent
    \jpword{まいにち}{毎日}{Mainichi} 
    \jpkana{どれぐらいの}{dore gurai no} 
    \jpword{じかん}{時間}{jikan} 
    \jpword{べんきょう}{勉強}{benkyou} \jpkana{していますか？}{shite imasu ka?}

    \vspace{1.5em}
    \noindent{\Large \color{articolor} \textbf{Arti:} Setiap hari kamu belajar kira-kira berapa lama (waktunya)?}
\end{convobox}

\begin{convobox}{7}
    \noindent
    \jpword{いちにち}{1日}{Ichi-nichi} 
    \jpword{いちじかん}{1時間}{ichi-jikan} 
    \jpkana{ぐらい}{gurai} 
    \jpword{べんきょう}{勉強}{benkyou} \jpkana{しています。}{shite imasu.}

    \vspace{1.5em}
    \noindent{\Large \color{articolor} \textbf{Arti:} Aku belajar kira-kira 1 jam sehari.}
\end{convobox}

\begin{convobox}{8}
    \noindent
    \jpkana{いつ}{Itsu} 
    \jpword{べんきょう}{勉強}{benkyou} \jpkana{していますか？}{shite imasu ka?}

    \vspace{1.5em}
    \noindent{\Large \color{articolor} \textbf{Arti:} Kapan kamu belajar?}
\end{convobox}

\begin{convobox}{9}
    \noindent
    \jpword{あさ}{朝}{Asa} \jpkana{ですか？}{desu ka?} 
    \jpword{よる}{夜}{Yoru} \jpkana{ですか？}{desu ka?}

    \vspace{1.5em}
    \noindent{\Large \color{articolor} \textbf{Arti:} Apakah pagi hari? Apakah malam hari?}
\end{convobox}

\begin{convobox}{10}
    \noindent
    \jpword{あさ}{朝}{Asa}\jpkana{ごはんを}{gohan o} 
    \jpword{た}{食}{ta}\jpkana{べた}{beta} 
    \jpword{あと}{後}{ato} \jpkana{と、}{to,}

    \vspace{1.5em}
    \noindent{\Large \color{articolor} \textbf{Arti:} Setelah makan sarapan (pagi), dan}
\end{convobox}

\begin{convobox}{11}
    \noindent
    \jpword{よる}{夜}{Yoru} 
    \jpword{ね}{寝}{ne}\jpkana{る}{ru} 
    \jpword{まえ}{前}{mae} \jpkana{に}{ni} 
    \jpword{べんきょう}{勉強}{benkyou} \jpkana{しています。}{shite imasu.}

    \vspace{1.5em}
    \noindent{\Large \color{articolor} \textbf{Arti:} sebelum tidur di malam hari, aku belajar.}
\end{convobox}

\begin{convobox}{12}
    \noindent
    \jpkana{どうやって}{Dou yatte} 
    \jpword{にほんご}{日本語}{Nihongo} \jpkana{を}{o} 
    \jpword{べんきょう}{勉強}{benkyou} \jpkana{を}{o} 
    \jpkana{していますか？}{shite imasu ka?}

    \vspace{1.5em}
    \noindent{\Large \color{articolor} \textbf{Arti:} Bagaimana caramu belajar bahasa Jepang?}
\end{convobox}

% --- LESSON BUFFER 2 ---
\begin{lessonbox}{Catatan Belajar: Sebelum dan Sesudah!}
    Waktu belajar setiap orang beda-beda ya! Ada yang pagi, ada yang malam. Nah, kalau kita mau bilang "Sesudah" dan "Sebelum", kita bisa pakai rumus ini:
    \begin{itemize}
        \item \textbf{〜た後 (〜ta ato)} = Setelah (melakukan sesuatu). Contoh: \textit{Tabeta ato} (Setelah makan).
        \item \textbf{〜る前 (〜ru mae)} = Sebelum (melakukan sesuatu). Contoh: \textit{Neru mae} (Sebelum tidur).
    \end{itemize}
    Coba deh ingat-ingat, kegiatan apa yang biasa kamu lakukan "Neru mae" (sebelum tidur)? Sikat gigi? Baca buku? Atau belajar bahasa Jepang?
\end{lessonbox}

% ================= PERCAKAPAN 13-18 =================
\begin{convobox}{13}
    \noindent
    \jpkana{お}{O}\jpword{き}{気}{ki} \jpkana{に}{ni} \jpword{い}{入}{i}\jpkana{りの}{ri no} 
    \jpword{きょうかしょ}{教科書}{kyoukasho} \jpkana{や、}{ya,} 
    \jpkana{Webサイト}{Web saito} \jpkana{は}{wa} 
    \jpkana{ありますか？}{arimasu ka?}

    \vspace{1.5em}
    \noindent{\Large \color{articolor} \textbf{Arti:} Apakah ada buku pelajaran atau website favorit?}
\end{convobox}

\begin{convobox}{14}
    \noindent
    \jpword{きょうかしょ}{教科書}{Kyoukasho} \jpkana{と}{to} 
    \jpkana{YouTube}{YouTube} \jpkana{を}{o} 
    \jpword{つか}{使}{tsuka}\jpkana{って}{tte} 
    \jpword{べんきょう}{勉強}{benkyou} \jpkana{しています。}{shite imasu.}

    \vspace{1.5em}
    \noindent{\Large \color{articolor} \textbf{Arti:} Aku belajar menggunakan buku pelajaran dan YouTube.}
\end{convobox}

\begin{convobox}{15}
    \noindent
    \jpkana{YouTube}{YouTube} \jpkana{は、}{wa,} 
    \jpkana{たなかさんの}{Tanaka-san no} 
    \jpword{どうが}{動画}{douga} \jpkana{が}{ga} 
    \jpkana{お}{o}\jpword{き}{気}{ki} \jpkana{に}{ni} \jpword{い}{入}{i}\jpkana{りです。}{ri desu.}

    \vspace{1.5em}
    \noindent{\Large \color{articolor} \textbf{Arti:} Di YouTube, video dari Tanaka-san adalah favoritku.}
\end{convobox}

\begin{convobox}{16}
    \noindent
    \jpword{さいきん}{最近}{Saikin} 
    \jpword{おぼ}{覚}{obo}\jpkana{えた}{eta} 
    \jpword{ことば}{言葉}{kotoba} \jpkana{は}{wa} 
    \jpword{なに}{何}{nani}\jpkana{か}{ka} 
    \jpkana{ありますか？}{arimasu ka?}

    \vspace{1.5em}
    \noindent{\Large \color{articolor} \textbf{Arti:} Apakah ada kosakata yang baru-baru ini kamu hafal/ingat?}
\end{convobox}

\begin{convobox}{17}
    \noindent
    \jpkana{「チョキチョキ」という}{"Choki-choki" to iu} 
    \jpword{ことば}{言葉}{kotoba} \jpkana{を}{o} 
    \jpword{おぼ}{覚}{obo}\jpkana{えました。}{emashita.}

    \vspace{1.5em}
    \noindent{\Large \color{articolor} \textbf{Arti:} Aku menghafal kata "Choki-choki".}
\end{convobox}

\begin{convobox}{18}
    \noindent
    \jpkana{はさみで}{Hasami de} 
    \jpword{なに}{何}{nani}\jpkana{か}{ka} \jpkana{を}{o} 
    \jpword{き}{切}{ki}\jpkana{る}{ru} 
    \jpword{おと}{音}{oto} \jpkana{を}{o} 
    \jpword{あらわ}{表}{arawa}\jpkana{す}{su} 
    \jpkana{オノマトペです。}{onomatope desu.}

    \vspace{1.5em}
    \noindent{\Large \color{articolor} \textbf{Arti:} Itu adalah Onomatope yang menggambarkan suara memotong sesuatu dengan gunting.}
\end{convobox}

% --- LESSON BUFFER 3 ---
\begin{lessonbox}{Catatan Belajar: Kata Benda Kesukaan \& Onomatope!}
    Kata \textbf{お気に入り (O-ki ni iri)} artinya "Favorit" atau kesukaan. Jadi kalau kamu mau bilang "Buku favorit", tinggal tambahkan \textbf{no} jadi \textit{O-ki ni iri no hon}! \\
    
    Nah, terus apa itu \textbf{オノマトペ (Onomatope)}? Itu adalah kata tiruan bunyi! Di Indonesia kita punya kata "Dor!" untuk suara tembakan. Di Jepang, suara gunting memotong kertas itu bunyinya adalah \textbf{チョキチョキ (Choki-choki)}! Lucu banget, ya?
\end{lessonbox}

% ================= PERCAKAPAN 19-24 =================
\begin{convobox}{19}
    \noindent
    \jpkana{ふだん}{Fudan} 
    \jpword{ひとり}{一人}{hitori} \jpkana{で}{de} 
    \jpword{べんきょう}{勉強}{benkyou} \jpkana{していますか？}{shite imasu ka?}

    \vspace{1.5em}
    \noindent{\Large \color{articolor} \textbf{Arti:} Biasanya kamu belajar sendirian?}
\end{convobox}

\begin{convobox}{20}
    \noindent
    \jpkana{それとも、}{Sore tomo,} 
    \jpword{せんせい}{先生}{sensei} \jpkana{や}{ya} 
    \jpword{ともだち}{友達}{tomodachi} \jpkana{と}{to} 
    \jpword{いっしょ}{一緒}{issho} \jpkana{に}{ni} 
    \jpword{べんきょう}{勉強}{benkyou} \jpkana{していますか？}{shite imasu ka?}

    \vspace{1.5em}
    \noindent{\Large \color{articolor} \textbf{Arti:} Atau, kamu belajar bersama guru dan teman-teman?}
\end{convobox}

\begin{convobox}{21}
    \noindent
    \jpword{ひとり}{一人}{Hitori} \jpkana{で}{de} 
    \jpword{べんきょう}{勉強}{benkyou} \jpkana{しています。}{shite imasu.}

    \vspace{1.5em}
    \noindent{\Large \color{articolor} \textbf{Arti:} Aku belajar sendirian.}
\end{convobox}

\begin{convobox}{22}
    \noindent
    \jpword{べんきょう}{勉強}{Benkyou} \jpkana{の}{no} 
    \jpkana{やる}{yaru}\jpword{き}{気}{ki} \jpkana{が}{ga} 
    \jpword{で}{出}{de}\jpkana{ない}{nai} 
    \jpword{とき}{時}{toki} \jpkana{は}{wa} 
    \jpkana{どうしていますか？}{dou shite imasu ka?}

    \vspace{1.5em}
    \noindent{\Large \color{articolor} \textbf{Arti:} Apa yang kamu lakukan saat sedang tidak ada motivasi belajar?}
\end{convobox}

\begin{convobox}{23}
    \noindent
    \jpword{べんきょう}{勉強}{Benkyou} \jpkana{を}{o} 
    \jpword{すこ}{少}{suko}\jpkana{し}{shi} 
    \jpword{やす}{休}{yasu}\jpkana{んで、}{nde,} 
    \jpword{さんぽ}{散歩}{sanpo} \jpkana{に}{ni} 
    \jpword{い}{行}{i}\jpkana{ったり、}{ttari,}

    \vspace{1.5em}
    \noindent{\Large \color{articolor} \textbf{Arti:} Aku istirahat belajar sebentar, lalu pergi jalan-jalan (santai), atau}
\end{convobox}

\begin{convobox}{24}
    \noindent
    \jpkana{ジムに}{Jimu ni} 
    \jpword{い}{行}{i}\jpkana{ったりしています。}{ttari shite imasu.}

    \vspace{1.5em}
    \noindent{\Large \color{articolor} \textbf{Arti:} pergi ke tempat gym.}
\end{convobox}

% --- LESSON BUFFER 4 ---
\begin{lessonbox}{Catatan Belajar: Kehilangan Semangat?}
    Ada kalanya kita merasa malas atau tidak semangat belajar. Dalam bahasa Jepang, motivasi atau kemauan disebut \textbf{やる気 (Yaruki)}. Kalau motivasinya tidak muncul, kita bilang \textit{Yaruki ga denai}. \\
    
    Kalau sedang begitu, tidak apa-apa beristirahat sebentar! Di percakapan ini, kita pakai tata bahasa \textbf{〜たり、〜たり (〜tari, 〜tari)} lagi lho, seperti yang kita pelajari di Bab 3. Gunanya untuk menyebutkan kegiatan-kegiatan acak yang kita lakukan saat sedang istirahat.
\end{lessonbox}

% ================= PERCAKAPAN 25-37 =================
\begin{convobox}{25}
    \noindent
    \jpword{にほんご}{日本語}{Nihongo} \jpkana{の}{no} 
    \jpword{うた}{歌}{uta} \jpkana{を}{o} 
    \jpkana{よく}{yoku} 
    \jpword{き}{聴}{ki}\jpkana{きますか？}{kimasu ka?}

    \vspace{1.5em}
    \noindent{\Large \color{articolor} \textbf{Arti:} Apakah kamu sering mendengarkan lagu bahasa Jepang?}
\end{convobox}

\begin{convobox}{26}
    \noindent
    \jpkana{おすすめの}{Osusume no} 
    \jpword{うた}{歌}{uta} \jpkana{が}{ga} 
    \jpkana{あったら、}{attara,} 
    \jpword{おし}{教}{oshi}\jpkana{えてください。}{ete kudasai.}

    \vspace{1.5em}
    \noindent{\Large \color{articolor} \textbf{Arti:} Kalau ada lagu rekomendasi, tolong beritahu ya.}
\end{convobox}

\begin{convobox}{27}
    \noindent
    \jpword{にほんご}{日本語}{Nihongo} \jpkana{の}{no} 
    \jpword{うた}{歌}{uta} \jpkana{は}{wa} 
    \jpkana{あまり}{amari} 
    \jpword{き}{聴}{ki}\jpkana{きません。}{kimasen.}

    \vspace{1.5em}
    \noindent{\Large \color{articolor} \textbf{Arti:} Aku tidak terlalu sering mendengarkan lagu bahasa Jepang.}
\end{convobox}

\begin{convobox}{28}
    \noindent
    \jpkana{おすすめを}{Osusume o} 
    \jpword{し}{知}{shi}\jpkana{りたいです。}{ritai desu.}

    \vspace{1.5em}
    \noindent{\Large \color{articolor} \textbf{Arti:} Aku yang ingin tahu rekomendasinya.}
\end{convobox}

\begin{convobox}{29}
    \noindent
    \jpword{にほんご}{日本語}{Nihongo} \jpkana{が}{ga} 
    \jpword{いま}{今}{ima} \jpkana{よりもっと}{yori motto} 
    \jpword{じょうず}{上手}{jouzu} \jpkana{になったら、}{ni nattara,}

    \vspace{1.5em}
    \noindent{\Large \color{articolor} \textbf{Arti:} Kalau bahasa Jepangmu nanti sudah menjadi lebih jago dari sekarang,}
\end{convobox}

\begin{convobox}{30}
    \noindent
    \jpword{なに}{何}{Nani} \jpkana{を}{o} 
    \jpkana{したいですか？}{shitai desu ka?}

    \vspace{1.5em}
    \noindent{\Large \color{articolor} \textbf{Arti:} apa yang ingin kamu lakukan?}
\end{convobox}

\begin{convobox}{31}
    \noindent
    \jpword{にほん}{日本}{Nihon}\jpword{りょこう}{旅行}{ryokou} \jpkana{をして、}{o shite,} 
    \jpword{げんち}{現地}{genchi} \jpkana{の}{no} 
    \jpword{ひと}{人}{hito} \jpkana{と}{to} 
    \jpkana{たくさん}{takusan} 
    \jpword{はな}{話}{hana}\jpkana{したいです。}{shitai desu.}

    \vspace{1.5em}
    \noindent{\Large \color{articolor} \textbf{Arti:} Aku ingin jalan-jalan (wisata) di Jepang, dan banyak berbicara dengan penduduk lokal.}
\end{convobox}

\begin{convobox}{32}
    \noindent
    \jpkana{とてもすてきな}{Totemo suteki na} 
    \jpword{もくひょう}{目標}{mokuhyou} \jpkana{ですね。}{desu ne.}

    \vspace{1.5em}
    \noindent{\Large \color{articolor} \textbf{Arti:} Itu adalah tujuan/target yang sangat indah ya.}
\end{convobox}

\begin{convobox}{33}
    \noindent
    \jpword{きょう}{今日}{Kyou} \jpkana{は}{wa} 
    \jpkana{お}{o}\jpword{はなし}{話}{hanashi} \jpkana{できて}{dekite} 
    \jpword{たの}{楽}{tano}\jpkana{しかったです。}{shikatta desu.}

    \vspace{1.5em}
    \noindent{\Large \color{articolor} \textbf{Arti:} Hari ini aku senang sekali bisa mengobrol.}
\end{convobox}

\begin{convobox}{34}
    \noindent
    \jpword{べんきょう}{勉強}{Benkyou} \jpkana{が}{ga} 
    \jpkana{うまくいかない}{umaku ikanai} 
    \jpword{とき}{時}{toki} \jpkana{や、}{ya,}

    \vspace{1.5em}
    \noindent{\Large \color{articolor} \textbf{Arti:} Saat belajarnya sedang tidak berjalan lancar, atau}
\end{convobox}

\begin{convobox}{35}
    \noindent
    \jpkana{ちょっと}{Chotto} 
    \jpword{きゅう}{休}{kyuu}\jpkana{けいしたい}{kei shitai} 
    \jpword{とき}{時}{toki} \jpkana{は、}{wa,}

    \vspace{1.5em}
    \noindent{\Large \color{articolor} \textbf{Arti:} saat ingin sedikit beristirahat,}
\end{convobox}

\begin{convobox}{36}
    \noindent
    \jpkana{また}{Mata} \jpkana{いつでも}{itsu demo} 
    \jpkana{ここに}{koko ni} 
    \jpword{き}{来}{ki}\jpkana{てくださいね。}{te kudasai ne.}

    \vspace{1.5em}
    \noindent{\Large \color{articolor} \textbf{Arti:} silakan datang ke sini kapan saja ya.}
\end{convobox}

\begin{convobox}{37}
    \noindent
    \jpkana{それでは。}{Sore dewa.}

    \vspace{1.5em}
    \noindent{\Large \color{articolor} \textbf{Arti:} Kalau begitu, (sampai jumpa).}
\end{convobox}

% --- SUMMARY BOX ---
\vspace{2em}
\begin{tcolorbox}[
    colback=blue!5!white,
    colframe=blue!80!black,
    boxrule=3pt,
    arc=5mm,
    title={\huge\bfseries \sffamily 🌟 Rangkuman Bab 6 🌟},
    fonttitle=\color{white},
    attach boxed title to top center={yshift=-4mm},
    boxed title style={colback=blue!80!black, rounded corners},
    left=5mm, right=5mm, top=7mm, bottom=5mm
]
\Large \textbf{Otsukaresama desu! Kamu telah menyelesaikan Bab 6!} Yuk kita cek apa saja yang sudah kita pelajari:
\begin{itemize}
    \item \textbf{Kalimat Pujian:} \textit{Otsukaresama desu} (Terima kasih atas kerja kerasmu).
    \item \textbf{Kata Keterangan Waktu:} \textit{\dots ta ato} (Sesudah) dan \textit{\dots ru mae ni} (Sebelum).
    \item \textbf{Kesukaan / Favorit:} \textit{O-ki ni iri}.
    \item \textbf{Tiruan Bunyi:} \textit{Onomatope}. Contoh: \textit{Choki-choki} (Suara gunting).
    \item \textbf{Motivasi Belajar:} \textit{Yaruki}. Kalau lagi malas, bilang saja \textit{Yaruki ga denai}.
\end{itemize}
Ingat pesan di akhir percakapan ya! Kalau kamu capek atau belajarnya terasa sulit, tidak apa-apa untuk istirahat sejenak. Kamu punya tujuan yang hebat, jadi pelan-pelan saja yang penting pasti!
\end{tcolorbox}

\newpage
\chapter{End Your Day in a Hot Spring}

% ================= PERCAKAPAN 1-2 =================
\begin{convobox}{1}
    \noindent
    \jpkana{こんばんは。}{Konbanwa.}

    \vspace{1.5em}
    \noindent{\Large \color{articolor} \textbf{Arti:} Selamat malam.}
\end{convobox}

\begin{convobox}{2}
    \noindent
    \jpword{きょう}{今日}{Kyou} \jpkana{も}{mo} 
    \jpkana{お}{o}\jpword{つか}{疲}{tsuka}\jpkana{れ}{re}\jpword{さま}{様}{sama} \jpkana{でした。}{deshita.}

    \vspace{1.5em}
    \noindent{\Large \color{articolor} \textbf{Arti:} Hari ini juga, terima kasih atas kerja kerasmu.}
\end{convobox}

\begin{lessonbox}{Catatan Belajar: Selamat Malam!}
    Adik-adik, kalau hari sudah gelap, kita sapa teman dengan \textbf{こんばんは (Konbanwa)} yang artinya "Selamat malam". Terus, perhatikan kata \textit{Otsukaresama deshita}. Di bab sebelumnya pakai \textit{desu}, sekarang pakai \textit{deshita} karena harinya sudah hampir berakhir!
\end{lessonbox}

% ================= PERCAKAPAN 3-4 =================
\begin{convobox}{3}
    \noindent
    \jpkana{ゆっくり}{Yukkuri} 
    \jpword{おんせん}{温泉}{onsen} \jpkana{に}{ni} 
    \jpword{つか}{浸}{tsuka}\jpkana{りながら、}{ri nagara,}

    \vspace{1.5em}
    \noindent{\Large \color{articolor} \textbf{Arti:} Sambil berendam santai di pemandian air panas (onsen),}
\end{convobox}

\begin{convobox}{4}
    \noindent
    \jpword{いちにち}{一日}{Ichi-nichi} \jpkana{を}{o} 
    \jpword{ふ}{振}{fu}\jpkana{り}{ri}\jpword{かえ}{返}{kae}\jpkana{りましょう。}{rimashou.}

    \vspace{1.5em}
    \noindent{\Large \color{articolor} \textbf{Arti:} mari kita ingat-ingat (renungkan) apa saja yang terjadi hari ini.}
\end{convobox}

\begin{lessonbox}{Catatan Belajar: Melakukan 2 Hal Sekaligus!}
    Coba lihat kata \textbf{〜ながら (〜nagara)}. Ini artinya "Sambil". Jadi kalau kamu mandi sambil nyanyi, atau makan sambil nonton TV, kamu bisa pakai \textit{nagara} di tengah kalimatnya. Seru kan?
\end{lessonbox}

% ================= PERCAKAPAN 5-6 =================
\begin{convobox}{5}
    \noindent
    \jpword{きょう}{今日}{Kyou} \jpkana{の}{no} 
    \jpword{てんき}{天気}{tenki} \jpkana{は}{wa} 
    \jpkana{どうでしたか？}{dou deshita ka?}

    \vspace{1.5em}
    \noindent{\Large \color{articolor} \textbf{Arti:} Bagaimana cuaca hari ini?}
\end{convobox}

\begin{convobox}{6}
    \noindent
    \jpword{いちにちじゅう}{一日中}{Ichi-nichi-juu} 
    \jpword{あめ}{雨}{ame} \jpkana{でした。}{deshita.}

    \vspace{1.5em}
    \noindent{\Large \color{articolor} \textbf{Arti:} Sepanjang hari hujan.}
\end{convobox}

\begin{lessonbox}{Catatan Belajar: Sepanjang Hari}
    Kalau kamu melakukan sesuatu dari pagi sampai malam, kamu bisa bilang \textbf{一日中 (Ichi-nichi-juu)}. Di cerita ini hujannya \textit{ichi-nichi-juu} alias nggak berhenti seharian!
\end{lessonbox}

% ================= PERCAKAPAN 7-8 =================
\begin{convobox}{7}
    \noindent
    \jpkana{そうですか。}{Sou desu ka.}

    \vspace{1.5em}
    \noindent{\Large \color{articolor} \textbf{Arti:} Oh, begitu.}
\end{convobox}

\begin{convobox}{8}
    \noindent
    \jpword{きょう}{今日}{Kyou} \jpkana{は}{wa} 
    \jpword{いそが}{忙}{isoga}\jpkana{しかったですか？}{shikatta desu ka?}

    \vspace{1.5em}
    \noindent{\Large \color{articolor} \textbf{Arti:} Apakah hari ini sibuk?}
\end{convobox}

\begin{lessonbox}{Catatan Belajar: Tanya Kabar Masa Lalu}
    Ingat kan kalau kata sifat bisa berubah jadi masa lalu? Kata "Sibuk" adalah \textit{Isogashii}. Karena bertanyanya malam hari untuk kejadian siang tadi, kata \textit{Isogashii} berubah bentuk menjadi \textbf{Isogashikatta}!
\end{lessonbox}

% ================= PERCAKAPAN 9-10 =================
\begin{convobox}{9}
    \noindent
    \jpkana{はい、}{Hai,} 
    \jpword{すこ}{少}{suko}\jpkana{し}{shi} 
    \jpword{いそが}{忙}{isoga}\jpkana{しかったです。}{shikatta desu.}

    \vspace{1.5em}
    \noindent{\Large \color{articolor} \textbf{Arti:} Ya, sedikit sibuk.}
\end{convobox}

\begin{convobox}{10}
    \noindent
    \jpword{きょう}{今日}{Kyou} \jpkana{は、}{wa,} 
    \jpword{なに}{何}{nani} \jpkana{を}{o} 
    \jpkana{しましたか？}{shimashita ka?}

    \vspace{1.5em}
    \noindent{\Large \color{articolor} \textbf{Arti:} Hari ini, (kamu) melakukan apa?}
\end{convobox}

\begin{lessonbox}{Catatan Belajar: "Melakukan Apa?"}
    Ini adalah pertanyaan penting kalau kamu mau kepo sama hari temanmu! \textbf{何をしましたか？ (Nani o shimashita ka?)}. Kamu bisa jawab dengan aktivitas yang kamu lakukan hari itu.
\end{lessonbox}

% ================= PERCAKAPAN 11-12 =================
\begin{convobox}{11}
    \noindent
    \jpkana{えーっと、}{Eetto,} 
    \jpword{せんたく}{洗濯}{sentaku} \jpkana{を}{o} \jpkana{して、}{shite,} 
    \jpword{りょうり}{料理}{ryouri} \jpkana{を}{o} \jpkana{して…}{shite...}

    \vspace{1.5em}
    \noindent{\Large \color{articolor} \textbf{Arti:} Emm, mencuci baju, lalu memasak...}
\end{convobox}

\begin{convobox}{12}
    \noindent
    \jpkana{その}{Sono} \jpword{あと}{後}{ato} 
    \jpkana{カフェ}{kafe} \jpkana{に}{ni} 
    \jpword{い}{行}{i}\jpkana{って}{tte} 
    \jpword{ほん}{本}{hon} \jpkana{を}{o} 
    \jpword{よ}{読}{yo}\jpkana{みました。}{mimashita.}

    \vspace{1.5em}
    \noindent{\Large \color{articolor} \textbf{Arti:} setelah itu pergi ke kafe dan membaca buku.}
\end{convobox}

\begin{lessonbox}{Catatan Belajar: Menyambung Kegiatan}
    Kalau kamu melakukan banyak hal berturut-turut, kamu bisa ubah akhirannya jadi \textbf{〜て (〜te)}. Seperti di atas: \textit{Sentaku o shi\textbf{te}}, \textit{Ryouri o shi\textbf{te}}. Artinya "Melakukan A, lalu melakukan B".
\end{lessonbox}

% ================= PERCAKAPAN 13-14 =================
\begin{convobox}{13}
    \noindent
    \jpword{わたし}{私}{Watashi} \jpkana{は、}{wa,} 
    \jpword{いちにちじゅう}{一日中}{ichi-nichi-juu} 
    \jpword{しごと}{仕事}{shigoto} \jpkana{を}{o} \jpkana{して}{shite}

    \vspace{1.5em}
    \noindent{\Large \color{articolor} \textbf{Arti:} Kalau aku, seharian bekerja,}
\end{convobox}

\begin{convobox}{14}
    \noindent
    \jpkana{ちょっと}{chotto} 
    \jpword{つか}{疲}{tsuka}\jpkana{れたので、}{reta node,} 
    \jpword{おんせん}{温泉}{onsen} \jpkana{に}{ni} 
    \jpword{き}{来}{ki}\jpkana{ました。}{mashita.}

    \vspace{1.5em}
    \noindent{\Large \color{articolor} \textbf{Arti:} karena sedikit lelah, jadi aku datang ke onsen.}
\end{convobox}

\begin{lessonbox}{Catatan Belajar: "Karena..." (Sekali Lagi!)}
    Di Bab 2 kita sudah bahas ini loh. \textbf{〜ので (〜node)} artinya "Karena". Di kalimat ini teman kita bilang \textit{Tsukareta node} yang artinya "Karena aku lelah". Makanya dia pergi berendam deh!
\end{lessonbox}

% ================= PERCAKAPAN 15-16 =================
\begin{convobox}{15}
    \noindent
    \jpword{おんせん}{温泉}{Onsen} \jpkana{に}{ni} 
    \jpword{はい}{入}{hai}\jpkana{ると、}{ru to,} 
    \jpword{つか}{疲}{tsuka}\jpkana{れが}{re ga} 
    \jpkana{よく}{yoku} 
    \jpword{と}{取}{to}\jpkana{れるんですよ。}{reru n desu yo.}

    \vspace{1.5em}
    \noindent{\Large \color{articolor} \textbf{Arti:} Kalau masuk ke onsen, rasa lelahnya bisa hilang lho.}
\end{convobox}

\begin{convobox}{16}
    \noindent
    \jpword{つか}{疲}{Tsuka}\jpkana{れた}{reta} 
    \jpword{とき}{時}{toki} \jpkana{や、}{ya,} 
    \jpkana{リラックス}{rirakkusu} \jpkana{したい}{shitai} 
    \jpword{とき}{時}{toki}\jpkana{、}{,}

    \vspace{1.5em}
    \noindent{\Large \color{articolor} \textbf{Arti:} Saat merasa lelah, atau saat ingin bersantai,}
\end{convobox}

\begin{lessonbox}{Catatan Belajar: Kata "Saat / Waktu"}
    Dalam bahasa Jepang, kalau kita mau bilang "Saat... / Ketika...", kita memakai kata \textbf{時 (Toki)}. Misalnya, \textit{Tsukareta toki} (Saat lelah). Mudah, kan?
\end{lessonbox}

% ================= PERCAKAPAN 17-18 =================
\begin{convobox}{17}
    \noindent
    \jpkana{よく}{Yoku} 
    \jpword{なに}{何}{nani} \jpkana{を}{o} 
    \jpkana{しますか？}{shimasu ka?}

    \vspace{1.5em}
    \noindent{\Large \color{articolor} \textbf{Arti:} biasanya kamu sering melakukan apa?}
\end{convobox}

\begin{convobox}{18}
    \noindent
    \jpkana{ハーブティー}{Haabutii} \jpkana{を}{o} 
    \jpword{の}{飲}{no}\jpkana{んだり、}{ndari,}

    \vspace{1.5em}
    \noindent{\Large \color{articolor} \textbf{Arti:} Meminum teh herbal, atau}
\end{convobox}

\begin{lessonbox}{Catatan Belajar: Grammar Favorit Kita!}
    Yey, kita ketemu lagi sama grammar \textbf{〜たり、〜たり (〜tari, 〜tari)}! Ini untuk menyebutkan hal acak yang kita lakukan. Salah satunya adalah \textit{Nondari} (Minum), berasal dari kata \textit{Nomimasu}.
\end{lessonbox}

% ================= PERCAKAPAN 19-20 =================
\begin{convobox}{19}
    \noindent
    \jpword{さんぽ}{散歩}{Sanpo} \jpkana{に}{ni} 
    \jpword{い}{行}{i}\jpkana{ったりします。}{ttari shimasu.}

    \vspace{1.5em}
    \noindent{\Large \color{articolor} \textbf{Arti:} pergi jalan-jalan santai.}
\end{convobox}

\begin{convobox}{20}
    \noindent
    \jpkana{いいですね。}{Ii desu ne.}

    \vspace{1.5em}
    \noindent{\Large \color{articolor} \textbf{Arti:} Wah bagus ya.}
\end{convobox}

\begin{lessonbox}{Catatan Belajar: "Wah, Bagus Ya!"}
    Kalau ada temanmu bercerita hal yang menyenangkan, kamu bisa memujinya dengan senyuman dan bilang \textbf{いいですね (Ii desu ne!)} yang berarti "Bagus/Hebat ya!".
\end{lessonbox}

% ================= PERCAKAPAN 21-22 =================
\begin{convobox}{21}
    \noindent
    \jpkana{ところで、}{Tokoro de,} 
    \jpword{しゅみ}{趣味}{shumi} \jpkana{は}{wa} 
    \jpword{なに}{何}{nani}\jpkana{か}{ka} 
    \jpkana{ありますか？}{arimasu ka?}

    \vspace{1.5em}
    \noindent{\Large \color{articolor} \textbf{Arti:} Ngomong-ngomong, apakah kamu punya suatu hobi?}
\end{convobox}

\begin{convobox}{22}
    \noindent
    \jpkana{ギター}{Gitaa} \jpkana{を}{o} 
    \jpword{ひ}{弾}{hi}\jpkana{くことです。}{ku koto desu.}

    \vspace{1.5em}
    \noindent{\Large \color{articolor} \textbf{Arti:} (Hobiku) adalah bermain gitar.}
\end{convobox}

\begin{lessonbox}{Catatan Belajar: Mengganti Topik \& Hobi}
    Kalau mau pindah ke topik obrolan baru, gunakan \textbf{ところで (Tokoro de)} yang artinya "Ngomong-ngomong". Untuk menjawab hobi, kita gunakan kata \textbf{こと (koto)} setelah kegiatannya supaya jadi kata benda. Contohnya \textit{Hiku koto} (Bermain alat musik).
\end{lessonbox}

% ================= PERCAKAPAN 23-24 =================
\begin{convobox}{23}
    \noindent
    \jpkana{そうなんですね。}{Sou nan desu ne.}

    \vspace{1.5em}
    \noindent{\Large \color{articolor} \textbf{Arti:} Oh begitu ya.}
\end{convobox}

\begin{convobox}{24}
    \noindent
    \jpkana{いつ}{Itsu} 
    \jpkana{その}{sono} 
    \jpword{しゅみ}{趣味}{shumi} \jpkana{を}{o} 
    \jpword{はじ}{始}{haji}\jpkana{めたんですか？}{meta n desu ka?}

    \vspace{1.5em}
    \noindent{\Large \color{articolor} \textbf{Arti:} Kapan kamu mulai hobi tersebut?}
\end{convobox}

\begin{lessonbox}{Catatan Belajar: "Sou nan desu ne!"}
    Kata ini mirip seperti \textit{Sou desu ka}, tapi mendengarnya terasa lebih hangat dan antusias! \textbf{そうなんですね (Sou nan desu ne)} biasa dipakai saat kita mendengarkan cerita teman dengan penuh perhatian.
\end{lessonbox}

% ================= PERCAKAPAN 25-26 =================
\begin{convobox}{25}
    \noindent
    \jpword{にねんまえ}{2年前}{Ni-nen mae} 
    \jpkana{ぐらいです。}{gurai desu.}

    \vspace{1.5em}
    \noindent{\Large \color{articolor} \textbf{Arti:} Kira-kira 2 tahun yang lalu.}
\end{convobox}

\begin{convobox}{26}
    \noindent
    \jpword{きょう}{今日}{Kyou} \jpkana{は、}{wa,} 
    \jpkana{それを}{sore o} 
    \jpkana{しましたか？}{shimashita ka?}

    \vspace{1.5em}
    \noindent{\Large \color{articolor} \textbf{Arti:} Hari ini, apakah kamu melakukannya?}
\end{convobox}

\begin{lessonbox}{Catatan Belajar: "Sore" artinya "Itu"}
    Daripada mengulang kata "bermain gitar" berulang kali, bahasa Jepang punya kata tunjuk \textbf{それ (Sore)} yang artinya "Itu". Jadi \textit{Sore o shimashita ka?} berarti "Apakah kamu melakukan (hobi) itu?".
\end{lessonbox}

% ================= PERCAKAPAN 27-28 =================
\begin{convobox}{27}
    \noindent
    \jpkana{はい、}{Hai,} 
    \jpword{いちじかん}{1時間}{ichi-jikan} \jpkana{ぐらい}{gurai} 
    \jpkana{ギター}{gitaa} \jpkana{を}{o} 
    \jpword{ひ}{弾}{hi}\jpkana{きました。}{kimashita.}

    \vspace{1.5em}
    \noindent{\Large \color{articolor} \textbf{Arti:} Ya, aku bermain gitar kira-kira 1 jam.}
\end{convobox}

\begin{convobox}{28}
    \noindent
    \jpword{きょう}{今日}{Kyou} \jpkana{は、}{wa,} 
    \jpword{なんじ}{何時}{nan-ji} \jpkana{に}{ni} 
    \jpword{ね}{寝}{ne}\jpkana{るつもりですか？}{ru tsumori desu ka?}

    \vspace{1.5em}
    \noindent{\Large \color{articolor} \textbf{Arti:} Hari ini, kamu berencana mau tidur jam berapa?}
\end{convobox}

\begin{lessonbox}{Catatan Belajar: Rencana / Niat}
    Kita bisa menyatakan niat kita melakukan sesuatu pakai kata \textbf{〜つもり (〜tsumori)}. Kalau kamu ditanya \textit{Neru tsumori desu ka?}, artinya temanmu nanya "Kamu ada niatan/rencana mau tidur jam berapa?".
\end{lessonbox}

% ================= PERCAKAPAN 29-30 =================
\begin{convobox}{29}
    \noindent
    \jpword{じゅうじ}{10時}{Juu-ji} \jpkana{に}{ni} 
    \jpword{ね}{寝}{ne}\jpkana{るつもりです。}{ru tsumori desu.}

    \vspace{1.5em}
    \noindent{\Large \color{articolor} \textbf{Arti:} Rencananya mau tidur jam 10.}
\end{convobox}

\begin{convobox}{30}
    \noindent
    \jpword{わたし}{私}{Watashi} \jpkana{は}{wa} 
    \jpword{ふだん}{普段}{fudan} 
    \jpword{じゅういちじ}{11時}{juu-ichi-ji} \jpkana{に}{ni} 
    \jpword{ね}{寝}{ne}\jpkana{ますが、}{masu ga,}

    \vspace{1.5em}
    \noindent{\Large \color{articolor} \textbf{Arti:} Kalau aku biasanya tidur jam 11, tapi}
\end{convobox}

\begin{lessonbox}{Catatan Belajar: "Biasanya"}
    Adik-adik, kalau mau cerita kebiasaan rutinitasmu, pakai kata \textbf{普段 (Fudan)} yang artinya "Biasanya". Misalnya: \textit{Fudan wa 9-ji ni nemasu} (Biasanya aku tidur jam 9).
\end{lessonbox}

% ================= PERCAKAPAN 31-33 =================
\begin{convobox}{31}
    \noindent
    \jpword{きょう}{今日}{Kyou} \jpkana{は}{wa} 
    \jpword{すこ}{少}{suko}\jpkana{し}{shi} 
    \jpword{はや}{早}{haya}\jpkana{めに}{me ni} 
    \jpword{ね}{寝}{ne}\jpkana{ようと}{you to} 
    \jpword{おも}{思}{omo}\jpkana{います。}{imasu.}

    \vspace{1.5em}
    \noindent{\Large \color{articolor} \textbf{Arti:} hari ini sepertinya aku ingin/pikir mau tidur lebih awal.}
\end{convobox}

\begin{convobox}{32}
    \noindent
    \jpkana{お}{o}\jpword{はなし}{話}{hanashi} \jpkana{できて、}{dekite,} 
    \jpword{たの}{楽}{tano}\jpkana{しかったです。}{shikatta desu.}

    \vspace{1.5em}
    \noindent{\Large \color{articolor} \textbf{Arti:} Bisa mengobrol (denganmu), rasanya menyenangkan.}
\end{convobox}

\begin{convobox}{33}
    \noindent
    \jpkana{それでは、}{Sore dewa,} 
    \jpkana{また}{mata} 
    \jpword{あした}{明日}{ashita}\jpkana{。}{.}

    \vspace{1.5em}
    \noindent{\Large \color{articolor} \textbf{Arti:} Kalau begitu, sampai jumpa besok.}
\end{convobox}

\begin{lessonbox}{Catatan Belajar: Berpisah di Malam Hari}
    Saat mau tidur atau berpamitan di ujung hari, kita bisa pakai \textbf{また明日 (Mata ashita)} yang artinya "Sampai jumpa besok!". Kamu juga bisa tidur lebih awal atau \textit{Hayame ni neru} biar besok pagi bangun dengan segar!
\end{lessonbox}

% --- SUMMARY BOX ---
\vspace{2em}
\begin{tcolorbox}[
    colback=blue!5!white,
    colframe=blue!80!black,
    boxrule=3pt,
    arc=5mm,
    title={\huge\bfseries \sffamily 🌟 Rangkuman Bab 7 🌟},
    fonttitle=\color{white},
    attach boxed title to top center={yshift=-4mm},
    boxed title style={colback=blue!80!black, rounded corners},
    left=5mm, right=5mm, top=7mm, bottom=5mm
]
\Large \textbf{Luar Biasa! Kita belajar banyak sekali malam ini!} Ini dia ringkasannya:
\begin{itemize}
    \item \textbf{Menyambung Kalimat:} Gunakan \textbf{〜て (〜te)} untuk kegiatan berturut-turut, dan \textbf{〜ながら (〜nagara)} untuk kegiatan yang dilakukan "Sambil".
    \item \textbf{Kata Keterangan:} \textbf{Ichi-nichi-juu} (Seharian), \textbf{Fudan} (Biasanya), dan \textbf{Tokoro de} (Ngomong-ngomong).
    \item \textbf{Waktu / Saat:} Gunakan \textbf{〜時 (〜toki)}.
    \item \textbf{Menanyakan Niat/Rencana:} Pakai \textbf{〜つもり (〜tsumori)}.
    \item \textbf{Salam:} \textbf{Konbanwa} (Malam), dan \textbf{Mata ashita} (Sampai besok).
\end{itemize}
Karena kita sudah belajar dengan giat dari tadi, yuk kita \textit{yukkuri} beristirahat dan jangan lupa \textit{hayame ni neyou} (tidur lebih cepat) ya! Mata ashita!
\end{tcolorbox}

\newpage
\chapter{Looking Back on the Year}

% ================= PERCAKAPAN 1-2 =================
\begin{convobox}{1}
    \noindent
    \jpkana{こんにちは。}{Konnichiwa.}

    \vspace{1.5em}
    \noindent{\Large \color{articolor} \textbf{Arti:} Halo.}
\end{convobox}

\begin{convobox}{2}
    \noindent
    \jpword{ことし}{今年}{Kotoshi} \jpkana{も}{mo} 
    \jpkana{あと}{ato} \jpword{すこ}{少}{suko}\jpkana{しで}{shi de} 
    \jpword{お}{終}{o}\jpkana{わりですね。}{wari desu ne.}

    \vspace{1.5em}
    \noindent{\Large \color{articolor} \textbf{Arti:} Tahun ini juga tinggal sedikit lagi selesai ya.}
\end{convobox}

\begin{lessonbox}{Catatan Belajar: Tinggal Sedikit Lagi!}
    Halo adik-adik! Kalau kamu sedang menunggu sesuatu yang hampir selesai, kamu bisa pakai kata ajaib \textbf{あと少しで (Ato sukoshi de)} yang artinya "Tinggal sedikit lagi". Contohnya: \textit{Ato sukoshi de liburan!} Wah, asik banget!
\end{lessonbox}

% ================= PERCAKAPAN 3-4 =================
\begin{convobox}{3}
    \noindent
    \jpkana{ゆっくり}{Yukkuri} 
    \jpword{ちゃ}{お茶}{o-cha} \jpkana{でも}{demo} 
    \jpword{の}{飲}{no}\jpkana{みながら、}{mi nagara,}

    \vspace{1.5em}
    \noindent{\Large \color{articolor} \textbf{Arti:} Sambil bersantai meminum teh atau semacamnya,}
\end{convobox}

\begin{convobox}{4}
    \noindent
    \jpword{いっしょ}{一緒}{Issho} \jpkana{に}{ni} 
    \jpkana{この}{kono} 
    \jpword{いちねん}{1年}{ichi-nen} \jpkana{を}{o} 
    \jpword{ふ}{振}{fu}\jpkana{り}{ri}\jpword{かえ}{返}{kae}\jpkana{りましょう。}{rimashou.}

    \vspace{1.5em}
    \noindent{\Large \color{articolor} \textbf{Arti:} mari kita kenang (ingat kembali) 1 tahun ini bersama-sama.}
\end{convobox}

\begin{lessonbox}{Catatan Belajar: Mari Kita...!}
    Kalau kamu mau mengajak temanmu melakukan sesuatu, ubah akhiran kata kerjanya menjadi \textbf{〜ましょう (〜mashou)}. Artinya adalah "Mari kita\dots". Contoh: \textit{Furikaerimashou!} (Mari kita ingat-ingat!).
\end{lessonbox}

% ================= PERCAKAPAN 5-6 =================
\begin{convobox}{5}
    \noindent
    \jpword{ことし}{今年}{Kotoshi} \jpkana{は、}{wa,} 
    \jpword{たの}{楽}{tano}\jpkana{しい}{shii} 
    \jpword{いちねん}{1年}{ichi-nen} \jpkana{でしたか？}{deshita ka?}

    \vspace{1.5em}
    \noindent{\Large \color{articolor} \textbf{Arti:} Apakah tahun ini adalah tahun yang menyenangkan?}
\end{convobox}

\begin{convobox}{6}
    \noindent
    \jpkana{それとも、}{Sore tomo,} 
    \jpword{たいへん}{大変}{taihen} \jpkana{な}{na} 
    \jpword{いちねん}{1年}{ichi-nen} \jpkana{でしたか？}{deshita ka?}

    \vspace{1.5em}
    \noindent{\Large \color{articolor} \textbf{Arti:} Atau, apakah tahun yang berat?}
\end{convobox}

\begin{lessonbox}{Catatan Belajar: Atau yang Mana?}
    Saat kamu memberikan 2 pilihan pertanyaan, gunakan kata penyambung \textbf{それとも (Soretomo)} yang artinya "Atau". Menyenangkan? \textit{Soretomo} (Atau) berat?
\end{lessonbox}

% ================= PERCAKAPAN 7-8 =================
\begin{convobox}{7}
    \noindent
    \jpkana{どちらかというと、}{Dochira ka to iu to,} 
    \jpword{たいへん}{大変}{taihen} \jpkana{な}{na} 
    \jpword{いちねん}{1年}{ichi-nen} \jpkana{でした。}{deshita.}

    \vspace{1.5em}
    \noindent{\Large \color{articolor} \textbf{Arti:} Kalau boleh dibilang, ini tahun yang berat.}
\end{convobox}

\begin{convobox}{8}
    \noindent
    \jpkana{そうですか。}{Sou desu ka.}

    \vspace{1.5em}
    \noindent{\Large \color{articolor} \textbf{Arti:} Oh begitu.}
\end{convobox}

\begin{lessonbox}{Catatan Belajar: "Kalau Disuruh Milih..."}
    Nah, \textbf{どちらかというと (Dochira ka to iu to)} ini adalah kalimat gaul orang Jepang kalau mereka ragu-ragu tapi harus memilih! Artinya kurang lebih: "Kalau disuruh milih sih\dots" atau "Bisa dibilang sih\dots".
\end{lessonbox}

% ================= PERCAKAPAN 9-10 =================
\begin{convobox}{9}
    \noindent
    \jpword{ことし}{今年}{Kotoshi}\jpkana{、}{,} 
    \jpword{いちばん}{一番}{ichiban} 
    \jpword{たの}{楽}{tano}\jpkana{しかった}{shikatta} 
    \jpkana{こと}{koto} \jpkana{は}{wa} 
    \jpword{なん}{何}{nan} \jpkana{でしたか？}{deshita ka?}

    \vspace{1.5em}
    \noindent{\Large \color{articolor} \textbf{Arti:} Tahun ini, hal apa yang paling menyenangkan?}
\end{convobox}

\begin{convobox}{10}
    \noindent
    \jpword{なつ}{夏}{Natsu} \jpkana{に、}{ni,} 
    \jpword{かぞく}{家族}{kazoku} \jpkana{で}{de} 
    \jpkana{キャンプ}{kyanpu} \jpkana{に}{ni} 
    \jpword{い}{行}{i}\jpkana{った}{tta} 
    \jpkana{こと}{koto} \jpkana{です。}{desu.}

    \vspace{1.5em}
    \noindent{\Large \color{articolor} \textbf{Arti:} (Hal yang menyenangkan adalah) Pergi berkemah bersama keluarga di musim panas.}
\end{convobox}

\begin{lessonbox}{Catatan Belajar: Paling Juara!}
    Adik-adik tahu juara nomor 1? Di Jepang, kata untuk "Paling" atau "Nomor Satu" adalah \textbf{一番 (Ichiban)}! Jadi \textit{Ichiban tanoshikatta} artinya "Yaaang paling menyenangkan!".
\end{lessonbox}

% ================= PERCAKAPAN 11-12 =================
\begin{convobox}{11}
    \noindent
    \jpkana{では、}{Dewa,} 
    \jpword{ことし}{今年}{kotoshi}\jpkana{、}{,} 
    \jpword{いちばん}{一番}{ichiban} 
    \jpword{たいへん}{大変}{taihen} \jpkana{だった}{datta} 
    \jpkana{こと}{koto} \jpkana{は}{wa} 
    \jpword{なん}{何}{nan} \jpkana{でしたか？}{deshita ka?}

    \vspace{1.5em}
    \noindent{\Large \color{articolor} \textbf{Arti:} Kalau begitu, tahun ini, hal apa yang paling berat?}
\end{convobox}

\begin{convobox}{12}
    \noindent
    \jpkana{バイク}{Baiku} \jpkana{で}{de} 
    \jpword{ころ}{転}{koro}\jpkana{んで、}{nde,} 
    \jpword{こっせつ}{骨折}{kossetsu} \jpkana{した}{shita} 
    \jpkana{こと}{koto} \jpkana{です。}{desu.}

    \vspace{1.5em}
    \noindent{\Large \color{articolor} \textbf{Arti:} Jatuh dari (menggunakan) motor dan patah tulang.}
\end{convobox}

\begin{lessonbox}{Catatan Belajar: Naik Apa?}
    Coba lihat partikel \textbf{で (De)} setelah kata motor (\textit{Baiku}). Di bahasa Jepang, kalau mau bilang kita pergi "Naik" atau "Menggunakan" suatu alat, kita pakai partikel ini! \textit{Baiku de} = Dengan motor. Jangan sampai jatuh ya!
\end{lessonbox}

% ================= PERCAKAPAN 13-14 =================
\begin{convobox}{13}
    \noindent
    \jpkana{それは}{Sore wa} 
    \jpword{たいへん}{大変}{taihen} \jpkana{でしたね。}{deshita ne.}

    \vspace{1.5em}
    \noindent{\Large \color{articolor} \textbf{Arti:} Itu pasti sangat berat (menyusahkan) ya.}
\end{convobox}

\begin{convobox}{14}
    \noindent
    \jpword{ことし}{今年}{Kotoshi}\jpkana{、}{,} 
    \jpword{いちばん}{一番}{ichiban} 
    \jpword{か}{買}{ka}\jpkana{って}{tte} 
    \jpkana{よかった}{yokatta} 
    \jpword{もの}{物}{mono} \jpkana{は}{wa} 
    \jpword{なん}{何}{nan} \jpkana{でしたか？}{deshita ka?}

    \vspace{1.5em}
    \noindent{\Large \color{articolor} \textbf{Arti:} Tahun ini, barang apa yang paling kamu syukuri karena telah membelinya?}
\end{convobox}

\begin{lessonbox}{Catatan Belajar: Syukurlah Sudah...!}
    Kalau kamu senang melakukan sesuatu, kamu bisa bilang "Syukurlah sudah...". Caranya gampang! Tambahkan \textbf{〜てよかった (〜te yokatta)} di belakang kata kerja. Contohnya \textit{Katte yokatta} = Syukurlah sudah beli (Tidak menyesal beli).
\end{lessonbox}

% ================= PERCAKAPAN 15-16 =================
\begin{convobox}{15}
    \noindent
    \jpkana{ノイズキャンセリング}{Noizu kyanseringu} \jpkana{が}{ga} 
    \jpkana{ついた}{tsuita} 
    \jpkana{イヤホン}{iyahon} \jpkana{です。}{desu.}

    \vspace{1.5em}
    \noindent{\Large \color{articolor} \textbf{Arti:} Earphone yang dilengkapi noise-canceling.}
\end{convobox}

\begin{convobox}{16}
    \noindent
    \jpword{ことし}{今年}{Kotoshi}\jpkana{、}{,} 
    \jpword{いちばん}{一番}{ichiban} 
    \jpkana{がんばった}{ganbatta} 
    \jpkana{こと}{koto} \jpkana{は}{wa} 
    \jpword{なん}{何}{nan} \jpkana{でしたか？}{deshita ka?}

    \vspace{1.5em}
    \noindent{\Large \color{articolor} \textbf{Arti:} Tahun ini, hal apa yang paling kamu usahakan (perjuangkan)?}
\end{convobox}

\begin{lessonbox}{Catatan Belajar: Berjuang Keras!}
    Kamu sering dengar kata "Ganbatte!" kan? Nah, kalau perjuangannya SUDAH selesai atau berlalu, kita ubah menjadi bentuk lampau yaitu \textbf{がんばった (Ganbatta)}! Kamu juga \textit{ganbatta} lho belajar bahasa Jepang tahun ini!
\end{lessonbox}

% ================= PERCAKAPAN 17-18 =================
\begin{convobox}{17}
    \noindent
    \jpkana{アルバイト}{Arubaito} \jpkana{です。}{desu.}

    \vspace{1.5em}
    \noindent{\Large \color{articolor} \textbf{Arti:} Kerja paruh waktu.}
\end{convobox}

\begin{convobox}{18}
    \noindent
    \jpword{ことし}{今年}{Kotoshi}\jpkana{、}{,} 
    \jpword{いちばん}{一番}{ichiban} \jpkana{の}{no} 
    \jpword{まな}{学}{mana}\jpkana{びは}{bi wa} 
    \jpword{なん}{何}{nan} \jpkana{でしたか？}{deshita ka?}

    \vspace{1.5em}
    \noindent{\Large \color{articolor} \textbf{Arti:} Tahun ini, apa pelajaran (hikmah) paling berharga?}
\end{convobox}

\begin{lessonbox}{Catatan Belajar: Pelajaran Berharga}
    Dari kata \textit{Manabu} (Belajar), kita bisa buat kata benda \textbf{学び (Manabi)} yang berarti "Pelajaran hidup" atau "Hikmah". Apa nih \textit{Manabi} kamu tahun ini?
\end{lessonbox}

% ================= PERCAKAPAN 19-20 =================
\begin{convobox}{19}
    \noindent
    \jpword{はたら}{働}{Hatara}\jpkana{きすぎて、}{ki-sugite,} 
    \jpword{たいちょう}{体調}{taichou} \jpkana{を}{o} 
    \jpword{くず}{崩}{kuzu}\jpkana{してしまった}{shite shimatta} 
    \jpword{とき}{時}{toki}\jpkana{、}{,}

    \vspace{1.5em}
    \noindent{\Large \color{articolor} \textbf{Arti:} Saat tubuhku jatuh sakit karena terlalu banyak bekerja,}
\end{convobox}

\begin{convobox}{20}
    \noindent
    \jpword{けんこう}{健康}{Kenkou} \jpkana{が}{ga} 
    \jpword{いちばん}{一番}{ichiban} 
    \jpword{だいじ}{大事}{daiji} \jpkana{だと}{da to} 
    \jpword{まな}{学}{mana}\jpkana{びました。}{bimashita.}

    \vspace{1.5em}
    \noindent{\Large \color{articolor} \textbf{Arti:} aku belajar bahwa kesehatan adalah yang paling penting.}
\end{convobox}

\begin{lessonbox}{Catatan Belajar: Jangan Terlalu Banyak!}
    Kalau kamu melakukan sesuatu secara berlebihan, tambahkan \textbf{〜すぎる (〜sugiru)} di akhir kata kerja! Makan terlalu banyak = \textit{Tabe-sugiru}. Kalau di sini, \textit{Hataraki-sugite} = Terlalu banyak bekerja. Ingat, \textit{kenkou} (kesehatan) itu penting lho!
\end{lessonbox}

% ================= PERCAKAPAN 21-22 =================
\begin{convobox}{21}
    \noindent
    \jpword{らいねん}{来年}{Rainen}\jpkana{、}{,} 
    \jpword{ことし}{今年}{kotoshi} \jpkana{よりも}{yori mo} 
    \jpkana{がんばりたい}{ganbaritai} 
    \jpkana{こと}{koto} \jpkana{は}{wa} 
    \jpword{なん}{何}{nan} \jpkana{ですか？}{desu ka?}

    \vspace{1.5em}
    \noindent{\Large \color{articolor} \textbf{Arti:} Tahun depan, hal apa yang ingin kamu usahakan lebih daripada tahun ini?}
\end{convobox}

\begin{convobox}{22}
    \noindent
    \jpword{らいねん}{来年}{Rainen} \jpkana{は、}{wa,} 
    \jpword{ことし}{今年}{kotoshi} \jpkana{よりも}{yori mo} 
    \jpkana{たくさん}{takusan} 
    \jpkana{ジム}{jimu} \jpkana{に}{ni} 
    \jpword{かよ}{通}{kayo}\jpkana{いたいです。}{itai desu.}

    \vspace{1.5em}
    \noindent{\Large \color{articolor} \textbf{Arti:} Tahun depan, aku ingin lebih banyak (sering) pergi ke gym dibandingkan tahun ini.}
\end{convobox}

\begin{lessonbox}{Catatan Belajar: "Daripada..."}
    Kamu mau membandingkan 2 hal? Pakai saja grammar \textbf{〜よりも (〜yori mo)}. Artinya "Dibandingkan / Daripada". Contoh: \textit{Kotoshi yori mo} = Daripada tahun ini.
\end{lessonbox}

% ================= PERCAKAPAN 23-24 =================
\begin{convobox}{23}
    \noindent
    \jpkana{では、}{Dewa,} 
    \jpword{らいねん}{来年}{rainen}\jpkana{、}{,} 
    \jpword{あたら}{新}{atara}\jpkana{しく}{shiku} 
    \jpword{ちょうせん}{挑戦}{chousen} \jpkana{してみたい}{shite mitai} 
    \jpkana{こと}{koto} \jpkana{は}{wa}

    \vspace{1.5em}
    \noindent{\Large \color{articolor} \textbf{Arti:} Kalau begitu, tahun depan, hal apa yang ingin kamu coba tantang (lakukan) yang baru}
\end{convobox}

\begin{convobox}{24}
    \noindent
    \jpword{なに}{何}{nani}\jpkana{か}{ka} 
    \jpkana{ありますか？}{arimasu ka?}

    \vspace{1.5em}
    \noindent{\Large \color{articolor} \textbf{Arti:} apakah ada?}
\end{convobox}

\begin{lessonbox}{Catatan Belajar: "Aku Ingin Mencoba!"}
    Kalau kamu penasaran dan ingin mencoba sesuatu yang baru, gunakan rumus \textbf{〜てみたい (〜te mitai)}. Berasal dari kata "Melihat", di sini artinya "Ingin mencoba melakukan". Keren kan?
\end{lessonbox}

% ================= PERCAKAPAN 25-26 =================
\begin{convobox}{25}
    \noindent
    \jpword{からて}{空手}{Karate} \jpkana{を}{o} 
    \jpword{なら}{習}{nara}\jpkana{ってみたいです。}{tte mitai desu.}

    \vspace{1.5em}
    \noindent{\Large \color{articolor} \textbf{Arti:} Aku ingin mencoba belajar Karate.}
\end{convobox}

\begin{convobox}{26}
    \noindent
    \jpword{ことし}{今年}{Kotoshi} 
    \jpword{いちねん}{1年}{ichi-nen}\jpkana{、}{,} 
    \jpword{つら}{辛}{tsura}\jpkana{かった}{katta} 
    \jpkana{こと}{koto} \jpkana{や}{ya}

    \vspace{1.5em}
    \noindent{\Large \color{articolor} \textbf{Arti:} Selama 1 tahun ini, hal-hal yang menyakitkan (sulit) dan}
\end{convobox}

\begin{lessonbox}{Catatan Belajar: Menyebutkan Banyak Hal!}
    Partikel \textbf{や (Ya)} dipakai kalau kamu mau bilang "Dan" atau "Atau" saat menyebutkan sebagian contoh! Misalnya, beli Apel \textit{ya} Jeruk \textit{ya} Pisang.
\end{lessonbox}

% ================= PERCAKAPAN 27-28 =================
\begin{convobox}{27}
    \noindent
    \jpword{たいへん}{大変}{taihen} \jpkana{だった}{datta} 
    \jpkana{こと}{koto} \jpkana{も}{mo} 
    \jpkana{あったと}{atta to} 
    \jpword{おも}{思}{omo}\jpkana{います。}{imasu.}

    \vspace{1.5em}
    \noindent{\Large \color{articolor} \textbf{Arti:} hal-hal yang berat, aku pikir pasti (pernah) ada.}
\end{convobox}

\begin{convobox}{28}
    \noindent
    \jpkana{でも、}{Demo,} 
    \jpkana{もうすぐ}{mou sugu} 
    \jpword{あたら}{新}{atara}\jpkana{しい}{shii} 
    \jpword{とし}{年}{toshi} \jpkana{が}{ga} 
    \jpword{はじ}{始}{haji}\jpkana{まります。}{marimasu.}

    \vspace{1.5em}
    \noindent{\Large \color{articolor} \textbf{Arti:} Tapi, sebentar lagi tahun yang baru akan dimulai.}
\end{convobox}

\begin{lessonbox}{Catatan Belajar: Sebentar Lagi!}
    Wah, kata \textbf{もうすぐ (Mou sugu)} sangat berguna loh! Artinya "Sebentar lagi". \textit{Mou sugu} makan malam! \textit{Mou sugu} tahun baru!
\end{lessonbox}

% ================= PERCAKAPAN 29-31 =================
\begin{convobox}{29}
    \noindent
    \jpword{らいねん}{来年}{Rainen} \jpkana{は、}{wa,} 
    \jpword{ことし}{今年}{kotoshi} \jpkana{よりも}{yori mo} 
    \jpkana{もっと}{motto} 
    \jpword{すてき}{素敵}{suteki} \jpkana{な}{na} 
    \jpkana{こと}{koto} \jpkana{が}{ga}

    \vspace{1.5em}
    \noindent{\Large \color{articolor} \textbf{Arti:} Di tahun depan, semoga ada lebih banyak hal yang indah}
\end{convobox}

\begin{convobox}{30}
    \noindent
    \jpkana{たくさん}{takusan} 
    \jpkana{ありますように。}{arimasu you ni.}

    \vspace{1.5em}
    \noindent{\Large \color{articolor} \textbf{Arti:} (lanjutan) daripada tahun ini.}
\end{convobox}

\begin{convobox}{31}
    \noindent
    \jpkana{それでは、}{Sore dewa,} 
    \jpkana{よいお}{yoi o}\jpword{とし}{年}{toshi} \jpkana{を。}{o.}

    \vspace{1.5em}
    \noindent{\Large \color{articolor} \textbf{Arti:} Kalau begitu, selamat menyambut tahun baru.}
\end{convobox}

\begin{lessonbox}{Catatan Belajar: Doa dan Harapan Tahun Baru}
    Bila kamu mau mendoakan hal baik, akhiri kalimatmu dengan \textbf{〜ように (〜you ni)}! Artinya "Semoga\dots". Dan yang terakhir, jangan lupa ucapkan \textbf{よいお年を (Yoi o-toshi o)} ke orang-orang sebelum pergantian tahun, ya!
\end{lessonbox}

% --- SUMMARY BOX ---
\vspace{2em}
\begin{tcolorbox}[
    colback=blue!5!white,
    colframe=blue!80!black,
    boxrule=3pt,
    arc=5mm,
    title={\huge\bfseries \sffamily 🌟 Rangkuman Bab 8 🌟},
    fonttitle=\color{white},
    attach boxed title to top center={yshift=-4mm},
    boxed title style={colback=blue!80!black, rounded corners},
    left=5mm, right=5mm, top=7mm, bottom=5mm
]
\Large \textbf{Selamat! Kamu telah menyelesaikan pelajaran tahun ini!} Banyak sekali grammar sakti yang kita pelajari di sini:
\begin{itemize}
    \item \textbf{Mengajak:} \textit{〜mashou} (Mari kita\dots)
    \item \textbf{Bersyukur:} \textit{〜te yokatta} (Syukurlah sudah\dots)
    \item \textbf{Berlebihan:} \textit{〜sugiru} (Terlalu banyak\dots)
    \item \textbf{Perbandingan:} \textit{〜yori mo} (Daripada\dots)
    \item \textbf{Keinginan baru:} \textit{〜te mitai} (Ingin mencoba\dots)
\end{itemize}
Terima kasih sudah belajar dengan rajin bersama Andi - Japanese Handbook. Tetap pertahankan \textbf{Yaruki} (semangat) kamu ya. \textit{Yoi o-toshi o!} Semoga tahun depan lebih indah!
\end{tcolorbox}

\newpage
\chapter{The Ultimate Guide to Japanese Convenience Store}

% ================= INTRO =================
\begin{convobox}{1}
    \noindent
    \jpkana{みなさん}{Minasan} \jpkana{こんにちは、}{konnichiwa,} \jpkana{たなか}{Tanaka} \jpkana{です。}{desu.}

    \vspace{1.5em}
    \noindent{\Large \color{articolor} \textbf{Arti:} Halo semuanya, saya Tanaka.}
\end{convobox}

\begin{convobox}{2}
    \noindent
    \jpword{きょう}{今日}{Kyou} \jpkana{は、}{wa,} 
    \jpkana{コンビニ}{konbini} \jpkana{での}{de no} 
    \jpword{にほんご}{日本語}{Nihongo} \jpword{かいわ}{会話}{kaiwa} \jpkana{を}{o} 
    \jpword{いっしょ}{一緒}{issho} \jpkana{に}{ni} 
    \jpword{れんしゅう}{練習}{renshuu} \jpkana{しましょう。}{shimashou.}

    \vspace{1.5em}
    \noindent{\Large \color{articolor} \textbf{Arti:} Hari ini, mari kita berlatih percakapan bahasa Jepang di minimarket bersama-sama.}
\end{convobox}

\begin{convobox}{3}
    \noindent
    \jpkana{みなさん}{Minasan} \jpkana{が}{ga} 
    \jpword{にほん}{日本}{Nihon} \jpkana{に}{ni} \jpword{き}{来}{ki}\jpkana{たら、}{tara,} 
    \jpkana{きっと}{kitto} \jpword{いっかい}{1回}{ikkai} \jpkana{は}{wa} 
    \jpkana{コンビニ}{konbini} \jpkana{に}{ni} 
    \jpword{い}{行}{i}\jpkana{くと}{ku to} 
    \jpword{おも}{思}{omo}\jpkana{います。}{imasu.}

    \vspace{1.5em}
    \noindent{\Large \color{articolor} \textbf{Arti:} Jika kalian datang ke Jepang, aku yakin pasti setidaknya 1 kali akan pergi ke minimarket.}
\end{convobox}

\begin{convobox}{4}
    \noindent
    \jpkana{コンビニ}{Konbini} \jpkana{での}{de no} 
    \jpword{かいわ}{会話}{kaiwa} \jpkana{は、}{wa,} 
    \jpword{いぜん}{以前}{izen} \jpkana{の}{no} 
    \jpword{どうが}{動画}{douga} \jpkana{でも}{demo} 
    \jpword{しょうかい}{紹介}{shoukai} \jpkana{したことがありますが、}{shita koto ga arimasu ga,}

    \vspace{1.5em}
    \noindent{\Large \color{articolor} \textbf{Arti:} Percakapan di minimarket sudah pernah aku perkenalkan di video sebelumnya, tapi}
\end{convobox}

\begin{convobox}{5}
    \noindent
    \jpword{こんかい}{今回}{Konkai} \jpkana{の}{no} 
    \jpword{どうが}{動画}{douga} \jpkana{では、}{de wa,} 
    \jpkana{レジで}{reji de} 
    \jpword{てんいん}{店員}{tenin}\jpkana{さんに}{san ni} 
    \jpkana{よく}{yoku} 
    \jpword{き}{聞}{ki}\jpkana{かれる}{kareru} 
    \jpword{しつもん}{質問}{shitsumon} \jpkana{を}{o} 
    \jpword{ひと}{一}{hito}\jpkana{つずつ}{tsu zutsu} 
    \jpword{いっしょ}{一緒}{issho} \jpkana{に}{ni} 
    \jpword{かくにん}{確認}{kakunin} \jpkana{します。}{shimasu.}

    \vspace{1.5em}
    \noindent{\Large \color{articolor} \textbf{Arti:} Di video kali ini, kita akan mengecek satu per satu pertanyaan yang sering ditanyakan oleh kasir di mesin kasir.}
\end{convobox}

\begin{convobox}{6}
    \noindent
    \jpword{どうが}{動画}{Douga} \jpkana{の}{no} 
    \jpword{こうはん}{後半}{kouhan} \jpkana{では、}{de wa,} 
    \jpkana{ロールプレイ}{rooru purei} \jpkana{で}{de} 
    \jpword{じっせんてき}{実践的}{jissenteki} \jpkana{な}{na} 
    \jpword{かいわ}{会話}{kaiwa} \jpkana{の}{no} 
    \jpword{れんしゅう}{練習}{renshuu} \jpkana{も}{mo} \jpkana{するので}{suru node}

    \vspace{1.5em}
    \noindent{\Large \color{articolor} \textbf{Arti:} Di paruh kedua video, karena kita juga akan berlatih percakapan praktis dengan bermain peran (roleplay),}
\end{convobox}

\begin{convobox}{7}
    \noindent
    \jpkana{ぜひ}{Zehi} 
    \jpword{さいご}{最後}{saigo} \jpkana{まで}{made} 
    \jpword{み}{観}{mi}\jpkana{てください。}{te kudasai.}

    \vspace{1.5em}
    \noindent{\Large \color{articolor} \textbf{Arti:} pastikan untuk menonton sampai akhir ya.}
\end{convobox}

\begin{convobox}{8}
    \noindent
    \jpkana{それでは}{Sore dewa} 
    \jpword{さっそく}{早速}{sassoku} 
    \jpword{い}{行}{i}\jpkana{ってみましょう！}{tte mimashou!}

    \vspace{1.5em}
    \noindent{\Large \color{articolor} \textbf{Arti:} Kalau begitu, ayo langsung kita mulai!}
\end{convobox}


% ================= TITLE: CONVERSATION =================
\vspace{2em}
\begin{tcolorbox}[colback=boxframe, colframe=boxframe, rounded corners, boxrule=0pt, arc=3mm, left=2mm, top=2mm, bottom=2mm]
    \Large \bfseries \color{white} 🏪 1. Conversation (Percakapan Kecepatan Normal)
\end{tcolorbox}
\vspace{1em}

\begin{convobox}{9}
    \noindent
    \jpword{まず}{先ず}{Mazu} \jpword{いちど}{一度}{ichido} 
    \jpkana{コンビニ}{konbini} \jpkana{での}{de no} 
    \jpword{かいわ}{会話}{kaiwa} \jpkana{を}{o} 
    \jpword{しぜん}{自然}{shizen} \jpkana{な}{na} 
    \jpkana{スピードで}{supiido de} 
    \jpword{き}{聞}{ki}\jpkana{いてみましょう。}{ite mimashou.}

    \vspace{1.5em}
    \noindent{\Large \color{articolor} \textbf{Arti:} Pertama-tama, mari kita dengarkan sekali percakapan di minimarket dengan kecepatan natural.}
\end{convobox}

\begin{convobox}{10}
    \noindent
    \jpkana{いらっしゃいませ。}{Irasshaimase.} 
    \jpkana{こちら}{Kochira} 
    \jpword{あたた}{温}{atata}\jpkana{めますか？}{memasu ka?}

    \vspace{1.5em}
    \noindent{\Large \color{articolor} \textbf{Arti:} Selamat datang. Apakah ini mau dipanaskan?}
\end{convobox}

\begin{convobox}{11}
    \noindent
    \jpkana{いえ、}{Ie,} \jpword{だいじょうぶ}{大丈夫}{daijoubu} \jpkana{です。}{desu.}

    \vspace{1.5em}
    \noindent{\Large \color{articolor} \textbf{Arti:} Tidak, tidak usah/tidak apa-apa.}
\end{convobox}

\begin{convobox}{12}
    \noindent
    \jpkana{お}{O}\jpword{はし}{箸}{hashi} \jpkana{は}{wa} 
    \jpword{りよう}{利用}{riyou} \jpkana{ですか？}{desu ka?}

    \vspace{1.5em}
    \noindent{\Large \color{articolor} \textbf{Arti:} Apakah menggunakan (membutuhkan) sumpit?}
\end{convobox}

\begin{convobox}{13}
    \noindent
    \jpkana{はい、}{Hai,} 
    \jpkana{お}{o}\jpword{ねが}{願}{nega}\jpkana{いします。}{ishimasu.}

    \vspace{1.5em}
    \noindent{\Large \color{articolor} \textbf{Arti:} Ya, tolong (berikan).}
\end{convobox}

\begin{convobox}{14}
    \noindent
    \jpword{ふくろ}{袋}{Fukuro} \jpkana{は}{wa} 
    \jpword{りよう}{利用}{riyou} \jpkana{ですか？}{desu ka?}

    \vspace{1.5em}
    \noindent{\Large \color{articolor} \textbf{Arti:} Apakah menggunakan (membutuhkan) kantong plastik?}
\end{convobox}

\begin{convobox}{15}
    \noindent
    \jpkana{はい、}{Hai,} 
    \jpkana{お}{o}\jpword{ねが}{願}{nega}\jpkana{いします。}{ishimasu.}

    \vspace{1.5em}
    \noindent{\Large \color{articolor} \textbf{Arti:} Ya, tolong.}
\end{convobox}

\begin{convobox}{16}
    \noindent
    \jpword{いちまい}{1枚}{Ichi-mai} 
    \jpword{さんえん}{3円}{san-en} \jpkana{ですが、}{desu ga,} 
    \jpkana{よろしいでしょうか？}{yoroshii deshou ka?}

    \vspace{1.5em}
    \noindent{\Large \color{articolor} \textbf{Arti:} 1 lembarnya 3 yen, apakah tidak apa-apa?}
\end{convobox}

\begin{convobox}{17}
    \noindent
    \jpkana{はい。}{Hai.}

    \vspace{1.5em}
    \noindent{\Large \color{articolor} \textbf{Arti:} Ya.}
\end{convobox}

\begin{convobox}{18}
    \noindent
    \jpkana{ポイントカード}{Pointo kaado} \jpkana{は}{wa} 
    \jpkana{お}{o}\jpword{も}{持}{mo}\jpkana{ちですか？}{chi desu ka?}

    \vspace{1.5em}
    \noindent{\Large \color{articolor} \textbf{Arti:} Apakah Anda memiliki/membawa kartu poin?}
\end{convobox}

\begin{convobox}{19}
    \noindent
    \jpkana{いえ。}{Ie.}

    \vspace{1.5em}
    \noindent{\Large \color{articolor} \textbf{Arti:} Tidak.}
\end{convobox}

\begin{convobox}{20}
    \noindent
    \jpkana{お}{o}\jpword{かいけい}{会計}{kaikei} 
    \jpword{ろっぴゃくえん}{600円}{roppyaku-en} 
    \jpkana{でございます。}{de gozaimasu.}

    \vspace{1.5em}
    \noindent{\Large \color{articolor} \textbf{Arti:} Total pembayarannya 600 yen.}
\end{convobox}

\begin{convobox}{21}
    \noindent
    \jpkana{Tanaka Pay}{Tanaka Pay} \jpkana{で}{de} 
    \jpkana{お}{o}\jpword{ねが}{願}{nega}\jpkana{いします。}{ishimasu.}

    \vspace{1.5em}
    \noindent{\Large \color{articolor} \textbf{Arti:} Pakai Tanaka Pay, tolong.}
\end{convobox}

\begin{convobox}{22}
    \noindent
    \jpkana{かしこまりました。}{Kashikomarimashita.}

    \vspace{1.5em}
    \noindent{\Large \color{articolor} \textbf{Arti:} Baik, saya mengerti (Sangat Sopan).}
\end{convobox}

\begin{convobox}{23}
    \noindent
    \jpkana{レシート}{Reshiito} 
    \jpword{りよう}{利用}{riyou} \jpkana{ですか？}{desu ka?}

    \vspace{1.5em}
    \noindent{\Large \color{articolor} \textbf{Arti:} Apakah butuh struk?}
\end{convobox}

\begin{convobox}{24}
    \noindent
    \jpkana{いえ、}{Ie,} \jpword{だいじょうぶ}{大丈夫}{daijoubu} \jpkana{です。}{desu.}

    \vspace{1.5em}
    \noindent{\Large \color{articolor} \textbf{Arti:} Tidak, tidak usah.}
\end{convobox}

\begin{convobox}{25}
    \noindent
    \jpkana{ありがとうございました。}{Arigatou gozaimashita.} 
    \jpkana{また}{Mata} 
    \jpkana{お}{o}\jpword{こ}{越}{ko}\jpkana{しくださいませ。}{shi kudasaimase.}

    \vspace{1.5em}
    \noindent{\Large \color{articolor} \textbf{Arti:} Terima kasih banyak. Silakan datang kembali.}
\end{convobox}

\begin{convobox}{26}
    \noindent
    \jpkana{それでは、}{Sore dewa,} 
    \jpword{じゅんばん}{順番}{junban} \jpkana{に}{ni} 
    \jpword{せつめい}{説明}{setsumei} \jpkana{しますね。}{shimasu ne.}

    \vspace{1.5em}
    \noindent{\Large \color{articolor} \textbf{Arti:} Kalau begitu, akan saya jelaskan secara berurutan ya.}
\end{convobox}

% ================= TITLE: PHRASE 1 =================
\vspace{2em}
\begin{tcolorbox}[colback=boxframe, colframe=boxframe, rounded corners, boxrule=0pt, arc=3mm, left=2mm, top=2mm, bottom=2mm]
    \Large \bfseries \color{white} 🍱 2. Phrase 1 (Memanaskan Makanan)
\end{tcolorbox}
\vspace{1em}

\begin{convobox}{27}
    \noindent
    \jpkana{「こちら}{"Kochira} \jpword{あたた}{温}{atata}\jpkana{めますか？」}{memasu ka?"}

    \vspace{1.5em}
    \noindent{\Large \color{articolor} \textbf{Arti:} "Apakah ini mau dipanaskan?"}
\end{convobox}

\begin{convobox}{28}
    \noindent
    \jpword{べんとう}{弁当}{Bentou} \jpkana{や}{ya} 
    \jpkana{パスタ、}{pasuta,} 
    \jpkana{ラーメンなど、}{raamen nado,} 
    \jpword{あたた}{温}{atata}\jpkana{めることができる}{meru koto ga dekiru} 
    \jpword{しょうひん}{商品}{shouhin} \jpkana{は、}{wa,}

    \vspace{1.5em}
    \noindent{\Large \color{articolor} \textbf{Arti:} Produk yang bisa dipanaskan seperti kotak bekal (Bento), pasta, atau ramen,}
\end{convobox}

\begin{convobox}{29}
    \noindent
    \jpword{てんいん}{店員}{Tenin}\jpkana{さんが}{san ga} 
    \jpword{でんし}{電子}{denshi}\jpkana{レンジで}{renji de} 
    \jpword{あたた}{温}{atata}\jpkana{めてくれます。}{mete kuremasu.}

    \vspace{1.5em}
    \noindent{\Large \color{articolor} \textbf{Arti:} akan dipanaskan oleh kasir di dalam microwave.}
\end{convobox}

\begin{convobox}{30}
    \noindent
    \jpword{こた}{答}{Kota}\jpkana{えるときは、}{eru toki wa,} 
    \jpkana{このように}{kono you ni} 
    \jpword{い}{言}{i}\jpkana{います。}{imasu.}

    \vspace{1.5em}
    \noindent{\Large \color{articolor} \textbf{Arti:} Saat menjawab, ucapkanlah seperti ini.}
\end{convobox}

\begin{convobox}{31}
    \noindent
    \jpkana{「はい、}{"Hai,} \jpkana{お}{o}\jpword{ねが}{願}{nega}\jpkana{いします。」}{ishimasu."}

    \vspace{1.5em}
    \noindent{\Large \color{articolor} \textbf{Arti:} "Ya, tolong."}
\end{convobox}

\begin{convobox}{32}
    \noindent
    \jpkana{「いえ、}{"Ie,} \jpword{だいじょうぶ}{大丈夫}{daijoubu} \jpkana{です。」}{desu."}

    \vspace{1.5em}
    \noindent{\Large \color{articolor} \textbf{Arti:} "Tidak, tidak usah."}
\end{convobox}

\begin{convobox}{33}
    \noindent
    \jpword{てんいん}{店員}{Tenin}\jpkana{さんに}{san ni} 
    \jpword{なに}{何}{nani}\jpkana{か}{ka} 
    \jpword{き}{聞}{ki}\jpkana{かれた}{kareta} 
    \jpword{とき}{時}{toki}\jpkana{は、}{wa,}

    \vspace{1.5em}
    \noindent{\Large \color{articolor} \textbf{Arti:} Saat ditanya sesuatu oleh kasir,}
\end{convobox}

\begin{convobox}{34}
    \noindent
    \jpword{きほんてき}{基本的}{Kihonteki} \jpkana{に}{ni} 
    \jpkana{この}{kono} \jpword{ふた}{二}{futa}\jpkana{つを}{tsu o} 
    \jpword{つか}{使}{tsuka}\jpkana{って}{tte} 
    \jpword{こた}{答}{kota}\jpkana{えれば}{ereba} 
    \jpword{だいじょうぶ}{大丈夫}{daijoubu} \jpkana{です。}{desu.}

    \vspace{1.5em}
    \noindent{\Large \color{articolor} \textbf{Arti:} pada dasarnya kalau kamu menjawab menggunakan 2 frasa ini saja sudah cukup (aman).}
\end{convobox}

\begin{convobox}{35}
    \noindent
    \jpword{しつもん}{質問}{Shitsumon} \jpkana{を}{o} 
    \jpword{き}{聞}{ki}\jpkana{き}{ki}\jpword{と}{取}{to}\jpkana{れなかった}{renakatta} 
    \jpword{とき}{時}{toki}\jpkana{は、}{wa,} 
    \jpkana{このように}{kono you ni} 
    \jpword{い}{言}{i}\jpkana{いましょう。}{imashou.}

    \vspace{1.5em}
    \noindent{\Large \color{articolor} \textbf{Arti:} Saat kamu tidak bisa mendengar pertanyaannya dengan jelas, ucapkanlah seperti ini.}
\end{convobox}

\begin{convobox}{36}
    \noindent
    \jpkana{「すみません、}{"Sumimasen,} 
    \jpword{もういちど}{もう一度}{mou ichido} 
    \jpkana{お}{o}\jpword{ねが}{願}{nega}\jpkana{いします。」}{ishimasu."}

    \vspace{1.5em}
    \noindent{\Large \color{articolor} \textbf{Arti:} "Maaf, tolong ulangi sekali lagi."}
\end{convobox}

\begin{lessonbox}{Aksi Pahlawan Konbini: Jurus 2 Jawaban Sakti!}
    \textbf{Penjelasan:} Di Jepang, kasir bicara sangaaat cepat dan sopan (Keigo). Jangan panik! Kamu TIDAK HARUS mengerti setiap kata.
    \\
    \textbf{🎬 ACTION / PRAKTEK:} 
    \begin{enumerate}
        \item Jika kasir menunjuk kotak bento/makananmu, mereka pasti bertanya soal \textit{Microwave} (Atatamemasu ka?).
        \item Cukup anggukkan kepala sedikit dan ucapkan keras-keras: \textbf{"Hai, onegaishimasu!"} (Ya tolong). 
        \item Jika kamu membeli minuman dingin atau tidak mau dipanaskan, gelengkan kepala dan ucapkan: \textbf{"Ie, daijoubu desu!"} (Tidak usah).
    \end{enumerate}
    Dua jurus ini (\textit{Hai, onegaishimasu} dan \textit{Ie, daijoubu desu}) adalah tameng utamamu di Konbini!
\end{lessonbox}

% ================= TITLE: PHRASE 2 =================
\vspace{2em}
\begin{tcolorbox}[colback=boxframe, colframe=boxframe, rounded corners, boxrule=0pt, arc=3mm, left=2mm, top=2mm, bottom=2mm]
    \Large \bfseries \color{white} 🥢 3. Phrase 2 (Sumpit \& Alat Makan)
\end{tcolorbox}
\vspace{1em}

\begin{convobox}{37}
    \noindent
    \jpkana{「お}{"O}\jpword{はし}{箸}{hashi} \jpkana{は}{wa} 
    \jpword{りよう}{利用}{riyou} \jpkana{ですか？」}{desu ka?"}

    \vspace{1.5em}
    \noindent{\Large \color{articolor} \textbf{Arti:} "Apakah membutuhkan sumpit?"}
\end{convobox}

\begin{convobox}{38}
    \noindent
    \jpword{おお}{多}{Oo}\jpkana{くの}{ku no} 
    \jpkana{コンビニ}{konbini} \jpkana{では、}{de wa,} 
    \jpkana{お}{o}\jpword{はし}{箸}{hashi}\jpkana{、}{,} 
    \jpkana{フォーク、}{fooku,} 
    \jpkana{スプーンなどを}{supuun nado o} 
    \jpword{むりょう}{無料}{muryou} \jpkana{でつけてくれます。}{de tsukete kuremasu.}

    \vspace{1.5em}
    \noindent{\Large \color{articolor} \textbf{Arti:} Di banyak minimarket, mereka akan menyertakan sumpit, garpu, sendok, dll secara gratis.}
\end{convobox}

\begin{convobox}{39}
    \noindent
    \jpword{こた}{答}{Kota}\jpkana{えるときは、}{eru toki wa,} 
    \jpword{さき}{先}{saki}\jpkana{ほどと}{hodo to} 
    \jpword{おな}{同}{ona}\jpkana{じように}{ji you ni} 
    \jpword{い}{言}{i}\jpkana{えば}{eba} 
    \jpword{だいじょうぶ}{大丈夫}{daijoubu} \jpkana{です。}{desu.}

    \vspace{1.5em}
    \noindent{\Large \color{articolor} \textbf{Arti:} Saat menjawab, jika kamu mengatakannya sama seperti sebelumnya, itu sudah cukup.}
\end{convobox}

\begin{convobox}{40}
    \noindent
    \jpkana{「はい、}{"Hai,} \jpkana{お}{o}\jpword{ねが}{願}{nega}\jpkana{いします。」}{ishimasu."}

    \vspace{1.5em}
    \noindent{\Large \color{articolor} \textbf{Arti:} "Ya, tolong."}
\end{convobox}

\begin{convobox}{41}
    \noindent
    \jpkana{「いえ、}{"Ie,} \jpword{だいじょうぶ}{大丈夫}{daijoubu} \jpkana{です。」}{desu."}

    \vspace{1.5em}
    \noindent{\Large \color{articolor} \textbf{Arti:} "Tidak, tidak usah."}
\end{convobox}

\begin{convobox}{42}
    \noindent
    \jpkana{もし}{Moshi} \jpword{てんいん}{店員}{tenin}\jpkana{さんが}{san ga} 
    \jpkana{お}{o}\jpword{はし}{箸}{hashi}\jpkana{を}{o} 
    \jpkana{つけるのを}{tsukeru no o} 
    \jpword{わす}{忘}{wasu}\jpkana{れている}{rete iru} 
    \jpword{とき}{時}{toki}\jpkana{は、}{wa,} 
    \jpkana{このように}{kono you ni} 
    \jpword{い}{言}{i}\jpkana{いましょう。}{imashou.}

    \vspace{1.5em}
    \noindent{\Large \color{articolor} \textbf{Arti:} Jika kasir lupa memberikan sumpitnya, ucapkanlah seperti ini.}
\end{convobox}

\begin{convobox}{43}
    \noindent
    \jpkana{「お}{"O}\jpword{はし}{箸}{hashi} \jpkana{お}{o}\jpword{ねが}{願}{nega}\jpkana{いします。」}{ishimasu."}

    \vspace{1.5em}
    \noindent{\Large \color{articolor} \textbf{Arti:} "Tolong berikan sumpitnya."}
\end{convobox}

\begin{lessonbox}{Aksi Pahlawan Konbini: Detektif Alat Makan!}
    \textbf{Penjelasan:} Biasanya kasir akan menebak alat makan apa yang kamu butuhkan. Beli mie? Dapat \textbf{O-hashi} (Sumpit). Beli puding? Dapat \textbf{Supuun} (Sendok kecil).
    \\
    \textbf{🎬 ACTION / PRAKTEK:}
    Jika kamu membeli mie instan untuk dibawa pulang ke hotel, jangan lupa minta sumpitnya JIKA kasir tidak bertanya! 
    Angkat tangan sedikit, tatap kasir sambil tersenyum, dan ucapkan dengan jelas: \textbf{"O-hashi, onegaishimasu!"}.
\end{lessonbox}

% ================= TITLE: PHRASE 3 =================
\vspace{2em}
\begin{tcolorbox}[colback=boxframe, colframe=boxframe, rounded corners, boxrule=0pt, arc=3mm, left=2mm, top=2mm, bottom=2mm]
    \Large \bfseries \color{white} 🛍️ 4. Phrase 3 (Kantong Plastik)
\end{tcolorbox}
\vspace{1em}

\begin{convobox}{44}
    \noindent
    \jpkana{「}{"} \jpword{ふくろ}{袋}{Fukuro} \jpkana{は}{wa} 
    \jpword{りよう}{利用}{riyou} \jpkana{ですか？」}{desu ka?"}

    \vspace{1.5em}
    \noindent{\Large \color{articolor} \textbf{Arti:} "Apakah membutuhkan kantong plastik?"}
\end{convobox}

\begin{convobox}{45}
    \noindent
    \jpkana{「}{"}\jpword{ふくろ}{袋}{Fukuro}\jpkana{」は、}{" wa,} 
    \jpkana{「レジ}{"Reji }\jpword{ぶくろ}{袋}{bukuro}\jpkana{」と}{" to} 
    \jpword{い}{言}{i}\jpkana{うこともあります。}{u koto mo arimasu.}

    \vspace{1.5em}
    \noindent{\Large \color{articolor} \textbf{Arti:} Kata "Fukuro" kadang-kadang juga disebut "Reji-bukuro" (kantong belanja).}
\end{convobox}

\begin{convobox}{46}
    \noindent
    \jpword{ふくろ}{袋}{Fukuro} \jpkana{は}{wa} 
    \jpword{すうねんまえ}{数年前}{suunen-mae} \jpkana{まで}{made} 
    \jpword{むりょう}{無料}{muryou} \jpkana{でしたが、}{deshita ga,}

    \vspace{1.5em}
    \noindent{\Large \color{articolor} \textbf{Arti:} Sampai beberapa tahun lalu plastik itu gratis, tapi}
\end{convobox}

\begin{convobox}{47}
    \noindent
    \jpword{さいきん}{最近}{Saikin} \jpkana{は}{wa} 
    \jpword{ほとん}{殆}{hoton}\jpkana{どの}{do no} 
    \jpkana{コンビニ}{konbini} \jpkana{で}{de} 
    \jpword{ゆうりょう}{有料}{yuuryou} \jpkana{になりました。}{ni narimashita.}

    \vspace{1.5em}
    \noindent{\Large \color{articolor} \textbf{Arti:} Akhir-akhir ini, di hampir semua minimarket menjadi berbayar.}
\end{convobox}

\begin{convobox}{48}
    \noindent
    \jpword{ねだん}{値段}{Nedan} \jpkana{は}{wa} 
    \jpword{さんえん}{3円}{san-en} \jpkana{から}{kara} 
    \jpword{ごえん}{5円}{go-en} \jpkana{ぐらいです。}{gurai desu.}

    \vspace{1.5em}
    \noindent{\Large \color{articolor} \textbf{Arti:} Harganya kira-kira dari 3 yen sampai 5 yen.}
\end{convobox}

\begin{convobox}{49}
    \noindent
    \jpkana{この}{Kono} \jpword{しつもん}{質問}{shitsumon} \jpkana{にも、}{ni mo,} 
    \jpword{さき}{先}{saki}\jpkana{ほどと}{hodo to} 
    \jpword{おな}{同}{ona}\jpkana{じように}{ji you ni} 
    \jpword{こた}{答}{kota}\jpkana{えれば}{ereba} 
    \jpword{だいじょうぶ}{大丈夫}{daijoubu} \jpkana{です。}{desu.}

    \vspace{1.5em}
    \noindent{\Large \color{articolor} \textbf{Arti:} Pada pertanyaan ini juga, jika menjawab dengan cara yang sama seperti tadi sudah cukup.}
\end{convobox}

\begin{convobox}{50}
    \noindent
    \jpkana{「はい、}{"Hai,} \jpkana{お}{o}\jpword{ねが}{願}{nega}\jpkana{いします。」}{ishimasu."}

    \vspace{1.5em}
    \noindent{\Large \color{articolor} \textbf{Arti:} "Ya, tolong."}
\end{convobox}

\begin{convobox}{51}
    \noindent
    \jpkana{「いえ、}{"Ie,} \jpword{だいじょうぶ}{大丈夫}{daijoubu} \jpkana{です。」}{desu."}

    \vspace{1.5em}
    \noindent{\Large \color{articolor} \textbf{Arti:} "Tidak, tidak usah."}
\end{convobox}

\begin{lessonbox}{Aksi Pahlawan Konbini: Siapkan Tas Belanjamu!}
    \textbf{Penjelasan:} Kata kunci yang harus kamu dengar dari kasir adalah \textbf{FUKURO} (Kantong plastik) atau \textbf{REJI-BUKURO}.
    \\
    \textbf{🎬 ACTION / PRAKTEK:} 
    Jadilah pembeli yang ramah lingkungan! Saat kasir menyebut kata \textit{Fukuro}, angkat tas kain yang kamu bawa dari rumah, senyum, dan katakan: \textbf{"Ie, daijoubu desu!"}. Kamu hemat 3 yen dan menyelamatkan bumi! 
\end{lessonbox}

% ================= TITLE: PHRASE 4 =================
\vspace{2em}
\begin{tcolorbox}[colback=boxframe, colframe=boxframe, rounded corners, boxrule=0pt, arc=3mm, left=2mm, top=2mm, bottom=2mm]
    \Large \bfseries \color{white} 💳 5. Phrase 4 (Kartu Poin)
\end{tcolorbox}
\vspace{1em}

\begin{convobox}{52}
    \noindent
    \jpkana{「ポイントカード}{"Pointo kaado} \jpkana{は}{wa} 
    \jpkana{お}{o}\jpword{も}{持}{mo}\jpkana{ちですか？」}{chi desu ka?"}

    \vspace{1.5em}
    \noindent{\Large \color{articolor} \textbf{Arti:} "Apakah Anda membawa kartu poin?"}
\end{convobox}

\begin{convobox}{53}
    \noindent
    \jpkana{コンビニ}{Konbini} \jpkana{には、}{ni wa,} 
    \jpkana{それぞれ}{sorezore} 
    \jpkana{ポイントカード}{pointo kaado} \jpkana{があります。}{ga arimasu.}

    \vspace{1.5em}
    \noindent{\Large \color{articolor} \textbf{Arti:} Di minimarket, masing-masing memiliki kartu poinnya sendiri.}
\end{convobox}

\begin{convobox}{54}
    \noindent
    \jpkana{お}{O}\jpword{かいけい}{会計}{kaikei} \jpkana{の}{no} 
    \jpword{まえ}{前}{mae} \jpkana{に、}{ni,} 
    \jpkana{このように}{kono you ni} 
    \jpword{き}{聞}{ki}\jpkana{かれることが}{kareru koto ga} 
    \jpword{おお}{多}{oo}\jpkana{いです。}{i desu.}

    \vspace{1.5em}
    \noindent{\Large \color{articolor} \textbf{Arti:} Sebelum melakukan pembayaran, kamu sering ditanya seperti ini.}
\end{convobox}

\begin{convobox}{55}
    \noindent
    \jpkana{ポイントカード}{Pointo kaado} \jpkana{を}{o} 
    \jpword{も}{持}{mo}\jpkana{っている}{tte iru} 
    \jpword{とき}{時}{toki}\jpkana{は、}{wa,} 
    \jpkana{「はい」と}{"Hai" to} \jpword{い}{言}{i}\jpkana{って}{tte}

    \vspace{1.5em}
    \noindent{\Large \color{articolor} \textbf{Arti:} Saat kamu membawa/memiliki kartu poin, katakan "Ya", lalu}
\end{convobox}

\begin{convobox}{56}
    \noindent
    \jpkana{カード}{Kaado} \jpkana{や}{ya} \jpkana{アプリの}{apuri no} 
    \jpkana{バーコードを}{baakoodo o} 
    \jpword{てんいん}{店員}{tenin}\jpkana{さんに}{san ni} 
    \jpword{み}{見}{mi}\jpkana{せます。}{semasu.}

    \vspace{1.5em}
    \noindent{\Large \color{articolor} \textbf{Arti:} tunjukkan kartu atau barcode aplikasi kepada kasir.}
\end{convobox}

\begin{convobox}{57}
    \noindent
    \jpword{も}{持}{Mo}\jpkana{っていない}{tte inai} 
    \jpword{とき}{時}{toki}\jpkana{は、}{wa,} 
    \jpkana{「いえ。」と}{"Ie." to} 
    \jpword{い}{言}{i}\jpkana{えば}{eba} 
    \jpword{だいじょうぶ}{大丈夫}{daijoubu} \jpkana{です。}{desu.}

    \vspace{1.5em}
    \noindent{\Large \color{articolor} \textbf{Arti:} Saat tidak punya, bilang "Tidak" saja sudah cukup.}
\end{convobox}

\begin{lessonbox}{Aksi Pahlawan Konbini: Turis Tidak Butuh Poin!}
    \textbf{Penjelasan:} Kartu poin biasanya dikumpulkan oleh orang lokal Jepang yang sering belanja ke sana (seperti T-Point atau d-Point).
    \\
    \textbf{🎬 ACTION / PRAKTEK:} 
    Karena kamu sedang jalan-jalan (turis), kamu tidak butuh ini! Jika kasir bilang kata \textbf{POINTO KAADO}, segera gelengkan kepalamu dan bilang dengan tegas tapi sopan: \textbf{"Ie!"} atau \textbf{"Nai desu"} (Tidak ada).
\end{lessonbox}

% ================= TITLE: PHRASE 5 =================
\vspace{2em}
\begin{tcolorbox}[colback=boxframe, colframe=boxframe, rounded corners, boxrule=0pt, arc=3mm, left=2mm, top=2mm, bottom=2mm]
    \Large \bfseries \color{white} 📱 6. Phrase 5 (Metode Pembayaran)
\end{tcolorbox}
\vspace{1em}

\begin{convobox}{58}
    \noindent
    \jpkana{「Tanaka Pay}{"Tanaka Pay} \jpkana{で}{de} 
    \jpkana{お}{o}\jpword{ねが}{願}{nega}\jpkana{いします。」}{ishimasu."}

    \vspace{1.5em}
    \noindent{\Large \color{articolor} \textbf{Arti:} "Pakai Tanaka Pay, tolong."}
\end{convobox}

\begin{convobox}{59}
    \noindent
    \jpkana{お}{O}\jpword{かいけい}{会計}{kaikei} \jpkana{の}{no} 
    \jpword{ほうほう}{方法}{houhou} \jpkana{を}{o} 
    \jpword{つた}{伝}{tsuta}\jpkana{える}{eru} 
    \jpkana{フレーズです。}{fureezu desu.}

    \vspace{1.5em}
    \noindent{\Large \color{articolor} \textbf{Arti:} Ini adalah frasa untuk menyampaikan metode pembayaran.}
\end{convobox}

\begin{convobox}{60}
    \noindent
    \jpkana{コンビニ}{Konbini} \jpkana{では、}{de wa,} 
    \jpword{げんきん}{現金}{genkin}\jpkana{、}{,} 
    \jpkana{クレジットカード、}{kurejitto kaado,} 
    \jpkana{スマホ}{sumaho}\jpword{けっさい}{決済}{kessai} \jpkana{など、}{nado,}

    \vspace{1.5em}
    \noindent{\Large \color{articolor} \textbf{Arti:} Di minimarket, ada uang tunai (Genkin), kartu kredit, pembayaran smartphone, dll,}
\end{convobox}

\begin{convobox}{61}
    \noindent
    \jpkana{いろいろな}{Iroiro na} 
    \jpword{ほうほう}{方法}{houhou} \jpkana{で}{de} 
    \jpword{しはら}{支払}{shihara}\jpkana{いができます。}{i ga dekimasu.}

    \vspace{1.5em}
    \noindent{\Large \color{articolor} \textbf{Arti:} kamu bisa membayar dengan berbagai macam metode.}
\end{convobox}

\begin{convobox}{62}
    \noindent
    \jpword{たと}{例}{Tato}\jpkana{えば}{eba} 
    \jpkana{クレジットカードで}{kurejitto kaado de} 
    \jpword{はら}{払}{hara}\jpkana{いたい}{itai} 
    \jpword{とき}{時}{toki}\jpkana{は、}{wa,} 
    \jpkana{このように}{kono you ni} 
    \jpword{い}{言}{i}\jpkana{います。}{imasu.}

    \vspace{1.5em}
    \noindent{\Large \color{articolor} \textbf{Arti:} Misalnya, saat ingin membayar dengan kartu kredit, ucapkanlah seperti ini.}
\end{convobox}

\begin{convobox}{63}
    \noindent
    \jpkana{「クレジットカードで}{"Kurejitto kaado de} 
    \jpkana{お}{o}\jpword{ねが}{願}{nega}\jpkana{いします。」}{ishimasu."}

    \vspace{1.5em}
    \noindent{\Large \color{articolor} \textbf{Arti:} "Pakai Kartu Kredit, tolong."}
\end{convobox}

\begin{lessonbox}{Aksi Pahlawan Konbini: Sebut Senjatamu!}
    \textbf{Penjelasan:} Di Jepang, mesin kasir harus ditekan sesuai alat bayarmu. Jadi kamu HARUS kasih tahu kasir mau bayar pakai apa!
    \\
    \textbf{🎬 ACTION / PRAKTEK:} 
    Pegang alat bayarmu (Uang cash, Kartu, atau HP). Tunjukkan pada kasir, dan gunakan rumus sakti ini: \textbf{[Alat Bayar] + de onegaishimasu!}
    \begin{itemize}
        \item Uang Tunai: \textbf{"Genkin de onegaishimasu!"}
        \item Kartu IC (Suica/Pasmo): \textbf{"Suica de onegaishimasu!"}
        \item Kartu Kredit: \textbf{"Kaado de onegaishimasu!"}
    \end{itemize}
\end{lessonbox}

% ================= TITLE: PHRASE 6 =================
\vspace{2em}
\begin{tcolorbox}[colback=boxframe, colframe=boxframe, rounded corners, boxrule=0pt, arc=3mm, left=2mm, top=2mm, bottom=2mm]
    \Large \bfseries \color{white} 🧾 7. Phrase 6 (Struk Belanja)
\end{tcolorbox}
\vspace{1em}

\begin{convobox}{64}
    \noindent
    \jpkana{「レシート}{"Reshiito} \jpword{りよう}{利用}{riyou} \jpkana{ですか？」}{desu ka?"}

    \vspace{1.5em}
    \noindent{\Large \color{articolor} \textbf{Arti:} "Apakah butuh struk?"}
\end{convobox}

\begin{convobox}{65}
    \noindent
    \jpkana{お}{O}\jpword{かいけい}{会計}{kaikei} \jpkana{の}{no} 
    \jpword{あと}{後}{ato} \jpkana{に、}{ni,} 
    \jpkana{よく}{yoku} 
    \jpword{き}{聞}{ki}\jpkana{かれる}{kareru} 
    \jpkana{フレーズです。}{fureezu desu.}

    \vspace{1.5em}
    \noindent{\Large \color{articolor} \textbf{Arti:} Ini adalah frasa yang sering ditanyakan setelah pembayaran selesai.}
\end{convobox}

\begin{convobox}{66}
    \noindent
    \jpkana{これも}{Kore mo} 
    \jpkana{さっきと}{sakki to} 
    \jpword{おな}{同}{ona}\jpkana{じように}{ji you ni} 
    \jpword{こた}{答}{kota}\jpkana{えれば}{ereba} 
    \jpword{だいじょうぶ}{大丈夫}{daijoubu} \jpkana{です。}{desu.}

    \vspace{1.5em}
    \noindent{\Large \color{articolor} \textbf{Arti:} Ini juga jika dijawab sama seperti tadi sudah cukup kok.}
\end{convobox}

\begin{convobox}{67}
    \noindent
    \jpkana{「はい、}{"Hai,} \jpkana{お}{o}\jpword{ねが}{願}{nega}\jpkana{いします。」}{ishimasu."}

    \vspace{1.5em}
    \noindent{\Large \color{articolor} \textbf{Arti:} "Ya, tolong."}
\end{convobox}

\begin{convobox}{68}
    \noindent
    \jpkana{「いえ、}{"Ie,} \jpword{だいじょうぶ}{大丈夫}{daijoubu} \jpkana{です。」}{desu."}

    \vspace{1.5em}
    \noindent{\Large \color{articolor} \textbf{Arti:} "Tidak, tidak usah."}
\end{convobox}

% ================= TITLE: TIPS =================
\vspace{2em}
\begin{tcolorbox}[colback=boxframe, colframe=boxframe, rounded corners, boxrule=0pt, arc=3mm, left=2mm, top=2mm, bottom=2mm]
    \Large \bfseries \color{white} 💡 8. Tips Rahasia Konbini (Toilet)
\end{tcolorbox}
\vspace{1em}

\begin{convobox}{69}
    \noindent
    \jpword{おお}{多}{Oo}\jpkana{くの}{ku no} 
    \jpkana{コンビニ}{konbini} \jpkana{には、}{ni wa,} 
    \jpkana{トイレが}{toire ga} 
    \jpkana{あります。}{arimasu.}

    \vspace{1.5em}
    \noindent{\Large \color{articolor} \textbf{Arti:} Di banyak minimarket, terdapat toilet.}
\end{convobox}

\begin{convobox}{70}
    \noindent
    \jpkana{トイレを}{Toire o} 
    \jpword{か}{借}{ka}\jpkana{りるときは、}{riru toki wa,} 
    \jpword{てんいん}{店員}{tenin}\jpkana{さんに}{san ni} 
    \jpword{ひとこえ}{一声}{hito-koe} 
    \jpkana{かけましょう。}{kakemashou.}

    \vspace{1.5em}
    \noindent{\Large \color{articolor} \textbf{Arti:} Saat meminjam (menggunakan) toilet, mari menyapa/meminta izin kasir terlebih dahulu.}
\end{convobox}

\begin{convobox}{71}
    \noindent
    \jpkana{「すみません、}{"Sumimasen,} 
    \jpkana{トイレを}{toire o} 
    \jpword{か}{借}{ka}\jpkana{りてもいいですか？」}{rite mo ii desu ka?"}

    \vspace{1.5em}
    \noindent{\Large \color{articolor} \textbf{Arti:} "Maaf (Permisi), bolehkah saya meminjam toilet?"}
\end{convobox}

\begin{convobox}{72}
    \noindent
    \jpkana{トイレの}{Toire no} 
    \jpword{りよう}{利用}{riyou} \jpkana{は}{wa} 
    \jpword{むりょう}{無料}{muryou} \jpkana{ですが、}{desu ga,}

    \vspace{1.5em}
    \noindent{\Large \color{articolor} \textbf{Arti:} Penggunaan toilet itu gratis, tapi}
\end{convobox}

\begin{convobox}{73}
    \noindent
    \jpkana{できたら}{dekitara} 
    \jpword{なに}{何}{nani}\jpkana{か}{ka} 
    \jpword{しょうひん}{商品}{shouhin} \jpkana{を}{o} 
    \jpword{か}{買}{ka}\jpkana{って}{tte} 
    \jpword{みせ}{店}{mise} \jpkana{を}{o} 
    \jpword{で}{出}{de}\jpkana{ると}{ru to} 
    \jpkana{いいですね。}{ii desu ne.}

    \vspace{1.5em}
    \noindent{\Large \color{articolor} \textbf{Arti:} jika memungkinkan, lebih baik membeli suatu produk (apa saja) sebelum keluar dari toko ya.}
\end{convobox}

\begin{lessonbox}{Aksi Pahlawan Konbini: Sopan Santun Toilet!}
    \textbf{Penjelasan:} Di Jepang, toilet Konbini adalah penyelamat turis! Tapi ingat tata kramanya ya.
    \\
    \textbf{🎬 ACTION / PRAKTEK:} 
    Jangan langsung nyelonong masuk! Hampiri kasir, katakan \textbf{"Toire o karite mo ii desu ka?"}. Setelah selesai buang air, jangan langsung pergi. Ambil air minum botol seharga 100 yen atau permen kecil, lalu bayar di kasir sebagai rasa terima kasih. Ini adalah etika yang sangat keren!
\end{lessonbox}

% ================= TITLE: ROLE PLAY =================
\vspace{2em}
\begin{tcolorbox}[colback=boxframe, colframe=boxframe, rounded corners, boxrule=0pt, arc=3mm, left=2mm, top=2mm, bottom=2mm]
    \Large \bfseries \color{white} 🗣️ 9. Waktunya Role Play! (Praktik Ucapkan)
\end{tcolorbox}
\vspace{1em}

\begin{convobox}{74}
    \noindent
    \jpkana{それでは、}{Sore dewa,} 
    \jpword{いま}{今}{ima} \jpkana{から}{kara} 
    \jpkana{ロールプレイを}{rooru purei o} \jpkana{します。}{shimasu.}

    \vspace{1.5em}
    \noindent{\Large \color{articolor} \textbf{Arti:} Kalau begitu, sekarang kita akan melakukan roleplay (bermain peran).}
\end{convobox}

\begin{convobox}{75}
    \noindent
    \jpkana{こんかいは、}{Konkai wa,} 
    \jpkana{みなさんは}{minasan wa} 
    \jpkana{お}{o}\jpword{きゃく}{客}{kyaku}\jpkana{さんに}{san ni} 
    \jpkana{なったつもりで、}{natta tsumori de,} 
    \jpword{こえ}{声}{koe} \jpkana{に}{ni} 
    \jpword{だ}{出}{da}\jpkana{して}{shite} 
    \jpword{はな}{話}{hana}\jpkana{してみてください。}{shite mite kudasai.}

    \vspace{1.5em}
    \noindent{\Large \color{articolor} \textbf{Arti:} Kali ini, bayangkanlah dirimu sebagai pembeli, dan cobalah berbicara dengan suara yang lantang.}
\end{convobox}

\begin{lessonbox}{⚠️ TUGAS TERAKHIR PAHLAWAN KONBINI! ⚠️}
    Anggap saja kamu sedang berdiri di depan kasir Jepang membawa kotak makan siang. Baca keras-keras Romaji di bagian bawah (yang menjadi jawabanmu). \textbf{Jangan dibaca dalam hati!} Mulutmu harus terbiasa mengucapkan kata-kata ini agar lancar saat di Jepang nanti. Ayo kita mulai!
\end{lessonbox}

\begin{convobox}{76}
    \noindent
    \jpkana{【店員】いらっしゃいませ。}{[Tenin] Irasshaimase.} 
    \jpkana{こちら}{Kochira} \jpword{あたた}{温}{atata}\jpkana{めますか？}{memasu ka?}

    \vspace{1.5em}
    \noindent{\Large \color{articolor} \textbf{Arti:} [Kasir] Selamat datang. Apakah mau dipanaskan?}
\end{convobox}

\begin{convobox}{77}
    \noindent
    \jpkana{【あなた】}{[Anata]} \jpkana{はい、}{Hai,} 
    \jpkana{お}{o}\jpword{ねが}{願}{nega}\jpkana{いします。}{ishimasu.}

    \vspace{1.5em}
    \noindent{\Large \color{articolor} \textbf{Arti:} [Kamu] Ya, tolong.}
\end{convobox}

\begin{convobox}{78}
    \noindent
    \jpkana{【店員】お}{[Tenin] O}\jpword{はし}{箸}{hashi} \jpkana{は}{wa} 
    \jpword{りよう}{利用}{riyou} \jpkana{ですか？}{desu ka?}

    \vspace{1.5em}
    \noindent{\Large \color{articolor} \textbf{Arti:} [Kasir] Pakai sumpit?}
\end{convobox}

\begin{convobox}{79}
    \noindent
    \jpkana{【あなた】}{[Anata]} \jpkana{はい、}{Hai,} 
    \jpkana{お}{o}\jpword{ねが}{願}{nega}\jpkana{いします。}{ishimasu.}

    \vspace{1.5em}
    \noindent{\Large \color{articolor} \textbf{Arti:} [Kamu] Ya, tolong.}
\end{convobox}

\begin{convobox}{80}
    \noindent
    \jpkana{【店員】}{[Tenin] }\jpword{ふくろ}{袋}{Fukuro} \jpkana{は}{wa} 
    \jpword{りよう}{利用}{riyou} \jpkana{ですか？}{desu ka?}

    \vspace{1.5em}
    \noindent{\Large \color{articolor} \textbf{Arti:} [Kasir] Butuh kantong plastik?}
\end{convobox}

\begin{convobox}{81}
    \noindent
    \jpkana{【あなた】}{[Anata]} \jpkana{いえ、}{Ie,} \jpword{だいじょうぶ}{大丈夫}{daijoubu} \jpkana{です。}{desu.}

    \vspace{1.5em}
    \noindent{\Large \color{articolor} \textbf{Arti:} [Kamu] (Aku bawa tas kain, jadi) Tidak usah.}
\end{convobox}

\begin{convobox}{82}
    \noindent
    \jpkana{【店員】お}{[Tenin] O}\jpword{かいけい}{会計}{kaikei} 
    \jpword{ろっぴゃくえん}{600円}{roppyaku-en} 
    \jpkana{でございます。}{de gozaimasu.}

    \vspace{1.5em}
    \noindent{\Large \color{articolor} \textbf{Arti:} [Kasir] Totalnya 600 yen.}
\end{convobox}

\begin{convobox}{83}
    \noindent
    \jpkana{【あなた】クレジットカードで}{[Anata] Kurejitto kaado de} 
    \jpkana{お}{o}\jpword{ねが}{願}{nega}\jpkana{いします。}{ishimasu.}

    \vspace{1.5em}
    \noindent{\Large \color{articolor} \textbf{Arti:} [Kamu] Pakai Kartu Kredit, tolong.}
\end{convobox}

\begin{convobox}{84}
    \noindent
    \jpkana{【店員】かしこまりました。}{[Tenin] Kashikomarimashita.} 
    \jpkana{レシート}{Reshiito} \jpword{りよう}{利用}{riyou} \jpkana{ですか？}{desu ka?}

    \vspace{1.5em}
    \noindent{\Large \color{articolor} \textbf{Arti:} [Kasir] Baik. Butuh struk?}
\end{convobox}

\begin{convobox}{85}
    \noindent
    \jpkana{【あなた】}{[Anata]} \jpkana{はい、}{Hai,} 
    \jpkana{お}{o}\jpword{ねが}{願}{nega}\jpkana{いします。}{ishimasu.}

    \vspace{1.5em}
    \noindent{\Large \color{articolor} \textbf{Arti:} [Kamu] Ya, tolong (buat kenang-kenangan!).}
\end{convobox}

\begin{convobox}{86}
    \noindent
    \jpkana{【店員】ありがとうございました。}{[Tenin] Arigatou gozaimashita.} 
    \jpkana{また}{Mata} \jpkana{お}{o}\jpword{こ}{越}{ko}\jpkana{しくださいませ。}{shi kudasaimase.}

    \vspace{1.5em}
    \noindent{\Large \color{articolor} \textbf{Arti:} [Kasir] Terima kasih banyak. Silakan datang kembali.}
\end{convobox}

% --- SUMMARY BOX ---
\vspace{2em}
\begin{tcolorbox}[
    colback=blue!5!white,
    colframe=blue!80!black,
    boxrule=3pt,
    arc=5mm,
    title={\huge\bfseries \sffamily 🌟 Rangkuman Misi Konbini 🌟},
    fonttitle=\color{white},
    attach boxed title to top center={yshift=-4mm},
    boxed title style={colback=blue!80!black, rounded corners},
    left=5mm, right=5mm, top=7mm, bottom=5mm
]
\Large \textbf{Sempurna! Kamu telah lulus ujian Konbini Jepang!}
\begin{itemize}
    \item Ingat 2 Jawaban Sakti: \textbf{Hai, onegaishimasu} (Ya tolong) dan \textbf{Ie, daijoubu desu} (Tidak usah).
    \item Dengarkan kata kunci: \textit{Atatamemasu ka} (Panaskan?), \textit{O-hashi} (Sumpit), \textit{Fukuro} (Plastik), \textit{Pointo Kaado} (Kartu poin), dan \textit{Reshiito} (Struk).
    \item Cara bayar: Sebutkan nama alat bayarnya + \textit{De onegaishimasu}.
    \item Toilet darurat: Selalu sapa kasir \textit{Toire o karitemo ii desu ka?} dan belilah barang kecil sebagai bentuk kesopanan.
\end{itemize}
Sekarang kamu sudah 100\% siap berbelanja di minimarket Jepang tanpa perlu merasa panik. Terus semangat belajarnya, dan \textit{Mata jikai no douga de o-ai shimashou!} (Sampai jumpa di video/materi selanjutnya!)
\end{tcolorbox}

\newpage
\chapter{Ordering Ramen in Japanese}

% ================= INTRO =================
\begin{convobox}{1}
    \noindent
    \jpkana{みなさん}{Minasan} \jpkana{こんにちは、}{konnichiwa,} 
    \jpkana{たなかです。}{Tanaka desu.}

    \vspace{1.5em}
    \noindent{\Large \color{articolor} \textbf{Arti:} Halo semuanya, saya Tanaka.}
\end{convobox}

\begin{convobox}{2}
    \noindent
    \jpword{きょう}{今日}{Kyou} \jpkana{は、}{wa,} 
    \jpkana{ラーメン}{raamen}\jpword{や}{屋}{ya} \jpkana{で}{de} 
    \jpword{ちゅうもん}{注文}{chuumon} \jpkana{するときの}{suru toki no} 
    \jpword{にほんご}{日本語}{Nihongo} \jpkana{を}{o} 
    \jpword{いっしょ}{一緒}{issho} \jpkana{に}{ni} 
    \jpword{べんきょう}{勉強}{benkyou} \jpkana{しましょう。}{shimashou.}

    \vspace{1.5em}
    \noindent{\Large \color{articolor} \textbf{Arti:} Hari ini, mari kita belajar bersama bahasa Jepang saat memesan di kedai ramen.}
\end{convobox}

\begin{convobox}{3}
    \noindent
    \jpword{とつぜん}{突然}{Totsuzen} \jpkana{ですが、}{desu ga,} 
    \jpkana{みなさんは}{minasan wa} 
    \jpkana{ラーメン}{raamen} \jpkana{が}{ga} 
    \jpword{す}{好}{su}\jpkana{きですか？}{ki desu ka?} 
    \jpword{わたし}{私}{Watashi} \jpkana{は}{wa} 
    \jpword{だいす}{大好}{daisu}\jpkana{きです。}{ki desu.}

    \vspace{1.5em}
    \noindent{\Large \color{articolor} \textbf{Arti:} Tiba-tiba (Oiya), apakah kalian suka ramen? Kalau saya sangat suka.}
\end{convobox}

\begin{convobox}{4}
    \noindent
    \jpword{にほん}{日本}{Nihon} \jpkana{に}{ni} 
    \jpword{き}{来}{ki}\jpkana{たら、}{tara,} 
    \jpkana{ラーメン}{raamen}\jpword{や}{屋}{ya}\jpword{めぐ}{巡}{megu}\jpkana{りを}{ri o} 
    \jpkana{したいという}{shitai to iu} 
    \jpword{かた}{方}{kata} \jpkana{も}{mo} 
    \jpword{おお}{多}{oo}\jpkana{いと}{i to} 
    \jpword{おも}{思}{omo}\jpkana{います。}{imasu.}

    \vspace{1.5em}
    \noindent{\Large \color{articolor} \textbf{Arti:} Jika datang ke Jepang, aku rasa banyak orang yang ingin berwisata kuliner (berkeliling) kedai ramen.}
\end{convobox}

\begin{convobox}{5}
    \noindent
    \jpkana{でも、}{Demo,} 
    \jpkana{「}{"}\jpword{にほんご}{日本語}{Nihongo} \jpkana{で}{de} 
    \jpkana{ラーメン}{raamen} \jpkana{を}{o} 
    \jpword{ちゅうもん}{注文}{chuumon} \jpkana{するのは、}{suru no wa,} 
    \jpword{すこ}{少}{suko}\jpkana{し}{shi} 
    \jpword{むずか}{難}{muzuka}\jpkana{しそうだなぁ」}{shisou da naa"}

    \vspace{1.5em}
    \noindent{\Large \color{articolor} \textbf{Arti:} Tapi, "Memesan ramen dalam bahasa Jepang, sepertinya sedikit sulit ya..."}
\end{convobox}

\begin{convobox}{6}
    \noindent
    \jpkana{...と}{...to} 
    \jpword{おも}{思}{omo}\jpkana{っている}{tte iru} 
    \jpword{かた}{方}{kata} \jpkana{も}{mo} 
    \jpkana{いるのではないでしょうか。}{iru no de wa nai deshou ka.}

    \vspace{1.5em}
    \noindent{\Large \color{articolor} \textbf{Arti:} ...mungkin ada juga yang berpikir seperti itu.}
\end{convobox}

\begin{convobox}{7}
    \noindent
    \jpkana{でも、}{Demo,} 
    \jpword{あんしん}{安心}{anshin} \jpkana{してください。}{shite kudasai.} 
    \jpkana{この}{Kono} \jpword{どうが}{動画}{douga} \jpkana{を}{o} 
    \jpword{み}{観}{mi}\jpkana{れば、}{reba,} 
    \jpkana{ラーメン}{raamen} \jpkana{の}{no} 
    \jpword{ちゅうもん}{注文}{chuumon} \jpword{ほうほう}{方法}{houhou} \jpkana{が}{ga} 
    \jpword{わ}{分}{wa}\jpkana{かります。}{karimasu.}

    \vspace{1.5em}
    \noindent{\Large \color{articolor} \textbf{Arti:} Tapi, tenang saja (jangan khawatir). Kalau menonton video ini, kamu akan mengerti cara memesan ramen.}
\end{convobox}

\begin{convobox}{8}
    \noindent
    \jpkana{さらに、}{Sara ni,} 
    \jpkana{トッピングを}{toppingu o} 
    \jpword{ついか}{追加}{tsuika} \jpkana{したり、}{shitari,} 
    \jpword{めん}{麺}{men} \jpkana{を}{o} 
    \jpword{この}{この}{kono}\jpkana{みの}{mi no} 
    \jpkana{かたさにしてもらうように}{katasa ni shite morau you ni} 
    \jpword{つた}{伝}{tsuta}\jpkana{えたりして、}{etari shite,}

    \vspace{1.5em}
    \noindent{\Large \color{articolor} \textbf{Arti:} Terlebih lagi, (kamu bisa) menambah topping, menyampaikan tingkat kekerasan mi sesuai seleramu,}
\end{convobox}

\begin{convobox}{9}
    \noindent
    \jpkana{もっと}{motto} \jpkana{ラーメン}{raamen} \jpkana{を}{o} 
    \jpkana{おいしく}{oishiku} 
    \jpword{た}{食}{ta}\jpkana{べることも}{beru koto mo} 
    \jpkana{できるようになりますよ。}{dekiru you ni narimasu yo.}

    \vspace{1.5em}
    \noindent{\Large \color{articolor} \textbf{Arti:} sehingga kamu akan bisa memakan ramen dengan lebih lezat lagi lho.}
\end{convobox}

\begin{convobox}{10}
    \noindent
    \jpkana{それでは}{Sore dewa} 
    \jpword{さっそく}{早速}{sassoku} 
    \jpword{い}{行}{i}\jpkana{ってみましょう！}{tte mimashou!}

    \vspace{1.5em}
    \noindent{\Large \color{articolor} \textbf{Arti:} Kalau begitu, ayo langsung kita mulai!}
\end{convobox}

% ================= TITLE: CONVERSATION =================
\vspace{2em}
\begin{tcolorbox}[colback=boxframe, colframe=boxframe, rounded corners, boxrule=0pt, arc=3mm, left=2mm, top=2mm, bottom=2mm]
    \Large \bfseries \color{white} 🍜 1. Conversation (Percakapan Dasar Memesan Ramen)
\end{tcolorbox}
\vspace{1em}

\begin{convobox}{11}
    \noindent
    \jpword{まず}{先ず}{Mazu} \jpword{いちど}{一度}{ichido} 
    \jpkana{ラーメン}{raamen}\jpword{や}{屋}{ya} \jpkana{での}{de no} 
    \jpword{かいわ}{会話}{kaiwa} \jpkana{を}{o} 
    \jpword{しぜん}{自然}{shizen} \jpkana{な}{na} 
    \jpkana{スピードで}{supiido de} 
    \jpword{き}{聞}{ki}\jpkana{いてみましょう。}{ite mimashou.}

    \vspace{1.5em}
    \noindent{\Large \color{articolor} \textbf{Arti:} Pertama-tama, mari kita dengarkan sekali percakapan di kedai ramen dengan kecepatan natural.}
\end{convobox}

\begin{convobox}{12}
    \noindent
    \jpkana{いらっしゃいませ、}{Irasshaimase,} 
    \jpword{なんめい}{何名}{nan-mei}\jpword{さま}{様}{sama} \jpkana{ですか？}{desu ka?}

    \vspace{1.5em}
    \noindent{\Large \color{articolor} \textbf{Arti:} Selamat datang, untuk berapa orang?}
\end{convobox}

\begin{convobox}{13}
    \noindent
    \jpword{ひとり}{1人}{Hitori} \jpkana{です。}{desu.}

    \vspace{1.5em}
    \noindent{\Large \color{articolor} \textbf{Arti:} 1 orang.}
\end{convobox}

\begin{convobox}{14}
    \noindent
    \jpkana{こちらへ}{Kochira e} \jpkana{どうぞ。}{douzo.}

    \vspace{1.5em}
    \noindent{\Large \color{articolor} \textbf{Arti:} Silakan ke sebelah sini.}
\end{convobox}

\begin{convobox}{15}
    \noindent
    \jpkana{しょうゆラーメン}{Shouyu raamen} 
    \jpkana{お}{o}\jpword{ねが}{願}{nega}\jpkana{いします。}{ishimasu.}

    \vspace{1.5em}
    \noindent{\Large \color{articolor} \textbf{Arti:} Tolong pesan ramen rasa Shoyu (kecap asin).}
\end{convobox}

\begin{convobox}{16}
    \noindent
    \jpword{めん}{麺}{Men} \jpkana{の}{no} 
    \jpkana{かたさは}{katasa wa} 
    \jpkana{どうなさいますか？}{dou nasaimasu ka?}

    \vspace{1.5em}
    \noindent{\Large \color{articolor} \textbf{Arti:} Tingkat kematangan (kekerasan) mi-nya bagaimana?}
\end{convobox}

\begin{convobox}{17}
    \noindent
    \jpword{めん}{麺}{Men} 
    \jpkana{かためで}{katame de} 
    \jpkana{お}{o}\jpword{ねが}{願}{nega}\jpkana{いします。}{ishimasu.}

    \vspace{1.5em}
    \noindent{\Large \color{articolor} \textbf{Arti:} Tolong mi-nya agak keras (setengah matang).}
\end{convobox}

\begin{convobox}{18}
    \noindent
    \jpkana{かしこまりました。}{Kashikomarimashita.}

    \vspace{1.5em}
    \noindent{\Large \color{articolor} \textbf{Arti:} Baik, saya mengerti.}
\end{convobox}

\begin{convobox}{19}
    \noindent
    \jpkana{トッピングは}{Toppingu wa} 
    \jpkana{いかがですか？}{ikaga desu ka?}

    \vspace{1.5em}
    \noindent{\Large \color{articolor} \textbf{Arti:} Bagaimana dengan tambahan topping-nya?}
\end{convobox}

\begin{convobox}{20}
    \noindent
    \jpword{あじたま}{味玉}{Ajitama} 
    \jpkana{トッピング}{toppingu} 
    \jpkana{お}{o}\jpword{ねが}{願}{nega}\jpkana{いします。}{ishimasu.}

    \vspace{1.5em}
    \noindent{\Large \color{articolor} \textbf{Arti:} Tolong tambah topping telur rebus (Ajitama).}
\end{convobox}

\begin{convobox}{21}
    \noindent
    \jpkana{かしこまりました。}{Kashikomarimashita.}

    \vspace{1.5em}
    \noindent{\Large \color{articolor} \textbf{Arti:} Baik.}
\end{convobox}

\begin{convobox}{22}
    \noindent
    \jpkana{それでは、}{Sore dewa,} 
    \jpword{じゅんばん}{順番}{junban} \jpkana{に}{ni} 
    \jpword{せつめい}{説明}{setsumei} \jpkana{しますね。}{shimasu ne.}

    \vspace{1.5em}
    \noindent{\Large \color{articolor} \textbf{Arti:} Kalau begitu, akan saya jelaskan secara berurutan ya.}
\end{convobox}

\begin{lessonbox}{Tips Belanja untuk Anak \& Orang Tua: Jurus "Kore, Onegaishimasu"}
    \textbf{TIPS UMUM:} Saat masuk ke restoran mana saja di Jepang, pelayan akan langsung menyambut dengan \textit{Irasshaimase} dan membawa menu. 
    \\
    \textbf{JURUS RAHASIA:} Kalau kamu belum tahu cara membaca menunya, jangan khawatir! Cukup \textbf{tunjuk gambarnya} di buku menu dan ucapkan: \textbf{"Kore, onegaishimasu"} (Tolong pesan yang INI). Begitu mudah!
\end{lessonbox}

% ================= TITLE: PHRASE 1 =================
\vspace{2em}
\begin{tcolorbox}[colback=boxframe, colframe=boxframe, rounded corners, boxrule=0pt, arc=3mm, left=2mm, top=2mm, bottom=2mm]
    \Large \bfseries \color{white} 👥 2. Phrase 1 (Cara Menghitung Jumlah Orang)
\end{tcolorbox}
\vspace{1em}

\begin{convobox}{23}
    \noindent
    \jpkana{「いらっしゃいませ、}{"Irasshaimase,} 
    \jpword{なんめい}{何名}{nan-mei}\jpword{さま}{様}{sama} \jpkana{ですか？」}{desu ka?"}

    \vspace{1.5em}
    \noindent{\Large \color{articolor} \textbf{Arti:} "Selamat datang, untuk berapa orang?"}
\end{convobox}

\begin{convobox}{24}
    \noindent
    \jpkana{「}{"}\jpword{ひとり}{1人}{Hitori} \jpkana{です。」}{desu."}

    \vspace{1.5em}
    \noindent{\Large \color{articolor} \textbf{Arti:} "1 orang."}
\end{convobox}

\begin{convobox}{25}
    \noindent
    \jpkana{お}{O}\jpword{みせ}{店}{mise} \jpkana{の}{no} 
    \jpword{ひと}{人}{hito} \jpkana{に}{ni} 
    \jpword{にんずう}{人数}{ninzuu} \jpkana{を}{o} 
    \jpword{つた}{伝}{tsuta}\jpkana{える}{eru} 
    \jpkana{フレーズです。}{fureezu desu.}

    \vspace{1.5em}
    \noindent{\Large \color{articolor} \textbf{Arti:} Ini adalah frasa untuk menyampaikan jumlah orang kepada staf toko/restoran.}
\end{convobox}

\begin{convobox}{26}
    \noindent
    \jpword{にんずう}{人数}{Ninzuu} \jpkana{は、}{wa,} 
    \jpkana{このように}{kono you ni} 
    \jpword{い}{言}{i}\jpkana{います。}{imasu.}

    \vspace{1.5em}
    \noindent{\Large \color{articolor} \textbf{Arti:} Jumlah orang, diucapkan seperti ini.}
\end{convobox}

\begin{convobox}{27}
    \noindent
    \jpword{ひとり}{1人}{Hitori} / \jpword{ふたり}{2人}{Futari} / \jpword{さんにん}{3人}{San-nin} / \jpword{よにん}{4人}{Yo-nin} / \jpword{ごにん}{5人}{Go-nin}

    \vspace{1.5em}
    \noindent{\Large \color{articolor} \textbf{Arti:} 1 orang / 2 orang / 3 orang / 4 orang / 5 orang}
\end{convobox}

\begin{convobox}{28}
    \noindent
    \jpkana{このように}{Kono you ni} 
    \jpword{ゆび}{指}{yubi} \jpkana{も}{mo} 
    \jpword{いっしょ}{一緒}{issho} \jpkana{に}{ni} 
    \jpword{だ}{出}{da}\jpkana{すと、}{su to,} 
    \jpkana{お}{o}\jpword{みせ}{店}{mise} \jpkana{の}{no} 
    \jpword{ひと}{人}{hito} \jpkana{に}{ni} 
    \jpword{わ}{分}{wa}\jpkana{かりやすいです。}{kariyasui desu.}

    \vspace{1.5em}
    \noindent{\Large \color{articolor} \textbf{Arti:} Jika kamu juga menunjukkan jari (sesuai jumlah) seperti ini, staf toko akan lebih mudah memahaminya.}
\end{convobox}

\begin{lessonbox}{Hafalan Wajib: Menghitung Anggota Keluarga!}
    \textbf{Orang Tua \& Anak:} Di Jepang, menghitung orang itu sedikit unik. 1 dan 2 orang panggilannya spesial!
    \begin{itemize}
        \item Sendiri = \textbf{Hitori}
        \item Berdua = \textbf{Futari}
        \item Bertiga = \textbf{San-nin} (Tinggal tambah "nin" di belakang angka)
        \item Berempat = \textbf{Yo-nin} (Bukan Yon-nin lho ya!)
    \end{itemize}
    \textbf{Aksi:} Saat pelayan tanya \textit{"Nan-mei sama desu ka?"}, angkat jarimu tinggi-tinggi dan ucapkan jumlah keluargamu dengan bangga! "Yo-nin desu!" (4 orang).
\end{lessonbox}

\begin{convobox}{29}
    \noindent
    \jpkana{「しょうゆラーメン}{"Shouyu raamen} 
    \jpkana{お}{o}\jpword{ねが}{願}{nega}\jpkana{いします。」}{ishimasu."}

    \vspace{1.5em}
    \noindent{\Large \color{articolor} \textbf{Arti:} "Tolong ramen rasa shoyu."}
\end{convobox}

\begin{convobox}{30}
    \noindent
    \jpkana{ラーメンの}{Raamen no} 
    \jpword{しゅるい}{種類}{shurui} \jpkana{を}{o} 
    \jpword{つた}{伝}{tsuta}\jpkana{える}{eru} 
    \jpkana{フレーズです。}{fureezu desu.}

    \vspace{1.5em}
    \noindent{\Large \color{articolor} \textbf{Arti:} Ini adalah frasa untuk menyebutkan jenis ramen.}
\end{convobox}

\begin{convobox}{31}
    \noindent
    \jpkana{ラーメンの}{Raamen no} \jpword{しゅるい}{種類}{shurui} \jpkana{には、}{ni wa,} 
    \jpword{たと}{例}{tato}\jpkana{えば}{eba} 
    \jpkana{このようなものが}{kono you na mono ga} 
    \jpkana{あります。}{arimasu.}

    \vspace{1.5em}
    \noindent{\Large \color{articolor} \textbf{Arti:} Untuk jenis ramen, misalnya ada yang seperti ini.}
\end{convobox}

\begin{convobox}{32}
    \noindent
    \jpkana{しょうゆラーメン}{Shouyu raamen} / 
    \jpkana{とんこつラーメン}{Tonkotsu raamen} / 
    \jpword{しお}{塩}{Shio} \jpkana{ラーメン}{raamen} / 
    \jpkana{みそラーメン}{Miso raamen}

    \vspace{1.5em}
    \noindent{\Large \color{articolor} \textbf{Arti:} Ramen Shoyu (Kecap asin) / Ramen Tonkotsu (Kaldu tulang) / Ramen Shio (Garam) / Ramen Miso (Pasta kedelai)}
\end{convobox}

\begin{convobox}{33}
    \noindent
    \jpkana{ひとつだけでなく、}{Hitotsu dake de naku,} 
    \jpkana{いくつか}{ikutsu ka} 
    \jpword{ちゅうもん}{注文}{chuumon} \jpkana{するときは、}{suru toki wa,} 
    \jpkana{このように}{kono you ni} 
    \jpword{い}{言}{i}\jpkana{うことができます。}{u koto ga dekimasu.}

    \vspace{1.5em}
    \noindent{\Large \color{articolor} \textbf{Arti:} Jika memesan bukan cuma satu (tapi beberapa), kamu bisa mengucapkannya seperti ini.}
\end{convobox}

\begin{convobox}{34}
    \noindent
    \jpkana{「しょうゆラーメン}{"Shouyu raamen} \jpkana{ひとつ}{hitotsu} \jpkana{と、}{to,} 
    \jpkana{とんこつラーメン}{tonkotsu raamen} \jpkana{ふたつ}{futatsu} 
    \jpkana{お}{o}\jpword{ねが}{願}{nega}\jpkana{いします。」}{ishimasu."}

    \vspace{1.5em}
    \noindent{\Large \color{articolor} \textbf{Arti:} "Tolong ramen shoyu 1 (mangkuk), dan ramen tonkotsu 2 (mangkuk)."}
\end{convobox}

\begin{convobox}{35}
    \noindent
    \jpkana{「しょうゆラーメン}{"Shouyu raamen} \jpkana{と、}{to,} 
    \jpword{しお}{塩}{shio}\jpkana{ラーメンと、}{raamen to,} 
    \jpkana{みそラーメン}{miso raamen} 
    \jpkana{お}{o}\jpword{ねが}{願}{nega}\jpkana{いします。」}{ishimasu."}

    \vspace{1.5em}
    \noindent{\Large \color{articolor} \textbf{Arti:} "Tolong ramen shoyu, dan ramen shio, dan ramen miso."}
\end{convobox}

\begin{lessonbox}{Hafalan Wajib: Menghitung Benda (Porsi/Mangkuk)!}
    Kalau tadi menghitung orang, sekarang menghitung benda bulat atau porsi mangkuk!
    \begin{itemize}
        \item 1 Porsi = \textbf{Hitotsu}
        \item 2 Porsi = \textbf{Futatsu}
        \item 3 Porsi = \textbf{Mittsu}
    \end{itemize}
    \textbf{Rumusnya:} [Nama Makanan] + [Jumlah] + [Onegaishimasu]. \\
    Contoh: "Ramen Hitotsu, Onegaishimasu!" (Tolong Ramen 1 porsi!)
\end{lessonbox}

% ================= TITLE: PHRASE 3 =================
\vspace{2em}
\begin{tcolorbox}[colback=boxframe, colframe=boxframe, rounded corners, boxrule=0pt, arc=3mm, left=2mm, top=2mm, bottom=2mm]
    \Large \bfseries \color{white} 🥢 4. Phrase 3 (Kekerasan Mi)
\end{tcolorbox}
\vspace{1em}

\begin{convobox}{36}
    \noindent
    \jpkana{「}{"}\jpword{めん}{麺}{Men} \jpkana{かためで}{katame de} 
    \jpkana{お}{o}\jpword{ねが}{願}{nega}\jpkana{いします。」}{ishimasu."}

    \vspace{1.5em}
    \noindent{\Large \color{articolor} \textbf{Arti:} "Tolong mi-nya agak keras (setengah matang)."}
\end{convobox}

\begin{convobox}{37}
    \noindent
    \jpkana{お}{O}\jpword{みせ}{店}{mise} \jpkana{によっては、}{ni yotte wa,} 
    \jpkana{このように}{kono you ni} 
    \jpword{めん}{麺}{men} \jpkana{の}{no} 
    \jpkana{かたさを}{katasa o} 
    \jpword{えら}{選}{era}\jpkana{ぶことができます。}{bu koto ga dekimasu.}

    \vspace{1.5em}
    \noindent{\Large \color{articolor} \textbf{Arti:} Tergantung tokonya, kamu bisa memilih tingkat kekerasan/kematangan mi seperti ini.}
\end{convobox}

\begin{convobox}{38}
    \noindent
    \jpkana{かためが}{Katame ga} \jpkana{いい}{ii} \jpword{とき}{時}{toki}\jpkana{は、}{wa,} 
    \jpkana{「}{"}\jpword{めん}{麺}{men} \jpkana{かためで}{katame de} 
    \jpkana{お}{o}\jpword{ねが}{願}{nega}\jpkana{いします。」}{ishimasu."}

    \vspace{1.5em}
    \noindent{\Large \color{articolor} \textbf{Arti:} Jika ingin yang agak keras (al dente), "Tolong mi-nya katame."}
\end{convobox}

\begin{convobox}{39}
    \noindent
    \jpkana{やわらかめが}{Yawarakame ga} \jpkana{いい}{ii} \jpword{とき}{時}{toki}\jpkana{は、}{wa,} 
    \jpkana{「}{"}\jpword{めん}{麺}{men} \jpkana{やわらかめで}{yawarakame de} 
    \jpkana{お}{o}\jpword{ねが}{願}{nega}\jpkana{いします。」}{ishimasu."}

    \vspace{1.5em}
    \noindent{\Large \color{articolor} \textbf{Arti:} Jika ingin yang lembek/lunak, "Tolong mi-nya yawarakame."}
\end{convobox}

\begin{convobox}{40}
    \noindent
    \jpword{ふつう}{普通}{Futsuu} \jpkana{の}{no} 
    \jpkana{かたさが}{katasa ga} \jpkana{いい}{ii} \jpword{とき}{時}{toki}\jpkana{は、}{wa,} 
    \jpkana{「}{"}\jpword{ふつう}{普通}{futsuu} \jpkana{で}{de} 
    \jpkana{お}{o}\jpword{ねが}{願}{nega}\jpkana{いします。」}{ishimasu."}

    \vspace{1.5em}
    \noindent{\Large \color{articolor} \textbf{Arti:} Jika ingin yang normal/standar, "Tolong yang futsuu (normal)."}
\end{convobox}

\begin{convobox}{41}
    \noindent
    \jpkana{ちなみに、}{Chinami ni,} 
    \jpword{わたし}{私}{watashi} \jpkana{は}{wa} 
    \jpword{めん}{麺}{men} \jpkana{かためが}{katame ga} 
    \jpword{す}{好}{su}\jpkana{きです。}{ki desu.}

    \vspace{1.5em}
    \noindent{\Large \color{articolor} \textbf{Arti:} Ngomong-ngomong, aku suka mi yang katame (agak keras/kenyal).}
\end{convobox}

\begin{lessonbox}{Aksi Jagoan Ramen: Pilih Kekerasan Mi!}
    Orang Jepang sangat serius soal mi ramen! Kalau pelayan menanyakan tentang *Katasa* (Kekerasan), kamu harus siap menjawab:
    \begin{itemize}
        \item \textbf{Katame:} Agak keras/kenyal
        \item \textbf{Futsuu:} Biasa/Standar.
        \item \textbf{Yawarakame:} Lembut/Lunak (Cocok untuk anak-anak kecil!).
    \end{itemize}
\end{lessonbox}

% ================= TITLE: PHRASE 4 =================
\vspace{2em}
\begin{tcolorbox}[colback=boxframe, colframe=boxframe, rounded corners, boxrule=0pt, arc=3mm, left=2mm, top=2mm, bottom=2mm]
    \Large \bfseries \color{white} 🥚 5. Phrase 4 (Tambah Topping)
\end{tcolorbox}
\vspace{1em}

\begin{convobox}{42}
    \noindent
    \jpkana{「}{"}\jpword{あじたま}{味玉}{Ajitama} \jpkana{トッピング}{toppingu} 
    \jpkana{お}{o}\jpword{ねが}{願}{nega}\jpkana{いします。」}{ishimasu."}

    \vspace{1.5em}
    \noindent{\Large \color{articolor} \textbf{Arti:} "Tolong topping telur rebus (kecap)."}
\end{convobox}

\begin{convobox}{43}
    \noindent
    \jpkana{ラーメンには、}{Raamen ni wa,} 
    \jpword{じぶん}{自分}{jibun} \jpkana{の}{no} 
    \jpword{この}{この}{kono}\jpkana{みに}{mi ni} 
    \jpword{あ}{合}{a}\jpkana{わせて}{wasete} 
    \jpkana{いろいろな}{iroiro na} 
    \jpkana{トッピングが}{toppingu ga} \jpkana{できます。}{dekimasu.}

    \vspace{1.5em}
    \noindent{\Large \color{articolor} \textbf{Arti:} Di ramen, kamu bisa melakukan bermacam topping sesuai dengan selera dirimu sendiri.}
\end{convobox}

\begin{convobox}{44}
    \noindent
    \jpkana{チャーシュー}{Chaashuu} (Daging iris) / \jpword{あじたま}{味玉}{Ajitama} (Telur rebus) / \jpkana{メンマ}{Menma} (Rebung) / \jpkana{ねぎ}{Negi} (Daun bawang) / \jpkana{のり}{Nori} (Rumput laut) / \jpkana{にんにく}{Ninniku} (Bawang putih) / \jpword{ほうれんそう}{ほうれん草}{Hourensou} (Bayam)
\end{convobox}

\begin{convobox}{45}
    \noindent
    \jpword{ちゅうもん}{注文}{Chuumon} \jpkana{するときは、}{suru toki wa,} 
    \jpkana{このように}{kono you ni} 
    \jpword{い}{言}{i}\jpkana{います。}{imasu.}

    \vspace{1.5em}
    \noindent{\Large \color{articolor} \textbf{Arti:} Saat memesan, ucapkan seperti ini.}
\end{convobox}

\begin{convobox}{46}
    \noindent
    \jpkana{「チャーシュー}{"Chaashuu} \jpkana{トッピング}{toppingu} 
    \jpkana{お}{o}\jpword{ねが}{願}{nega}\jpkana{いします。」}{ishimasu."}

    \vspace{1.5em}
    \noindent{\Large \color{articolor} \textbf{Arti:} "Tolong topping chaashuu."}
\end{convobox}

\begin{convobox}{47}
    \noindent
    \jpkana{「}{"}\jpword{あじたま}{味玉}{Ajitama} \jpkana{と、}{to,} 
    \jpkana{メンマと、}{menma to,} 
    \jpkana{にんにく}{ninniku} 
    \jpkana{トッピング}{toppingu} 
    \jpkana{お}{o}\jpword{ねが}{願}{nega}\jpkana{いします。」}{ishimasu."}

    \vspace{1.5em}
    \noindent{\Large \color{articolor} \textbf{Arti:} "Tolong topping ajitama, menma, dan ninniku."}
\end{convobox}

\begin{lessonbox}{Aksi Jagoan Ramen: Partikel "To" (Dan)}
    Untuk menggabungkan banyak topping, gunakan partikel \textbf{と (to)} yang artinya "Dan".\\
    Contoh: Ajitama \textbf{to}, Nori \textbf{to}, Negi \textbf{to}... onegaishimasu!
\end{lessonbox}

% ================= TITLE: TIPS =================
\vspace{2em}
\begin{tcolorbox}[colback=boxframe, colframe=boxframe, rounded corners, boxrule=0pt, arc=3mm, left=2mm, top=2mm, bottom=2mm]
    \Large \bfseries \color{white} 💡 6. Tips Rahasia Membayar: Mesin Tiket (Kenbaiki)
\end{tcolorbox}
\vspace{1em}

\begin{convobox}{48}
    \noindent
    \jpword{こんかい}{今回}{Konkai} \jpkana{は、}{wa,} 
    \jpword{てんいん}{店員}{tenin}\jpkana{さんに}{san ni} 
    \jpword{ちょくせつ}{直接}{chokusetsu} 
    \jpword{ちゅうもん}{注文}{chuumon} \jpkana{する}{suru} 
    \jpword{ほうほう}{方法}{houhou} \jpkana{を}{o} 
    \jpword{つた}{伝}{tsuta}\jpkana{えましたが、}{emashita ga,}

    \vspace{1.5em}
    \noindent{\Large \color{articolor} \textbf{Arti:} Kali ini, aku mengajarkan cara memesan langsung ke kasir/pelayan, namun}
\end{convobox}

\begin{convobox}{49}
    \noindent
    \jpkana{お}{O}\jpword{みせ}{店}{mise} \jpkana{によっては、}{ni yotte wa,} 
    \jpkana{このような}{kono you na} 
    \jpword{きかい}{機械}{kikai} \jpkana{で}{de} 
    \jpword{ちゅうもん}{注文}{chuumon} \jpkana{することもあります。}{suru koto mo arimasu.}

    \vspace{1.5em}
    \noindent{\Large \color{articolor} \textbf{Arti:} Tergantung tokonya, ada kalanya kamu memesan lewat mesin seperti ini.}
\end{convobox}

\begin{convobox}{50}
    \noindent
    \jpkana{この}{Kono} \jpword{きかい}{機械}{kikai} \jpkana{のことを、}{no koto o,} 
    \jpkana{「}{"}\jpword{けんばいき}{券売機}{kenbaiki}\jpkana{」と}{" to} 
    \jpword{い}{言}{i}\jpkana{います。}{imasu.}

    \vspace{1.5em}
    \noindent{\Large \color{articolor} \textbf{Arti:} Mesin ini, disebut "Kenbaiki" (Mesin penjual tiket otomatis).}
\end{convobox}

\begin{convobox}{51}
    \noindent
    \jpword{けんばいき}{券売機}{Kenbaiki} \jpkana{で}{de} 
    \jpword{ちゅうもん}{注文}{chuumon} \jpkana{するときは、}{suru toki wa,} 
    \jpword{まず}{先ず}{mazu} \jpkana{お}{o}\jpword{かね}{金}{kane} \jpkana{を}{o} 
    \jpword{い}{入}{i}\jpkana{れて、}{rete,}

    \vspace{1.5em}
    \noindent{\Large \color{articolor} \textbf{Arti:} Saat memesan di mesin, pertama masukkan uangnya,}
\end{convobox}

\begin{convobox}{52}
    \noindent
    \jpword{ちゅうもん}{注文}{chuumon} \jpkana{したい}{shitai} 
    \jpkana{ラーメンや}{raamen ya} 
    \jpkana{トッピングの}{toppingu no} 
    \jpkana{ボタンを}{botan o} \jpword{お}{押}{o}\jpkana{して、}{shite,}

    \vspace{1.5em}
    \noindent{\Large \color{articolor} \textbf{Arti:} lalu pencet tombol ramen atau topping yang ingin dipesan,}
\end{convobox}

\begin{convobox}{53}
    \noindent
    \jpword{で}{出}{de}\jpkana{てきた}{tekita} 
    \jpword{かみ}{紙}{kami} \jpkana{を}{o} 
    \jpkana{お}{o}\jpword{みせ}{店}{mise} \jpkana{の}{no} 
    \jpword{ひと}{人}{hito} \jpkana{に}{ni} 
    \jpword{わた}{渡}{wata}\jpkana{します。}{shimasu.}

    \vspace{1.5em}
    \noindent{\Large \color{articolor} \textbf{Arti:} lalu serahkan kertas (tiket) yang keluar itu kepada staf toko.}
\end{convobox}

\begin{convobox}{54}
    \noindent
    \jpkana{あとは、}{Ato wa,} 
    \jpkana{ラーメンが}{raamen ga} 
    \jpword{く}{来}{ku}\jpkana{るのを}{ru no o} 
    \jpword{ま}{待}{ma}\jpkana{つだけです。}{tsu dake desu.}

    \vspace{1.5em}
    \noindent{\Large \color{articolor} \textbf{Arti:} Setelahnya, kamu tinggal menunggu ramennya datang.}
\end{convobox}

\begin{lessonbox}{Aksi Penting Kenbaiki (Mesin Tiket) 🚨}
    \textbf{AWAS TERBALIK!} Ini kesalahan turis paling umum! \\
    Di Indonesia atau mesin minuman, kita tekan tombol dulu baru masukin uang kan? Di mesin ramen Jepang (\textbf{Kenbaiki}), kamu harus \textbf{MASUKKAN UANG DULU}, barulah lampu-lampu tombol menunya menyala untuk bisa ditekan! Setelah tiketnya keluar, berikan tiketnya ke pelayan tanpa perlu bicara apa-apa. Gampang banget!
\end{lessonbox}

% ================= TITLE: ROLE PLAY =================
\vspace{2em}
\begin{tcolorbox}[colback=boxframe, colframe=boxframe, rounded corners, boxrule=0pt, arc=3mm, left=2mm, top=2mm, bottom=2mm]
    \Large \bfseries \color{white} 🗣️ 7. Waktunya Role Play! (Praktik Ucapkan)
\end{tcolorbox}
\vspace{1em}

\begin{convobox}{55}
    \noindent
    \jpkana{【店員】いらっしゃいませ、}{[Tenin] Irasshaimase,} 
    \jpword{なんめい}{何名}{nan-mei}\jpword{さま}{様}{sama} \jpkana{ですか？}{desu ka?}

    \vspace{1.5em}
    \noindent{\Large \color{articolor} \textbf{Arti:} [Pelayan] Selamat datang, untuk berapa orang?}
\end{convobox}

\begin{convobox}{56}
    \noindent
    \jpkana{【あなた】}{[Anata]} \jpword{ひとり}{1人}{Hitori} \jpkana{です。}{desu.}

    \vspace{1.5em}
    \noindent{\Large \color{articolor} \textbf{Arti:} [Kamu] 1 orang. (Jangan lupa angkat 1 jari!)}
\end{convobox}

\begin{convobox}{57}
    \noindent
    \jpkana{【店員】こちらへ}{[Tenin] Kochira e} \jpkana{どうぞ。}{douzo.}

    \vspace{1.5em}
    \noindent{\Large \color{articolor} \textbf{Arti:} [Pelayan] Silakan ke sebelah sini.}
\end{convobox}

\begin{convobox}{58}
    \noindent
    \jpkana{【あなた】とんこつラーメン}{[Anata] Tonkotsu raamen} 
    \jpkana{お}{o}\jpword{ねが}{願}{nega}\jpkana{いします。}{ishimasu.}

    \vspace{1.5em}
    \noindent{\Large \color{articolor} \textbf{Arti:} [Kamu] Tolong Ramen Tonkotsu.}
\end{convobox}

\begin{convobox}{59}
    \noindent
    \jpkana{【店員】}{[Tenin] }\jpword{めん}{麺}{Men} \jpkana{の}{no} 
    \jpkana{かたさは}{katasa wa} 
    \jpkana{どうなさいますか？}{dou nasaimasu ka?}

    \vspace{1.5em}
    \noindent{\Large \color{articolor} \textbf{Arti:} [Pelayan] Tingkat kekerasan mi-nya bagaimana?}
\end{convobox}

\begin{convobox}{60}
    \noindent
    \jpkana{【あなた】}{[Anata]} \jpword{ふつう}{普通}{Futsuu} \jpkana{で}{de} 
    \jpkana{お}{o}\jpword{ねが}{願}{nega}\jpkana{いします。}{ishimasu.}

    \vspace{1.5em}
    \noindent{\Large \color{articolor} \textbf{Arti:} [Kamu] Tolong yang biasa (futsuu).}
\end{convobox}

\begin{convobox}{61}
    \noindent
    \jpkana{【店員】かしこまりました。}{[Tenin] Kashikomarimashita.} 
    \jpkana{トッピングは}{Toppingu wa} 
    \jpkana{いかがですか？}{ikaga desu ka?}

    \vspace{1.5em}
    \noindent{\Large \color{articolor} \textbf{Arti:} [Pelayan] Baik. Bagaimana dengan topping?}
\end{convobox}

\begin{convobox}{62}
    \noindent
    \jpkana{【あなた】}{[Anata]} \jpword{あじたま}{味玉}{Ajitama} \jpkana{と}{to} 
    \jpkana{ねぎ}{negi} \jpkana{トッピング}{toppingu} 
    \jpkana{お}{o}\jpword{ねが}{願}{nega}\jpkana{いします。}{ishimasu.}

    \vspace{1.5em}
    \noindent{\Large \color{articolor} \textbf{Arti:} [Kamu] Tolong topping Ajitama (telur) dan Negi (daun bawang).}
\end{convobox}

\begin{convobox}{63}
    \noindent
    \jpkana{【店員】かしこまりました。}{[Tenin] Kashikomarimashita.}

    \vspace{1.5em}
    \noindent{\Large \color{articolor} \textbf{Arti:} [Pelayan] Baik, saya mengerti pesanan Anda.}
\end{convobox}

% --- SUMMARY BOX ---
\vspace{2em}
\begin{tcolorbox}[
    colback=blue!5!white,
    colframe=blue!80!black,
    boxrule=3pt,
    arc=5mm,
    title={\huge\bfseries \sffamily 🌟 Rangkuman Misi Ramen 🌟},
    fonttitle=\color{white},
    attach boxed title to top center={yshift=-4mm},
    boxed title style={colback=blue!80!black, rounded corners},
    left=5mm, right=5mm, top=7mm, bottom=5mm
]
\Large \textbf{Omedetou! (Selamat!) Kamu siap makan ramen sepuasnya di Jepang!}
\begin{itemize}
    \item \textbf{Pesan Cepat:} Tunjuk gambar dan bilang \textit{Kore, onegaishimasu!}
    \item \textbf{Menghitung Orang:} Hitori (1), Futari (2), San-nin (3), Yo-nin (4).
    \item \textbf{Menghitung Porsi:} Hitotsu (1), Futatsu (2).
    \item \textbf{Custom Ramen:} \textit{Katame} (Keras/Kenyal), \textit{Futsuu} (Biasa), \textit{Yawarakame} (Lembek/Lunak).
    \item \textbf{Mesin Tiket (Kenbaiki):} Masukkan uang DULU, baru tekan tombol!
\end{itemize}
Jangan lupa berlatih pesan dengan orang tuamu di rumah ya! \textit{Oishii raamen o tabete kudasai ne!} (Silakan makan ramen yang enak ya!).
\end{tcolorbox}
\chapter{Ordering at a Cafe in Japanese}

% ================= PAIR 1 =================
\begin{convobox}{1}
    \noindent
    \jpkana{みなさん}{Minasan} \jpkana{こんにちは、}{konnichiwa,} 
    \jpkana{たなかです。}{Tanaka desu.}

    \vspace{1.5em}
    \noindent{\Large \color{articolor} \textbf{Arti:} Halo semuanya, saya Tanaka.}
\end{convobox}

\begin{convobox}{2}
    \noindent
    \jpword{きょう}{今日}{Kyou} \jpkana{は、}{wa,} 
    \jpkana{カフェ}{kafe} \jpkana{で}{de} 
    \jpword{ちゅうもん}{注文}{chuumon} \jpkana{するときの}{suru toki no} 
    \jpword{にほんご}{日本語}{Nihongo} \jpkana{を}{o} 
    \jpword{いっしょ}{一緒}{issho} \jpkana{に}{ni} 
    \jpword{れんしゅう}{練習}{renshuu} \jpkana{しましょう。}{shimashou.}

    \vspace{1.5em}
    \noindent{\Large \color{articolor} \textbf{Arti:} Hari ini, mari kita berlatih bahasa Jepang untuk memesan di kafe bersama-sama.}
\end{convobox}

\begin{lessonbox}{✨ Selamat Datang di Kafe Jepang! (Tips Keluarga)}
    Halo Ayah, Bunda, dan Adik-adik! Kafe di Jepang sangat nyaman dan ramah anak. Selain kopi, biasanya selalu ada menu untuk si kecil seperti \textbf{オレンジジュース (Orenji Juusu - Jus Jeruk)} atau \textbf{ミルク (Miruku - Susu sapi)}. Jangan ragu untuk mampir ke kafe saat kalian butuh istirahat sejenak setelah lelah berjalan-jalan!
\end{lessonbox}

% ================= PAIR 2 =================
\begin{convobox}{3}
    \noindent
    \jpkana{カフェ}{Kafe} \jpkana{で}{de} 
    \jpword{ちゅうもん}{注文}{chuumon} \jpkana{するのって、}{suru no tte,} 
    \jpkana{ちょっと}{chotto} 
    \jpword{きんちょう}{緊張}{kinchou} \jpkana{しませんか？}{shimasen ka?}

    \vspace{1.5em}
    \noindent{\Large \color{articolor} \textbf{Arti:} Memesan di kafe itu, rasanya sedikit bikin gugup kan?}
\end{convobox}

\begin{convobox}{4}
    \noindent
    \jpword{じぶん}{自分}{Jibun} \jpkana{の}{no} 
    \jpword{ちゅうもん}{注文}{chuumon} \jpkana{したい}{shitai} \jpword{もの}{物}{mono} \jpkana{が、}{ga,} 
    \jpkana{うまく}{umaku} \jpword{つた}{伝}{tsuta}\jpkana{わらなかった}{waranakatta} 
    \jpword{けいけん}{経験}{keiken} \jpkana{のある}{no aru} 
    \jpword{かた}{方}{kata} \jpkana{もいるかもしれません。}{mo iru kamo shiremasen.}

    \vspace{1.5em}
    \noindent{\Large \color{articolor} \textbf{Arti:} Mungkin ada juga dari kalian yang punya pengalaman (di mana) pesanan yang kalian inginkan tidak tersampaikan dengan baik.}
\end{convobox}

\begin{lessonbox}{💡 Jangan Gugup, Gunakan Bahasa Tubuh!}
    Kalau kamu takut pesananmu salah, jangan panik! Hampir semua kafe di Jepang punya menu bergambar atau etalase kaca berisi kue. Kamu cukup menunjuk gambar atau kuenya dan bilang: \textbf{「これ、お願いします」 (Kore, onegaishimasu - Tolong yang INI)}. Kasir pasti akan langsung mengerti!
\end{lessonbox}

% ================= PAIR 3 =================
\begin{convobox}{5}
    \noindent
    \jpword{てんいん}{店員}{Tenin}\jpkana{さんは}{san wa} 
    \jpword{はやぐち}{早口}{hayaguchi} \jpkana{だし、}{dashi,} 
    \jpword{けいご}{敬語}{keigo} \jpkana{を}{o} 
    \jpword{つか}{使}{tsuka}\jpkana{うので、}{u node,}

    \vspace{1.5em}
    \noindent{\Large \color{articolor} \textbf{Arti:} Kasir (pelayan toko) itu bicaranya cepat, dan karena mereka menggunakan bahasa super sopan (Keigo),}
\end{convobox}

\begin{convobox}{6}
    \noindent
    \jpword{しつもん}{質問}{Shitsumon} \jpkana{を}{o} 
    \jpword{き}{聞}{ki}\jpkana{き}{ki}\jpword{と}{取}{to}\jpkana{るのが}{ru no ga} 
    \jpword{むずか}{難}{muzuka}\jpkana{しいと}{shii to} 
    \jpword{おも}{思}{omo}\jpkana{います。}{imasu.}

    \vspace{1.5em}
    \noindent{\Large \color{articolor} \textbf{Arti:} aku pikir akan sulit untuk menangkap (mendengar) pertanyaan mereka.}
\end{convobox}

\begin{lessonbox}{🛡️ Tameng Anti-Bingung: Dengarkan Kata Kuncinya!}
    Bahasa sopan kasir Jepang (\textit{Keigo}) memang sangat panjang dan cepat. Rahasianya: \textbf{Jangan coba menerjemahkan setiap kata!} Dengarkan saja kata kuncinya. Kalau terdengar kata \textbf{"Hotto"} atau \textbf{"Aisu"}, berarti mereka tanya suhu minuman. Fokus pada kata kunci saja ya!
\end{lessonbox}

% ================= PAIR 4 =================
\begin{convobox}{7}
    \noindent
    \jpkana{そこで}{Soko de} 
    \jpword{きょう}{今日}{kyou} \jpkana{は、}{wa,} 
    \jpkana{カフェ}{kafe} \jpkana{で}{de} \jpkana{よく}{yoku} 
    \jpword{き}{聞}{ki}\jpkana{かれる}{kareru} 
    \jpword{しつもん}{質問}{shitsumon} \jpkana{と}{to} 
    \jpkana{その}{sono} \jpword{こた}{答}{kota}\jpkana{え}{e}\jpword{かた}{方}{kata} \jpkana{を}{o} 
    \jpword{ようい}{用意}{youi} \jpkana{しました。}{shimashita.}

    \vspace{1.5em}
    \noindent{\Large \color{articolor} \textbf{Arti:} Oleh karena itu hari ini, aku menyiapkan pertanyaan yang sering ditanyakan di kafe beserta cara menjawabnya.}
\end{convobox}

\begin{convobox}{8}
    \noindent
    \jpkana{この}{Kono} \jpword{どうが}{動画}{douga} \jpkana{を}{o} 
    \jpword{み}{観}{mi}\jpkana{て}{te} 
    \jpword{れんしゅう}{練習}{renshuu} \jpkana{すれば、}{sureba,} 
    \jpword{じしん}{自信}{jishin} \jpkana{を}{o} 
    \jpword{も}{持}{mo}\jpkana{って}{tte} 
    \jpword{ちゅうもん}{注文}{chuumon} \jpkana{できるように}{dekiru you ni} \jpkana{なります。}{narimasu.}

    \vspace{1.5em}
    \noindent{\Large \color{articolor} \textbf{Arti:} Jika kamu menonton video ini dan berlatih, kamu akan bisa memesan dengan penuh percaya diri.}
\end{convobox}

\begin{lessonbox}{🌟 Percaya Diri Adalah Kunci Utama!}
    Orang Jepang sangat menghargai turis yang mencoba bicara bahasa Jepang, betapapun terbata-batanya. Ajarkan anak-anak untuk selalu tersenyum dan mengucapkan \textbf{「ありがとう」 (Arigatou - Terima kasih)} setelah menerima struk atau minuman. Itu saja sudah membuat kasir sangat senang!
\end{lessonbox}

% ================= PAIR 5 =================
\begin{convobox}{9}
    \noindent
    \jpkana{それでは}{Sore dewa} 
    \jpword{さっそく}{早速}{sassoku} 
    \jpword{い}{行}{i}\jpkana{ってみましょう！}{tte mimashou!}

    \vspace{1.5em}
    \noindent{\Large \color{articolor} \textbf{Arti:} Kalau begitu, ayo langsung kita mulai!}
\end{convobox}

\vspace{1em}
\begin{tcolorbox}[colback=boxframe, colframe=boxframe, rounded corners, boxrule=0pt, arc=3mm, left=2mm, top=2mm, bottom=2mm]
    \Large \bfseries \color{white} ☕ 1. Conversation (Percakapan Dasar)
\end{tcolorbox}
\vspace{1em}

\begin{convobox}{10}
    \noindent
    \jpword{まず}{先ず}{Mazu} \jpword{いちど}{一度}{ichido} 
    \jpkana{カフェ}{kafe} \jpkana{での}{de no} 
    \jpword{かいわ}{会話}{kaiwa} \jpkana{を}{o} 
    \jpword{しぜん}{自然}{shizen} \jpkana{な}{na} 
    \jpkana{スピードで}{supiido de} 
    \jpword{き}{聞}{ki}\jpkana{いてみましょう。}{ite mimashou.}

    \vspace{1.5em}
    \noindent{\Large \color{articolor} \textbf{Arti:} Pertama-tama, mari dengarkan sekali percakapan di kafe dengan kecepatan natural.}
\end{convobox}

\begin{lessonbox}{🤫 Aturan Emas Kafe Jepang: Amankan Kursi Dulu!}
    Sebelum kita memesan di kasir, ada aturan penting di kafe Jepang (seperti Starbucks atau Tully's): \textbf{Cari dan amankan kursi/meja kalian terlebih dahulu!} Kamu bisa meletakkan syal, payung, atau saputangan di atas meja. Setelah meja aman, barulah seluruh keluarga pergi ke kasir untuk memesan.
\end{lessonbox}

% ================= PAIR 6 =================
\begin{convobox}{11}
    \noindent
    \jpkana{いらっしゃいませ、}{Irasshaimase,} 
    \jpword{てんない}{店内}{ten-nai} 
    \jpword{りよう}{利用}{riyou} \jpkana{ですか？}{desu ka?}

    \vspace{1.5em}
    \noindent{\Large \color{articolor} \textbf{Arti:} Selamat datang, apakah dimakan di dalam (Dine-in)?}
\end{convobox}

\begin{convobox}{12}
    \noindent
    \jpkana{いえ、}{Ie,} 
    \jpword{も}{持}{mo}\jpkana{ち}{chi}\jpword{かえ}{帰}{kae}\jpkana{りで}{ri de} 
    \jpkana{お}{o}\jpword{ねが}{願}{nega}\jpkana{いします。}{ishimasu.}

    \vspace{1.5em}
    \noindent{\Large \color{articolor} \textbf{Arti:} Tidak, tolong dibungkus (Take-away).}
\end{convobox}

\begin{lessonbox}{💰 Kenapa Pertanyaan Ini Sangat Penting?}
    Kasir wajib bertanya ini karena \textbf{Pajak di Jepang berbeda!} 
    \begin{itemize}
        \item Bawa pulang (\textit{Mochikaeri}) pajaknya \textbf{8\%}.
        \item Makan di tempat (\textit{Ten-nai}) pajaknya \textbf{10\%}.
    \end{itemize}
    Jadi, pastikan kamu menjawab dengan benar ya agar kasir bisa menghitung total harganya dengan tepat!
\end{lessonbox}

% ================= PAIR 7 =================
\begin{convobox}{13}
    \noindent
    \jpkana{かしこまりました。}{Kashikomarimashita.}

    \vspace{1.5em}
    \noindent{\Large \color{articolor} \textbf{Arti:} Baik, saya mengerti (Sangat Sopan).}
\end{convobox}

\begin{convobox}{14}
    \noindent
    \jpkana{コーヒー}{Koohii} \jpkana{ひとつ}{hitotsu} 
    \jpkana{お}{o}\jpword{ねが}{願}{nega}\jpkana{いします。}{ishimasu.}

    \vspace{1.5em}
    \noindent{\Large \color{articolor} \textbf{Arti:} Tolong pesan satu kopi.}
\end{convobox}

\begin{lessonbox}{🎀 Kata Sopan Andalan Pelayan: Kashikomarimashita}
    Kata \textbf{かしこまりました (Kashikomarimashita)} adalah bahasa yang sangaaaat sopan dari "Saya mengerti" (\textit{Wakarimashita}). Kamu akan mendengar kata ini di hotel, restoran, dan kafe! Sebagai tamu, kamu tidak perlu memakai kata ini, cukup mendengarkan saja.
\end{lessonbox}

% ================= PAIR 8 =================
\begin{convobox}{15}
    \noindent
    \jpkana{ホットか}{Hotto ka} \jpkana{アイス、}{aisu,} 
    \jpkana{どちらに}{dochira ni} \jpkana{なさいますか？}{nasaimasu ka?}

    \vspace{1.5em}
    \noindent{\Large \color{articolor} \textbf{Arti:} Hot (panas) atau Ice (dingin), Anda mau pilih yang mana?}
\end{convobox}

\begin{convobox}{16}
    \noindent
    \jpkana{ホットで}{Hotto de} 
    \jpkana{お}{o}\jpword{ねが}{願}{nega}\jpkana{いします。}{ishimasu.}

    \vspace{1.5em}
    \noindent{\Large \color{articolor} \textbf{Arti:} Tolong yang panas (Hot).}
\end{convobox}

\begin{lessonbox}{🧊 Bahasa Inggris ala Jepang (Japanglish)}
    Hati-hati saat mengucapkan kata bahasa Inggris di Jepang! 
    \begin{itemize}
        \item Jangan bilang "Hot", bilangnya \textbf{ホット (Hotto)}.
        \item Jangan bilang "Ice", bilangnya \textbf{アイス (Aisu)}.
    \end{itemize}
    Aisu juga bisa berarti "Es Krim" di situasi tertentu. Pastikan lafalnya seperti bahasa Jepang ya!
\end{lessonbox}

% ================= PAIR 9 =================
\begin{convobox}{17}
    \noindent
    \jpkana{お}{O}\jpword{さとう}{砂糖}{satou} \jpkana{・}{・} \jpkana{ミルクは}{miruku wa} 
    \jpword{りよう}{利用}{riyou} \jpkana{ですか？}{desu ka?}

    \vspace{1.5em}
    \noindent{\Large \color{articolor} \textbf{Arti:} Apakah menggunakan gula dan susu?}
\end{convobox}

\begin{convobox}{18}
    \noindent
    \jpkana{ミルクだけ}{Miruku dake} 
    \jpkana{お}{o}\jpword{ねが}{願}{nega}\jpkana{いします。}{ishimasu.}

    \vspace{1.5em}
    \noindent{\Large \color{articolor} \textbf{Arti:} Tolong susu saja.}
\end{convobox}

\begin{lessonbox}{☕ Self-Service Area di Kafe}
    Terkadang, kasir tidak menanyakan ini karena ada \textbf{Area Self-Service} (Meja khusus di mana kamu bisa mengambil gula, sedotan, tisu, dan susu cair kecil sendiri). Jika kamu melihat meja penuh toples kecil di dekat kasir, kamu bisa meracik gulamu sendiri di sana!
\end{lessonbox}

% ================= PAIR 10 =================
\begin{convobox}{19}
    \noindent
    \jpkana{お}{O}\jpword{かいけい}{会計}{kaikei}\jpkana{、}{,} 
    \jpword{ごひゃくえん}{500円}{gohyaku-en} 
    \jpkana{でございます。}{de gozaimasu.}

    \vspace{1.5em}
    \noindent{\Large \color{articolor} \textbf{Arti:} Total pembayarannya, 500 yen.}
\end{convobox}

\begin{convobox}{20}
    \noindent
    \jpword{すこ}{少}{suko}\jpkana{し}{shi} 
    \jpword{はや}{速}{haya}\jpkana{かったですか？}{katta desu ka?}

    \vspace{1.5em}
    \noindent{\Large \color{articolor} \textbf{Arti:} Apakah tadi sedikit terlalu cepat?}
\end{convobox}

\begin{lessonbox}{💳 Siapkan Alat Bayar Andalan Keluarga}
    Saat liburan keluarga, membayar dengan uang receh akan memakan waktu. Paling praktis adalah menggunakan kartu IC Card seperti \textbf{Suica} atau \textbf{Pasmo}! Tempelkan kartunya di mesin, berbunyi "Pi!", dan selesai! Sangat praktis untuk beli minuman dan kereta.
\end{lessonbox}

% ================= PAIR 11 =================
\begin{convobox}{21}
    \noindent
    \jpword{いま}{今}{Ima} \jpkana{から}{kara} 
    \jpword{ひと}{一}{hito}\jpkana{つずつ}{tsu zutsu} 
    \jpword{せつめい}{説明}{setsumei} \jpkana{します。}{shimasu.}

    \vspace{1.5em}
    \noindent{\Large \color{articolor} \textbf{Arti:} Mulai sekarang, saya akan menjelaskannya satu per satu.}
\end{convobox}

\vspace{1em}
\begin{tcolorbox}[colback=boxframe, colframe=boxframe, rounded corners, boxrule=0pt, arc=3mm, left=2mm, top=2mm, bottom=2mm]
    \Large \bfseries \color{white} 🍰 2. Phrase 1 (Makan di Tempat / Bawa Pulang)
\end{tcolorbox}
\vspace{1em}

\begin{convobox}{22}
    \noindent
    \jpkana{「いらっしゃいませ、}{"Irasshaimase,} 
    \jpword{てんない}{店内}{ten-nai} 
    \jpword{りよう}{利用}{riyou} \jpkana{ですか？」}{desu ka?"}

    \vspace{1.5em}
    \noindent{\Large \color{articolor} \textbf{Arti:} "Selamat datang, apakah makan/minum di dalam?"}
\end{convobox}

\begin{lessonbox}{🏫 Ten-nai = Di dalam toko}
    Kata \textbf{店内 (Ten-nai)} terdiri dari kanji 'Mise' (Toko) dan 'Uchi' (Dalam). Jadi ini adalah istilah baku untuk menyebut "Makan di tempat" (Dine-in).
\end{lessonbox}

% ================= PAIR 12 =================
\begin{convobox}{23}
    \noindent
    \jpkana{カフェ}{Kafe} \jpkana{で}{de} 
    \jpword{ちゅうもん}{注文}{chuumon} \jpkana{するとき、}{suru toki,} 
    \jpword{いちばん}{一番}{ichiban} \jpword{さいしょ}{最初}{saisho} \jpkana{に}{ni} 

    \vspace{1.5em}
    \noindent{\Large \color{articolor} \textbf{Arti:} Saat memesan di kafe, hal yang paling pertama}
\end{convobox}

\begin{convobox}{24}
    \noindent
    \jpword{き}{聞}{ki}\jpkana{かれる}{kareru} \jpkana{フレーズです。}{fureezu desu.} 
    \jpkana{「}{"}\jpword{てんない}{店内}{Ten-nai} 
    \jpword{りよう}{利用}{riyou} \jpkana{ですか？」}{desu ka?"}

    \vspace{1.5em}
    \noindent{\Large \color{articolor} \textbf{Arti:} didengar (ditanyakan) adalah frasa ini. "Apakah menggunakan area di dalam toko?"}
\end{convobox}

\begin{lessonbox}{🚫 Jangan Makan Sambil Jalan!}
    Jika keluarga memilih "Bawa Pulang" (Mochikaeri), \textbf{Jangan makan/minum sambil berjalan} ya! Di Jepang, makan sambil berjalan di trotoar dianggap kurang sopan. Carilah bangku taman, atau menyingkir ke sisi jalan untuk meminum jus/kopi kalian!
\end{lessonbox}

% ================= PAIR 13 =================
\begin{convobox}{25}
    \noindent
    \jpkana{もしくは}{Moshiku wa}

    \vspace{1.5em}
    \noindent{\Large \color{articolor} \textbf{Arti:} Atau (bisa juga)}
\end{convobox}

\begin{convobox}{26}
    \noindent
    \jpkana{「お}{"O}\jpword{も}{持}{mo}\jpkana{ち}{chi}\jpword{かえ}{帰}{kae}\jpkana{りですか？」}{ri desu ka?"}

    \vspace{1.5em}
    \noindent{\Large \color{articolor} \textbf{Arti:} "Apakah dibungkus bawa pulang?"}
\end{convobox}

\begin{lessonbox}{🛍️ "O-mochi-kaeri" = Bawa pulang}
    Kata ini asalnya dari kata "Motsu" (Membawa) dan "Kaeru" (Pulang). Ingat baik-baik kata \textbf{持ち帰り (Mochikaeri)} karena kamu akan melihat tulisan ini di kafe, restoran sushi, hingga kedai burger!
\end{lessonbox}

% ================= PAIR 14 =================
\begin{convobox}{27}
    \noindent
    \jpword{こた}{答}{Kota}\jpkana{えるときは、}{eru toki wa,} 
    \jpkana{このように}{kono you ni} \jpword{い}{言}{i}\jpkana{います。}{imasu.}

    \vspace{1.5em}
    \noindent{\Large \color{articolor} \textbf{Arti:} Saat menjawab, ucapkanlah seperti ini.}
\end{convobox}

\begin{convobox}{28}
    \noindent
    \jpkana{「}{"}\jpword{てんない}{店内}{Ten-nai} \jpkana{で}{de} 
    \jpkana{お}{o}\jpword{ねが}{願}{nega}\jpkana{いします。」}{ishimasu."}

    \vspace{1.5em}
    \noindent{\Large \color{articolor} \textbf{Arti:} "Tolong di dalam sini (Makan di tempat)."}
\end{convobox}

\begin{lessonbox}{🪄 Kekuatan Partikel "De"}
    Jangan lupa tambahkan partikel \textbf{で (De)} setelah pilihanmu! \textit{Ten-nai DE} (DI dalam), atau \textit{Mochikaeri DE} (DENGAN cara dibawa pulang). Lalu tutup dengan jurus \textit{Onegaishimasu} agar terdengar manis dan sopan!
\end{lessonbox}

% ================= PAIR 15 =================
\begin{convobox}{29}
    \noindent
    \jpkana{「お}{"O}\jpword{も}{持}{mo}\jpkana{ち}{chi}\jpword{かえ}{帰}{kae}\jpkana{りで}{ri de} 
    \jpkana{お}{o}\jpword{ねが}{願}{nega}\jpkana{いします。」}{ishimasu."}

    \vspace{1.5em}
    \noindent{\Large \color{articolor} \textbf{Arti:} "Tolong dibungkus (Bawa pulang)."}
\end{convobox}

\vspace{1em}
\begin{tcolorbox}[colback=boxframe, colframe=boxframe, rounded corners, boxrule=0pt, arc=3mm, left=2mm, top=2mm, bottom=2mm]
    \Large \bfseries \color{white} 🍩 3. Phrase 2 (Sebut Pesanan \& Jumlah)
\end{tcolorbox}
\vspace{1em}

\begin{convobox}{30}
    \noindent
    \jpkana{「コーヒー}{"Koohii} \jpkana{ひとつ}{hitotsu} 
    \jpkana{お}{o}\jpword{ねが}{願}{nega}\jpkana{いします。」}{ishimasu."}

    \vspace{1.5em}
    \noindent{\Large \color{articolor} \textbf{Arti:} "Tolong 1 buah kopi."}
\end{convobox}

\begin{lessonbox}{🥤 Cara Memesan Jus untuk Anak}
    Jika anak-anak ingin memesan sendiri, ajarkan mereka menyebut minumannya lalu angka \textit{hitotsu}! Contoh:
    "Ringo juusu (Jus apel) hitotsu, onegaishimasu!" Pasti kasirnya akan merasa sangat gemas!
\end{lessonbox}

% ================= PAIR 16 =================
\begin{convobox}{31}
    \noindent
    \jpword{ちゅうもん}{注文}{Chuumon} \jpkana{するものの}{suru mono no} 
    \jpword{なまえ}{名前}{namae} \jpkana{と}{to} 
    \jpword{かず}{数}{kazu} \jpkana{を}{o} 
    \jpword{つた}{伝}{tsuta}\jpkana{える}{eru} 
    \jpkana{フレーズです。}{fureezu desu.}

    \vspace{1.5em}
    \noindent{\Large \color{articolor} \textbf{Arti:} Ini adalah frasa untuk menyampaikan nama dan jumlah barang yang dipesan.}
\end{convobox}

\begin{convobox}{32}
    \noindent
    \jpword{かず}{数}{Kazu} \jpkana{は、}{wa,} 
    \jpkana{このように}{kono you ni} 
    \jpword{い}{言}{i}\jpkana{います。}{imasu.}

    \vspace{1.5em}
    \noindent{\Large \color{articolor} \textbf{Arti:} Untuk jumlahnya, diucapkan seperti ini.}
\end{convobox}

\begin{lessonbox}{📝 Urutan yang Benar}
    Berbeda dengan bahasa Indonesia "Satu Kopi", di Jepang urutannya dibalik: \textbf{BARANGNYA DULU, baru JUMLAHNYA}. Jadi: Kopi 1 (Koohii hitotsu), Kue 2 (Keeki futatsu). Ingat-ingat rumus ini ya!
\end{lessonbox}

% ================= PAIR 17 =================
\begin{convobox}{33}
    \noindent
    \jpkana{ひとつ}{Hitotsu} (1) / \jpkana{ふたつ}{Futatsu} (2) / \jpkana{みっつ}{Mittsu} (3) / \jpkana{よっつ}{Yottsu} (4) / \jpkana{いつつ}{Itsutsu} (5) / \jpkana{むっつ}{Muttsu} (6)
\end{convobox}

\begin{convobox}{34}
    \noindent
    \jpkana{「ひとつ」と}{"Hitotsu" to} 
    \jpkana{「ふたつ」は}{"futatsu" wa} 
    \jpword{はつおん}{発音}{hatsuon} \jpkana{が}{ga} 
    \jpkana{よく}{yoku} \jpword{に}{似}{ni}\jpkana{ているので、}{te iru node,}

    \vspace{1.5em}
    \noindent{\Large \color{articolor} \textbf{Arti:} Karena pengucapan "Hitotsu (1)" dan "Futatsu (2)" itu sangat mirip,}
\end{convobox}

\begin{lessonbox}{Hitotsu dan Futatsu}
    "Hitotsu" dan "Futatsu" terdengar hampir sama, terutama jika kafe sedang ramai dan berisik. Kadang kasir salah dengar kamu pesan 1 jadi pesan 2! Untuk menghindarinya, \textbf{Selalu angkat jarimu} saat menyebutkan angkanya ya!
\end{lessonbox}

% ================= PAIR 18 =================
\begin{convobox}{35}
    \noindent
    \jpword{ゆび}{指}{Yubi} \jpkana{も}{mo} 
    \jpword{だ}{出}{da}\jpkana{すと}{su to} 
    \jpkana{いいでしょう。}{ii deshou.}

    \vspace{1.5em}
    \noindent{\Large \color{articolor} \textbf{Arti:} ada baiknya jika kamu mengeluarkan (menunjukkan) jari juga.}
\end{convobox}

\begin{convobox}{36}
    \noindent
    \jpword{ちゅうもん}{注文}{Chuumon} \jpkana{するものが}{suru mono ga} 
    \jpkana{いくつか}{ikutsu ka} \jpkana{ある}{aru} 
    \jpword{とき}{時}{toki}\jpkana{は、}{wa,} 
    \jpkana{このように}{kono you ni} \jpword{い}{言}{i}\jpkana{います。}{imasu.}

    \vspace{1.5em}
    \noindent{\Large \color{articolor} \textbf{Arti:} Saat barang yang dipesan ada beberapa, ucapkanlah seperti ini.}
\end{convobox}

\begin{lessonbox}{👨‍👩‍👧‍👦 Pesan untuk Seluruh Keluarga}
    Biasanya Ayah atau Bunda yang akan memesan untuk seluruh anggota keluarga kan? Nah, siap-siap merangkai kata dengan partikel penghubung di bawah ini ya!
\end{lessonbox}

% ================= PAIR 19 =================
\begin{convobox}{37}
    \noindent
    \jpkana{「コーヒー}{"Koohii} \jpkana{ひとつ}{hitotsu} \jpkana{と、}{to,} 
    \jpkana{チーズケーキ}{chiizu keeki} \jpkana{ひとつ}{hitotsu} 
    \jpkana{お}{o}\jpword{ねが}{願}{nega}\jpkana{いします。」}{ishimasu."}

    \vspace{1.5em}
    \noindent{\Large \color{articolor} \textbf{Arti:} "Tolong 1 kopi DAN 1 cheesecake."}
\end{convobox}

\begin{convobox}{38}
    \noindent
    \jpkana{「カフェオレ}{"Kafe ore} \jpkana{ふたつ}{futatsu} \jpkana{と、}{to,} 
    \jpkana{サンドイッチ}{sandoitchi} \jpkana{ひとつ}{hitotsu} 
    \jpkana{お}{o}\jpword{ねが}{願}{nega}\jpkana{いします。」}{ishimasu."}

    \vspace{1.5em}
    \noindent{\Large \color{articolor} \textbf{Arti:} "Tolong 2 cafe au lait (kopi susu) DAN 1 sandwich."}
\end{convobox}

\begin{lessonbox}{🔗 Lem Perekat Kalimat: Partikel "To"}
    Partikel \textbf{と (to)} fungsinya sama seperti kata "Dan". Kamu bisa menyambung sebanyak apapun pesananmu! \textit{A to, B to, C to, D, onegaishimasu!}. Kasir akan dengan sabar mencatat semuanya!
\end{lessonbox}

% ================= PAIR 20 =================
\vspace{1em}
\begin{tcolorbox}[colback=boxframe, colframe=boxframe, rounded corners, boxrule=0pt, arc=3mm, left=2mm, top=2mm, bottom=2mm]
    \Large \bfseries \color{white} 🥤 4. Phrase 3 (Hot or Ice?)
\end{tcolorbox}
\vspace{1em}

\begin{convobox}{39}
    \noindent
    \jpkana{「ホットか}{"Hotto ka} \jpkana{アイス、}{aisu,} 
    \jpkana{どちらに}{dochira ni} \jpkana{なさいますか？」}{nasaimasu ka?"}

    \vspace{1.5em}
    \noindent{\Large \color{articolor} \textbf{Arti:} "Hot atau Ice, Anda ingin yang mana?"}
\end{convobox}

\begin{convobox}{40}
    \noindent
    \jpkana{コーヒーや}{Koohii ya} 
    \jpword{こうちゃ}{紅茶}{koucha} \jpkana{のように、}{no you ni,} 
    \jpkana{ホットか}{hotto ka} \jpkana{アイスを}{aisu o} 
    \jpword{えら}{選}{era}\jpkana{べる}{beru} 

    \vspace{1.5em}
    \noindent{\Large \color{articolor} \textbf{Arti:} Seperti kopi atau teh hitam, (minuman) yang bisa dipilih Hot atau Ice-nya}
\end{convobox}

\begin{lessonbox}{🧊 Es di Jepang Itu BANYAK Banget!}
    Sedikit tips: Kalau kamu memesan minuman dingin (\textbf{Aisu}), kafe Jepang biasanya mengisi gelasmu dengan ES BATU yang saaaaaangat penuh! Kalau kamu atau si kecil tidak mau terlalu dingin, bilang saja: \textbf{「氷 少なめで」 (Koori sukuname de - Esnya sedikit saja ya)}.
\end{lessonbox}

% ================= PAIR 21 =================
\begin{convobox}{41}
    \noindent
    \jpword{の}{飲}{no}\jpkana{み}{mi}\jpword{もの}{物}{mono} \jpkana{を}{o} 
    \jpword{ちゅうもん}{注文}{chuumon} \jpkana{した}{shita} 
    \jpword{とき}{時}{toki}\jpkana{は}{wa} 
    \jpkana{このように}{kono you ni} 
    \jpword{き}{聞}{ki}\jpkana{かれます。}{karemasu.}

    \vspace{1.5em}
    \noindent{\Large \color{articolor} \textbf{Arti:} Saat memesan minuman, kamu akan ditanya seperti ini.}
\end{convobox}

\begin{convobox}{42}
    \noindent
    \jpword{こた}{答}{Kota}\jpkana{えるときは、}{eru toki wa,} 
    \jpkana{こう}{kou} \jpword{い}{言}{i}\jpkana{いましょう。}{imashou.}

    \vspace{1.5em}
    \noindent{\Large \color{articolor} \textbf{Arti:} Saat menjawab, mari ucapkan seperti ini.}
\end{convobox}

\begin{lessonbox}{☕ Pilihan Cerdas Musim Dingin}
    Kalau kamu liburan ke Jepang saat musim dingin (Winter), memesan \textbf{Hotto} (Minuman panas) adalah pilihan terbaik! Gelas kopimu juga bisa berfungsi ganda sebagai penghangat tangan (Kairo) dadakan!
\end{lessonbox}

% ================= PAIR 22 =================
\begin{convobox}{43}
    \noindent
    \jpkana{「ホットで}{"Hotto de} \jpkana{お}{o}\jpword{ねが}{願}{nega}\jpkana{いします。」}{ishimasu."}

    \vspace{1.5em}
    \noindent{\Large \color{articolor} \textbf{Arti:} "Tolong yang panas."}
\end{convobox}

\begin{convobox}{44}
    \noindent
    \jpkana{「アイスで}{"Aisu de} \jpkana{お}{o}\jpword{ねが}{願}{nega}\jpkana{いします。」}{ishimasu."}

    \vspace{1.5em}
    \noindent{\Large \color{articolor} \textbf{Arti:} "Tolong yang es (dingin)."}
\end{convobox}

\begin{lessonbox}{⚠️ Alergi Susu? Peringatan Penting!}
    Kalau ada anggota keluarga yang alergi susu sapi (Lactose Intolerant), pastikan untuk memesan kopi hitam (Burakku Koohii), dan jangan memesan "Kafe Rate" atau "Cafe o-lait". Atau bilang: \textbf{ミルクのアレルギーがあります (Miruku no arerugii ga arimasu - Saya ada alergi susu)}.
\end{lessonbox}

% ================= PAIR 23 =================
\vspace{1em}
\begin{tcolorbox}[colback=boxframe, colframe=boxframe, rounded corners, boxrule=0pt, arc=3mm, left=2mm, top=2mm, bottom=2mm]
    \Large \bfseries \color{white} 🍯 5. Phrase 4 (Gula \& Susu)
\end{tcolorbox}
\vspace{1em}

\begin{convobox}{45}
    \noindent
    \jpkana{「お}{"O}\jpword{さとう}{砂糖}{satou} \jpkana{・}{・} \jpkana{ミルクは}{miruku wa} 
    \jpword{りよう}{利用}{riyou} \jpkana{ですか？」}{desu ka?"}

    \vspace{1.5em}
    \noindent{\Large \color{articolor} \textbf{Arti:} "Apakah membutuhkan gula dan susu?"}
\end{convobox}

\begin{convobox}{46}
    \noindent
    \jpword{こた}{答}{Kota}\jpkana{えるときは、}{eru toki wa,} 
    \jpkana{このように}{kono you ni} \jpword{い}{言}{i}\jpkana{います。}{imasu.}

    \vspace{1.5em}
    \noindent{\Large \color{articolor} \textbf{Arti:} Saat menjawab, ucapkanlah seperti ini.}
\end{convobox}

\begin{lessonbox}{🍯 Apa Itu "Susu" untuk Kopi Hitam?}
    Di Jepang, kalau kamu pesan kopi hitam lalu kasir menawarkan "Miruku" (Susu), ini BUKAN susu cair literan lho! Ini adalah krimer cair kecil di dalam cup plastik mungil yang disebut \textbf{コーヒーフレッシュ (Koohii Furesshu - Coffee Fresh)}. Gratis kok!
\end{lessonbox}

% ================= PAIR 24 =================
\begin{convobox}{47}
    \noindent
    \jpkana{「はい、}{"Hai,} \jpkana{お}{o}\jpword{ねが}{願}{nega}\jpkana{いします。」}{ishimasu."}

    \vspace{1.5em}
    \noindent{\Large \color{articolor} \textbf{Arti:} "Ya, tolong (berikan keduanya)."}
\end{convobox}

\begin{convobox}{48}
    \noindent
    \jpkana{「}{"}\jpword{さとう}{砂糖}{Satou} \jpkana{だけ}{dake} 
    \jpkana{お}{o}\jpword{ねが}{願}{nega}\jpkana{いします。」}{ishimasu."}

    \vspace{1.5em}
    \noindent{\Large \color{articolor} \textbf{Arti:} "Tolong gula saja."}
\end{convobox}

\begin{lessonbox}{💎 Kata Ajaib "Dake" (Hanya)}
    \textbf{だけ (Dake)} berarti "Hanya/Saja". Kalau anakmu hanya suka sirup strawberry tanpa krim di atas minumannya, kamu bisa bilang \textit{"Ichigo shiroppu \textbf{dake} onegaishimasu"}. Sangat berguna untuk pilih-pilih makanan!
\end{lessonbox}

% ================= PAIR 25 =================
\begin{convobox}{49}
    \noindent
    \jpkana{「ミルクだけ}{"Miruku dake} \jpkana{お}{o}\jpword{ねが}{願}{nega}\jpkana{いします。」}{ishimasu."}

    \vspace{1.5em}
    \noindent{\Large \color{articolor} \textbf{Arti:} "Tolong susu saja."}
\end{convobox}

\begin{convobox}{50}
    \noindent
    \jpkana{「いえ、}{"Ie,} \jpword{だいじょうぶ}{大丈夫}{daijoubu} \jpkana{です。」}{desu."}

    \vspace{1.5em}
    \noindent{\Large \color{articolor} \textbf{Arti:} "Tidak, tidak usah."}
\end{convobox}

\begin{lessonbox}{♻️ Mengembalikan Nampan Sendiri}
    \textbf{BUDAYA PENTING!} Di kafe Jepang seperti Starbucks, setelah selesai makan, kita \textbf{TIDAK BOLEH} meninggalkan gelas dan nampan di meja! Carilah tempat bernama \textbf{返却口 (Henkyakuguchi - Tempat Pengembalian)}. Buang sampahmu dan letakkan nampan di sana. Ini bisa jadi pelajaran disiplin yang seru buat anak-anak!
\end{lessonbox}

% ================= PAIR 26 =================
\vspace{1em}
\begin{tcolorbox}[colback=boxframe, colframe=boxframe, rounded corners, boxrule=0pt, arc=3mm, left=2mm, top=2mm, bottom=2mm]
    \Large \bfseries \color{white} 🗣️ 6. Role Play 1 (Simulasi Memesan Kopi)
\end{tcolorbox}
\vspace{1em}

\begin{convobox}{51}
    \noindent
    \jpkana{それでは、}{Sore dewa,} 
    \jpword{いま}{今}{ima} \jpkana{から}{kara} 
    \jpkana{ロールプレイを}{rooru purei o} \jpkana{します。}{shimasu.}

    \vspace{1.5em}
    \noindent{\Large \color{articolor} \textbf{Arti:} Kalau begitu, sekarang kita akan mulai bermain peran (Roleplay).}
\end{convobox}

\begin{convobox}{52}
    \noindent
    \jpword{さき}{先}{Sakki} \jpkana{ほどの}{hodo no} 
    \jpkana{カフェ}{kafe} \jpkana{での}{de no} 
    \jpword{かいわ}{会話}{kaiwa} \jpkana{を}{o} 
    \jpword{もういちど}{もう一度}{mou ichido} 
    \jpkana{なが}{流}{naga}\jpkana{します。}{shimasu.}

    \vspace{1.5em}
    \noindent{\Large \color{articolor} \textbf{Arti:} Aku akan memutar kembali percakapan kafe yang tadi.}
\end{convobox}

\begin{lessonbox}{🎭 Waktunya Akting! Bersiaplah!}
    Ayo semuanya tarik napas panjang! Tegakkan badanmu! Bayangkan kamu sedang berdiri di depan kasir kafe Jepang di Tokyo. Berbicaralah dengan lantang dan senyum yang ramah ya!
\end{lessonbox}

% ================= PAIR 27 =================
\begin{convobox}{53}
    \noindent
    \jpword{こんかい}{今回}{Konkai} \jpkana{は、}{wa,} 
    \jpkana{みなさんは}{minasan wa} 
    \jpkana{お}{o}\jpword{きゃく}{客}{kyaku}\jpkana{さんに}{san ni} 
    \jpkana{なったつもりで}{natta tsumori de} 

    \vspace{1.5em}
    \noindent{\Large \color{articolor} \textbf{Arti:} Kali ini, bayangkanlah kalian menjadi pelanggannya, dan}
\end{convobox}

\begin{convobox}{54}
    \noindent
    \jpword{こえ}{声}{koe} \jpkana{に}{ni} \jpword{だ}{出}{da}\jpkana{して}{shite} 
    \jpword{はな}{話}{hana}\jpkana{してみてください。}{shite mite kudasai.}

    \vspace{1.5em}
    \noindent{\Large \color{articolor} \textbf{Arti:} cobalah berbicara dengan mengeluarkan suara.}
\end{convobox}

\begin{lessonbox}{🎙️ Suara Lantang itu Bagus!}
    Kafe di Jepang biasanya memutar musik BGM (Background Music) dan ada suara mesin espresso yang bising. Jadi pastikan volume suaramu jelas ya agar kasir tidak meminta kita mengulang pesanan.
\end{lessonbox}

% ================= PAIR 28 =================
\begin{convobox}{55}
    \noindent
    \jpkana{【店員】いらっしゃいませ、}{[Tenin] Irasshaimase,} 
    \jpword{てんない}{店内}{ten-nai} 
    \jpword{りよう}{利用}{riyou} \jpkana{ですか？}{desu ka?}

    \vspace{1.5em}
    \noindent{\Large \color{articolor} \textbf{Arti:} [Kasir] Selamat datang, apakah dimakan di dalam?}
\end{convobox}

\begin{convobox}{56}
    \noindent
    \jpkana{【あなた】いえ、}{[Anata] Ie,} 
    \jpword{も}{持}{mo}\jpkana{ち}{chi}\jpword{かえ}{帰}{kae}\jpkana{りで}{ri de} 
    \jpkana{お}{o}\jpword{ねが}{願}{nega}\jpkana{いします。}{ishimasu.}

    \vspace{1.5em}
    \noindent{\Large \color{articolor} \textbf{Arti:} [Kamu] Tidak, tolong dibungkus bawa pulang.}
\end{convobox}

\begin{lessonbox}{🚶‍♂️ Mochikaeri: Cara Menghemat Waktu!}
    Kalau keluarga harus segera mengejar jadwal kereta Shinkansen, opsi \textit{Mochikaeri} (Bawa Pulang) adalah pilihan terbaik! Kopimu akan dimasukkan ke dalam kantong kertas (Kami-bukuro) yang aman agar tidak tumpah.
\end{lessonbox}

% ================= PAIR 29 =================
\begin{convobox}{57}
    \noindent
    \jpkana{【店員】かしこまりました。}{[Tenin] Kashikomarimashita.}

    \vspace{1.5em}
    \noindent{\Large \color{articolor} \textbf{Arti:} [Kasir] Baik.}
\end{convobox}

\begin{convobox}{58}
    \noindent
    \jpkana{【あなた】コーヒー}{[Anata] Koohii} \jpkana{ひとつ}{hitotsu} 
    \jpkana{お}{o}\jpword{ねが}{願}{nega}\jpkana{いします。}{ishimasu.}

    \vspace{1.5em}
    \noindent{\Large \color{articolor} \textbf{Arti:} [Kamu] Tolong 1 buah kopi.}
\end{convobox}

\begin{lessonbox}{☝️ Jangan Lupa Angkat 1 Jari!}
    Sambil bilang "Hitotsu", angkat satu jarimu (jari telunjuk)! Bahasa tubuh ini sangat universal dan langsung dipahami oleh kasir Jepang.
\end{lessonbox}

% ================= PAIR 30 =================
\begin{convobox}{59}
    \noindent
    \jpkana{【店員】ホットか}{[Tenin] Hotto ka} \jpkana{アイス、}{aisu,} 
    \jpkana{どちらに}{dochira ni} \jpkana{なさいますか？}{nasaimasu ka?}

    \vspace{1.5em}
    \noindent{\Large \color{articolor} \textbf{Arti:} [Kasir] Hot atau Ice, mau pilih yang mana?}
\end{convobox}

\begin{convobox}{60}
    \noindent
    \jpkana{【あなた】ホットで}{[Anata] Hotto de} 
    \jpkana{お}{o}\jpword{ねが}{願}{nega}\jpkana{いします。}{ishimasu.}

    \vspace{1.5em}
    \noindent{\Large \color{articolor} \textbf{Arti:} [Kamu] Tolong yang Hot (Panas).}
\end{convobox}

\begin{lessonbox}{❄️ Hati-Hati Cangkir Panas!}
    Jika kamu memesan \textit{Hotto}, biasanya cangkir kertasnya akan sangat panas. Pelayan kadang akan membungkusnya dengan \textbf{スリーブ (Suriibu - Sleeve / Pelindung kardus)}. Jangan biarkan anak-anak yang memegangnya ya!
\end{lessonbox}

% ================= PAIR 31 =================
\begin{convobox}{61}
    \noindent
    \jpkana{【店員】お}{[Tenin] O}\jpword{さとう}{砂糖}{satou} \jpkana{・}{・} \jpkana{ミルクは}{miruku wa} 
    \jpword{りよう}{利用}{riyou} \jpkana{ですか？}{desu ka?}

    \vspace{1.5em}
    \noindent{\Large \color{articolor} \textbf{Arti:} [Kasir] Apakah butuh gula dan susu?}
\end{convobox}

\begin{convobox}{62}
    \noindent
    \jpkana{【あなた】ミルクだけ}{[Anata] Miruku dake} 
    \jpkana{お}{o}\jpword{ねが}{願}{nega}\jpkana{いします。}{ishimasu.}

    \vspace{1.5em}
    \noindent{\Large \color{articolor} \textbf{Arti:} [Kamu] Tolong susu saja.}
\end{convobox}

\begin{lessonbox}{🍭 Meminta Sedotan atau Tisu Basah}
    Kalau kamu butuh sesuatu yang lain, kamu bisa minta ke kasir!
    Sedotan = \textbf{Sutoro}. Tisu basah untuk lap tangan = \textbf{Oshibori}. 
    Bilang saja: \textit{"Sutoro kudasai!"} atau \textit{"Oshibori kudasai!"}
\end{lessonbox}

% ================= PAIR 32 =================
\begin{convobox}{63}
    \noindent
    \jpkana{【店員】お}{[Tenin] O}\jpword{かいけい}{会計}{kaikei}\jpkana{、}{,} 
    \jpword{ごひゃくえん}{500円}{gohyaku-en} \jpkana{でございます。}{de gozaimasu.}

    \vspace{1.5em}
    \noindent{\Large \color{articolor} \textbf{Arti:} [Kasir] Totalnya 500 yen.}
\end{convobox}

\begin{convobox}{64}
    \noindent
    \jpkana{いい}{Ii} \jpword{かん}{感}{kan}\jpkana{じです！}{ji desu!}

    \vspace{1.5em}
    \noindent{\Large \color{articolor} \textbf{Arti:} [Tanaka Sensei] Kelihatan bagus / Kamu melakukannya dengan baik!}
\end{convobox}

\begin{lessonbox}{📱 Saat Kasir Menunjuk Layar Kecil}
    Jika kamu tidak mengerti angka dalam bahasa Jepang (seperti Gohyaku-en), jangan khawatir! Mesin kasir pasti punya \textbf{Layar Kecil yang menghadap ke arah pembeli}. Cukup lihat angka di layar itu untuk tahu berapa yang harus dibayar. Canggih, kan?
\end{lessonbox}

% ================= PAIR 33 =================
\vspace{1em}
\begin{tcolorbox}[colback=boxframe, colframe=boxframe, rounded corners, boxrule=0pt, arc=3mm, left=2mm, top=2mm, bottom=2mm]
    \Large \bfseries \color{white} 🗣️ 7. Role Play 2 (Pesanan Keluarga Besar!)
\end{tcolorbox}
\vspace{1em}

\begin{convobox}{65}
    \noindent
    \jpword{さいご}{最後}{Saigo} \jpkana{に、}{ni,} 
    \jpword{もういちど}{もう一度}{mou ichido} 
    \jpkana{ロールプレイを}{rooru purei o} \jpkana{します。}{shimasu.}

    \vspace{1.5em}
    \noindent{\Large \color{articolor} \textbf{Arti:} Terakhir, mari kita lakukan roleplay sekali lagi.}
\end{convobox}

\begin{convobox}{66}
    \noindent
    \jpword{こんど}{今度}{Kondo} \jpkana{は、}{wa,} 
    \jpkana{イラストを}{irasuto o} \jpword{み}{見}{mi}\jpkana{て}{te} 
    \jpword{ちゅうもん}{注文}{chuumon} \jpkana{してください。}{shite kudasai.}

    \vspace{1.5em}
    \noindent{\Large \color{articolor} \textbf{Arti:} Kali ini, silakan memesan dengan melihat ilustrasi gambarnya.}
\end{convobox}

\begin{lessonbox}{📸 Menu Jepang Sangat Membantu Visual!}
    Ilustrasi atau foto di menu Jepang 99\% sama dengan bentuk aslinya! Jadi kalau kamu lihat kue cantik dengan stroberi, yang datang nanti persis seperti itu. Jangan ragu pesan dari gambar!
\end{lessonbox}

% ================= PAIR 34 =================
\begin{convobox}{67}
    \noindent
    \jpword{じっさい}{実際}{Jissai} \jpkana{に}{ni} 
    \jpkana{カフェにいるのを}{kafe ni iru no o} 
    \jpkana{イメージして}{imeeji shite} 
    \jpword{はな}{話}{hana}\jpkana{してみてくださいね。}{shite mite kudasai ne.}

    \vspace{1.5em}
    \noindent{\Large \color{articolor} \textbf{Arti:} Cobalah berbicara sambil mengimajinasikan bahwa kamu benar-benar sedang berada di kafe (sungguhan).}
\end{convobox}

\begin{convobox}{68}
    \noindent
    \jpkana{【店員】いらっしゃいませ、}{[Tenin] Irasshaimase,} 
    \jpword{てんない}{店内}{ten-nai} 
    \jpword{りよう}{利用}{riyou} \jpkana{ですか？}{desu ka?}

    \vspace{1.5em}
    \noindent{\Large \color{articolor} \textbf{Arti:} [Kasir] Selamat datang, dimakan di sini?}
\end{convobox}

\begin{lessonbox}{🎒 Anak-anak Berbaris di Belakang}
    Tips keluarga: Saat memesan di kasir yang sempit, biarkan satu orang tua saja yang maju berbicara ke kasir (membawa catatan pesanan), sedangkan anak-anak dan orang tua satunya bisa menunggu atau mencari kursi terlebih dahulu!
\end{lessonbox}

% ================= PAIR 35 =================
\begin{convobox}{69}
    \noindent
    \jpkana{【あなた】はい、}{[Anata] Hai,} 
    \jpword{てんない}{店内}{ten-nai} \jpkana{で}{de} 
    \jpkana{お}{o}\jpword{ねが}{願}{nega}\jpkana{いします。}{ishimasu.}

    \vspace{1.5em}
    \noindent{\Large \color{articolor} \textbf{Arti:} [Kamu] Ya, tolong dimakan di sini.}
\end{convobox}

\begin{convobox}{70}
    \noindent
    \jpkana{【店員】かしこまりました。}{[Tenin] Kashikomarimashita.}

    \vspace{1.5em}
    \noindent{\Large \color{articolor} \textbf{Arti:} [Kasir] Baik.}
\end{convobox}

\begin{lessonbox}{🍵 Air Putih Gratis (O-hiya)}
    Saat kamu makan di tempat (\textit{Ten-nai}), biasanya akan ada gelas dan teko air putih di sudut ruangan. Air putih gratis ini disebut \textbf{お冷 (O-hiya)}. Silakan ambil sendiri untuk seluruh keluargamu!
\end{lessonbox}

% ================= PAIR 36 =================
\begin{convobox}{71}
    \noindent
    \jpkana{【あなた】コーヒー}{[Anata] Koohii} \jpkana{ひとつ}{hitotsu} \jpkana{と、}{to,} 
    \jpkana{アップルジュース}{appuru juusu} \jpkana{ひとつ}{hitotsu} \jpkana{と}{to} 
    \jpkana{ブルーベリーマフィン}{buruuberii mafin} \jpkana{ふたつ}{futatsu} 
    \jpkana{お}{o}\jpword{ねが}{願}{nega}\jpkana{いします。}{ishimasu.}

    \vspace{1.5em}
    \noindent{\Large \color{articolor} \textbf{Arti:} [Kamu] Tolong 1 kopi, dan 1 jus apel, dan 2 muffin blueberry.}
\end{convobox}

\begin{convobox}{72}
    \noindent
    \jpkana{【店員】コーヒーは}{[Tenin] Koohii wa} 
    \jpkana{ホットか}{hotto ka} \jpkana{アイス、}{aisu,} 
    \jpkana{どちらに}{dochira ni} \jpkana{なさいますか？}{nasaimasu ka?}

    \vspace{1.5em}
    \noindent{\Large \color{articolor} \textbf{Arti:} [Kasir] Untuk kopinya, Hot atau Ice, Anda ingin pilih yang mana?}
\end{convobox}

\begin{lessonbox}{🧁 Memesan Kue yang Menggoda}
    Lihat kan panjangnya pesanan di atas? Kamu menggabungkan Kopi, Jus Apel, dan 2 Muffin sekaligus memakai partikel \textbf{と (To)}. Ini adalah pesanan sempurna untuk Ibu, Anak, dan cemilan bersama!
\end{lessonbox}

% ================= PAIR 37 =================
\begin{convobox}{73}
    \noindent
    \jpkana{【あなた】アイスで}{[Anata] Aisu de} 
    \jpkana{お}{o}\jpword{ねが}{願}{nega}\jpkana{いします。}{ishimasu.}

    \vspace{1.5em}
    \noindent{\Large \color{articolor} \textbf{Arti:} [Kamu] Tolong yang Es.}
\end{convobox}

\begin{convobox}{74}
    \noindent
    \jpkana{【店員】お}{[Tenin] O}\jpword{さとう}{砂糖}{satou} \jpkana{・}{・} \jpkana{ミルクは}{miruku wa} 
    \jpword{りよう}{利用}{riyou} \jpkana{ですか？}{desu ka?}

    \vspace{1.5em}
    \noindent{\Large \color{articolor} \textbf{Arti:} [Kasir] Apakah menggunakan gula dan susu?}
\end{convobox}

\begin{lessonbox}{🧋 Pesan Minuman Kekinian: "Matcha Frapucino"}
    Kalau kamu masuk ke gerai kopi besar, minuman es kocok yang populer bukan cuma disebut 'Aisu'. Mereka sering memakai kata 'Frape' atau nama merek. Anak-anak pasti sangat suka minuman dingin manis ini di musim panas Jepang!
\end{lessonbox}

% ================= PAIR 38 =================
\begin{convobox}{75}
    \noindent
    \jpkana{【あなた】いえ、}{[Anata] Ie,} 
    \jpword{だいじょうぶ}{大丈夫}{daijoubu} \jpkana{です。}{desu.}

    \vspace{1.5em}
    \noindent{\Large \color{articolor} \textbf{Arti:} [Kamu] Tidak, tidak usah.}
\end{convobox}

\begin{convobox}{76}
    \noindent
    \jpkana{【店員】お}{[Tenin] O}\jpword{かいけい}{会計}{kaikei}\jpkana{、}{,} 
    \jpword{にせんえん}{2,000円}{ni-sen-en} 
    \jpkana{でございます。}{de gozaimasu.}

    \vspace{1.5em}
    \noindent{\Large \color{articolor} \textbf{Arti:} [Kasir] Totalnya 2,000 yen.}
\end{convobox}

\begin{lessonbox}{💴 Membayar Tanpa Ribet!}
    Angka 2.000 Yen dibaca \textbf{Ni-sen-en}. Kalau kamu pakai uang kertas pecahan \textbf{5.000 Yen (Go-sen-en)} atau \textbf{10.000 Yen (Ichi-man-en)}, berikan saja ke kasir. Mereka punya mesin penghitung otomatis yang akan mengeluarkan uang kembalianmu dengan sangat akurat! 
\end{lessonbox}

% ================= PAIR 39 =================
\begin{convobox}{77}
    \noindent
    \jpkana{いかがでしたか？}{Ikaga deshita ka?}

    \vspace{1.5em}
    \noindent{\Large \color{articolor} \textbf{Arti:} Bagaimana (pelajaran hari ini)?}
\end{convobox}

\begin{convobox}{78}
    \noindent
    \jpword{こんど}{今度}{Kondo} \jpword{にほん}{日本}{Nihon} \jpkana{で}{de} 
    \jpkana{カフェに}{kafe ni} 
    \jpword{い}{行}{i}\jpkana{くことがあったら、}{ku koto ga attara,}

    \vspace{1.5em}
    \noindent{\Large \color{articolor} \textbf{Arti:} Jika lain kali kamu ada kesempatan pergi ke kafe di Jepang,}
\end{convobox}

\begin{lessonbox}{✈️ Petualangan Menantimu!}
    Jepang memiliki banyak sekali 'Themed Cafe' (Kafe bertema) lucu seperti Pokemon Cafe, Kirby Cafe, atau Sanrio Cafe yang pastiiiiii bikin anak-anak melonjak kegirangan! Jangan lupa reservasi (pesan tempat) dari jauh hari karena kafe tematik selalu penuh!
\end{lessonbox}

% ================= PAIR 40 =================
\begin{convobox}{79}
    \noindent
    \jpkana{ぜひ}{Zehi} 
    \jpkana{この}{kono} \jpword{どうが}{動画}{douga} \jpkana{で}{de} 
    \jpword{まな}{学}{mana}\jpkana{んだ}{nda} 
    \jpkana{フレーズを}{fureezu o} 
    \jpword{つか}{使}{tsuka}\jpkana{って}{tte} 
    \jpword{ちゅうもん}{注文}{chuumon} \jpkana{してみてくださいね。}{shite mite kudasai ne.}

    \vspace{1.5em}
    \noindent{\Large \color{articolor} \textbf{Arti:} Pastikan (jangan ragu) untuk memesan dengan menggunakan frasa-frasa yang telah dipelajari di video ini ya.}
\end{convobox}

\begin{convobox}{80}
    \noindent
    \jpkana{それでは}{Sore dewa} 
    \jpkana{また}{mata} 
    \jpword{じかい}{次回}{jikai} \jpkana{の}{no} 
    \jpword{どうが}{動画}{douga} \jpkana{で}{de} 
    \jpkana{お}{o}\jpword{あ}{会}{a}\jpkana{いしましょう！}{i shimashou!}

    \vspace{1.5em}
    \noindent{\Large \color{articolor} \textbf{Arti:} Kalau begitu, sampai jumpa lagi di video (materi) berikutnya!}
\end{convobox}

% --- SUMMARY BOX ---
\vspace{2em}
\begin{tcolorbox}[
    colback=blue!5!white,
    colframe=blue!80!black,
    boxrule=3pt,
    arc=5mm,
    title={\huge\bfseries \sffamily 🌟 Rangkuman Misi Kafe Keluarga 🌟},
    fonttitle=\color{white},
    attach boxed title to top center={yshift=-4mm},
    boxed title style={colback=blue!80!black, rounded corners},
    left=5mm, right=5mm, top=7mm, bottom=5mm
]
\Large \textbf{Hebat! Keluarga Kalian Sudah Siap Nongkrong di Kafe Tokyo!}
\begin{itemize}
    \item \textbf{Cari Meja Dulu:} Selalu amankan mejamu sebelum maju memesan.
    \item \textbf{Pajaknya Beda:} \textit{Ten-nai} (Makan di tempat) dan \textit{Mochikaeri} (Bawa pulang) memiliki aturan pajak yang berbeda di kasir.
    \item \textbf{Rumus Pesanan:} [Nama Menu] + [Jumlah: Hitotsu/Futatsu] + \textit{Onegaishimasu}.
    \item \textbf{Pertanyaan Suhu:} \textit{Hotto} (Panas) atau \textit{Aisu} (Dingin).
    \item \textbf{Customizing:} \textit{... Dake} (Hanya ...), atau \textit{Ie, daijoubu desu} (Tolak dengan halus).
    \item \textbf{Etika Kebersihan:} Kembalikan nampan dan gelas kotor ke area 'Henkyakuguchi'.
\end{itemize}
Dengan berlatih mengucapkan ini keras-keras di rumah, kalian pasti akan memesan makanan seperti warga lokal Jepang sungguhan. Selamat berlibur dan menikmati kopi serta jus yang lezat!
\end{tcolorbox}

\newpage
\chapter{Phrases at Japanese Convenience Store}

% ================= TITLE: CONVERSATION =================
\vspace{1em}
\begin{tcolorbox}[colback=boxframe, colframe=boxframe, rounded corners, boxrule=0pt, arc=3mm, left=2mm, top=2mm, bottom=2mm]
    \Large \bfseries \color{white} 🏪 1. The Ultimate Konbini Flow (Alur Kasir)
\end{tcolorbox}
\vspace{1em}

% ================= PAIR 1 =================
\begin{convobox}{1}
    \noindent
    \jpword{つぎ}{次}{Tsugi} \jpkana{の}{no} 
    \jpword{かた}{方}{kata} \jpkana{どうぞ。}{douzo.}

    \vspace{1.5em}
    \noindent{\Large \color{articolor} \textbf{Arti:} Orang berikutnya, silakan (maju).}
\end{convobox}

\begin{convobox}{2}
    \noindent
    \jpkana{いらっしゃいませ。}{Irasshaimase.} 
    \jpkana{お}{o}\jpword{あず}{預}{azu}\jpkana{かりいたします。}{kari itashimasu.}

    \vspace{1.5em}
    \noindent{\Large \color{articolor} \textbf{Arti:} Selamat datang. Saya akan menerima/menyimpan (barang belanjaan) Anda.}
\end{convobox}

\begin{lessonbox}{⛩️ Panduan Turis: Garis Antrean \& Panggilan Kasir!}
    Halo Ayah, Bunda, dan teman-teman! Di Jepang, **kita tidak boleh langsung maju dan meletakkan barang di meja kasir**. Kamu harus berdiri di belakang garis atau stiker jejak kaki di lantai. \\
    
    Tunggu sampai kasir menatapmu dan berkata \textbf{「次の方どうぞ」 (Tsugi no kata douzo - Orang berikutnya silakan)}. Setelah mendengar kata ajaib itu, barulah kamu boleh maju ke meja kasir!
\end{lessonbox}

% ================= PAIR 2 =================
\begin{convobox}{3}
    \noindent
    \jpkana{お}{O}\jpword{べんとう}{弁当}{bentou} \jpword{あたた}{温}{atata}\jpkana{めますか？}{memasu ka?}

    \vspace{1.5em}
    \noindent{\Large \color{articolor} \textbf{Arti:} Apakah kotak bekalnya (Bento) mau dipanaskan?}
\end{convobox}

\begin{convobox}{4}
    \noindent
    \jpkana{はい、}{Hai,} \jpkana{お}{o}\jpword{ねが}{願}{nega}\jpkana{いします。}{ishimasu.}

    \vspace{1.5em}
    \noindent{\Large \color{articolor} \textbf{Arti:} Ya, tolong.}
\end{convobox}

\begin{lessonbox}{🍱 Panduan Turis: Memanaskan Makanan}
    Minimarket Jepang punya makanan instan terenak di dunia! Kalau kamu membeli Nasi (Bento), Pasta, atau Onigiri tertentu, kasir otomatis akan bertanya \textit{"Atatamemasu ka?"} (Mau dipanaskan?). \\
    
    Jawab dengan lantang: \textbf{「はい、お願いします」(Hai, onegaishimasu)} dan mereka akan memasukkannya ke microwave khusus berkecepatan tinggi!
\end{lessonbox}

% ================= PAIR 3 =================
\begin{convobox}{5}
    \noindent
    \jpkana{お}{O}\jpword{はし}{箸}{hashi} \jpkana{ご}{go}\jpword{りよう}{利用}{riyou} \jpkana{ですか？}{desu ka?}

    \vspace{1.5em}
    \noindent{\Large \color{articolor} \textbf{Arti:} Apakah Anda menggunakan (membutuhkan) sumpit?}
\end{convobox}

\begin{convobox}{6}
    \noindent
    \jpkana{はい、}{Hai,} \jpkana{お}{o}\jpword{ねが}{願}{nega}\jpkana{いします。}{ishimasu.}

    \vspace{1.5em}
    \noindent{\Large \color{articolor} \textbf{Arti:} Ya, tolong.}
\end{convobox}

\begin{lessonbox}{🥢 Panduan Turis: Alat Makan Sesuai Kebutuhan}
    Kadang kasir tidak menyebut "O-hashi" (Sumpit). Kalau kamu beli yogurt, mereka akan bilang \textbf{「スプーン」 (Supuun - Sendok)}. Kalau beli pasta, mereka bilang \textbf{「フォーク」 (Fooku - Garpu)}. Dengarkan awalannya ya, dan cukup jawab "Hai, onegaishimasu" untuk mendapatkannya gratis!
\end{lessonbox}

% ================= PAIR 4 =================
\begin{convobox}{7}
    \noindent
    \jpkana{ポイントカード}{Pointo kaado} \jpkana{お}{o}\jpword{も}{持}{mo}\jpkana{ちですか？}{chi desu ka?}

    \vspace{1.5em}
    \noindent{\Large \color{articolor} \textbf{Arti:} Apakah Anda membawa kartu poin?}
\end{convobox}

\begin{convobox}{8}
    \noindent
    \jpkana{いえ、}{Ie,} \jpword{も}{持}{mo}\jpkana{っていません。}{tte imasen.}

    \vspace{1.5em}
    \noindent{\Large \color{articolor} \textbf{Arti:} Tidak, saya tidak membawanya.}
\end{convobox}

\begin{lessonbox}{💳 Panduan Turis: Lewati Saja Bagian Ini!}
    Turis tidak butuh kartu poin (T-Point, d-Point, Ponta). Agar prosesnya cepat, kamu bisa memotong dengan sopan dengan menggelengkan kepala dan tersenyum sambil bilang \textbf{「いえ、大丈夫です」 (Ie, daijoubu desu - Tidak usah/Tidak ada)}.
\end{lessonbox}

% ================= PAIR 5 =================
\begin{convobox}{9}
    \noindent
    \jpkana{そちらの}{Sochira no} 
    \jpword{ねんれい}{年齢}{nenrei} \jpword{かくにん}{確認}{kakunin} \jpkana{ボタンを}{botan o} 
    \jpword{お}{押}{o}\jpkana{していただけますでしょうか？}{shite itadakemasu deshou ka?}

    \vspace{1.5em}
    \noindent{\Large \color{articolor} \textbf{Arti:} Bisakah Anda menekan tombol konfirmasi umur yang ada di sebelah sana?}
\end{convobox}

\begin{lessonbox}{🚨 PANDUAN KRITIS: Jebakan Layar Sentuh!}
    \textbf{Perhatian khusus untuk orang tua!} Jika Ayah/Bunda membeli bir, minuman alkohol kalengan (seperti Strong Zero, Chuhai), atau rokok, tiba-tiba mesin kasir akan berbunyi dan kasir menunjuk layar sentuh ke arahmu. \\
    
    Itu adalah **Nenrei Kakunin (Konfirmasi Umur)**. Di Jepang, kamu harus berusia di atas 20 tahun. Di layar akan muncul tulisan \textbf{「私は20歳以上です」 (Saya di atas 20 tahun)}. \textbf{KAMU HARUS MENYENTUH TOMBOL DI LAYAR ITU!} Jika tidak disentuh, kasir tidak akan bisa melanjutkan transaksi. Sentuh saja layarnya, dan transaksi lanjut!
\end{lessonbox}

% ================= PAIR 6 =================
\begin{convobox}{10}
    \noindent
    \jpword{ふくろ}{袋}{Fukuro} \jpkana{ご}{go}\jpword{りよう}{利用}{riyou} \jpkana{ですか？}{desu ka?}

    \vspace{1.5em}
    \noindent{\Large \color{articolor} \textbf{Arti:} Apakah Anda menggunakan kantong plastik?}
\end{convobox}

\begin{convobox}{11}
    \noindent
    \jpkana{はい、}{Hai,} \jpkana{お}{o}\jpword{ねが}{願}{nega}\jpkana{いします。}{ishimasu.}

    \vspace{1.5em}
    \noindent{\Large \color{articolor} \textbf{Arti:} Ya, tolong.}
\end{convobox}

\begin{lessonbox}{🛍️ Panduan Turis: Kantong Plastik (Fukuro)}
    Selalu ingat kata kunci \textbf{FUKURO}. Karena tanganmu mungkin sudah penuh memegang payung atau kamera, meminta plastik adalah ide bagus meskipun berbayar.
\end{lessonbox}

% ================= PAIR 7 =================
\begin{convobox}{12}
    \noindent
    \jpword{いちまい}{1枚}{Ichi-mai} \jpword{さんえん}{3円}{san-en} 
    \jpkana{ですが}{desu ga} 
    \jpkana{よろしいでしょうか？}{yoroshii deshou ka?}

    \vspace{1.5em}
    \noindent{\Large \color{articolor} \textbf{Arti:} 1 lembarnya seharga 3 yen, apakah tidak apa-apa (apakah Anda setuju)?}
\end{convobox}

\begin{convobox}{13}
    \noindent
    \jpkana{はい。}{Hai.}

    \vspace{1.5em}
    \noindent{\Large \color{articolor} \textbf{Arti:} Ya.}
\end{convobox}

\begin{lessonbox}{💴 Panduan Turis: Kasir yang Transparan}
    Kasir Jepang sangat jujur. Jika ada biaya tambahan (seperti plastik 3 yen sampai 5 yen), mereka akan selalu meminta persetujuanmu dengan kata \textbf{「よろしいでしょうか」 (Yoroshii deshou ka - Apakah tidak apa-apa?)}. Jawab saja "Hai!" dengan percaya diri.
\end{lessonbox}

% ================= PAIR 8 =================
\begin{convobox}{14}
    \noindent
    \jpkana{お}{O}\jpword{かいけい}{会計}{kaikei} 
    \jpword{ななひゃくななじゅうさんえん}{773円}{nana-hyaku nana-juu san-en} 
    \jpkana{でございます。}{de gozaimasu.}

    \vspace{1.5em}
    \noindent{\Large \color{articolor} \textbf{Arti:} Total pembayarannya adalah 773 yen.}
\end{convobox}

\vspace{2em}
\begin{tcolorbox}[colback=boxframe, colframe=boxframe, rounded corners, boxrule=0pt, arc=3mm, left=2mm, top=2mm, bottom=2mm]
    \Large \bfseries \color{white} 💬 2. Alternatif Jawaban (Saat Berubah Pikiran)
\end{tcolorbox}
\vspace{1em}

\begin{convobox}{15}
    \noindent
    \jpkana{いえ、}{Ie,} \jpword{だいじょうぶ}{大丈夫}{daijoubu} \jpkana{です。}{desu.}

    \vspace{1.5em}
    \noindent{\Large \color{articolor} \textbf{Arti:} Tidak, tidak usah (Tidak apa-apa).}
\end{convobox}

\begin{convobox}{16}
    \noindent
    \jpkana{やっぱり}{Yappari} 
    \jpword{い}{要}{i}\jpkana{りません。}{rimasen.}

    \vspace{1.5em}
    \noindent{\Large \color{articolor} \textbf{Arti:} Sepertinya aku tidak jadi butuh.}
\end{convobox}

\begin{lessonbox}{🔄 Panduan Turis: "Eh, Gak Jadi Deh!"}
    Bagaimana kalau kamu awalnya minta plastik, lalu ingat kamu bawa tas kain sendiri? 
    Gunakan kata ajaib \textbf{やっぱり (Yappari)} yang artinya "Sebenarnya/Gak jadi deh". Lalu sambung dengan \textbf{要りません (Irimasen - Saya tidak butuh)}. 
    
    "Yappari irimasen!" = "Eh, gak jadi butuh deh plastikannya!". Kasir akan langsung paham dan membatalkan plastiknya!
\end{lessonbox}


% --- SUMMARY BOX ---
\vspace{2em}
\begin{tcolorbox}[
    colback=blue!5!white,
    colframe=blue!80!black,
    boxrule=3pt,
    arc=5mm,
    title={\huge\bfseries \sffamily 🌟 Rangkuman Kelulusan Panduan Jepang! 🌟},
    fonttitle=\color{white},
    attach boxed title to top center={yshift=-4mm},
    boxed title style={colback=blue!80!black, rounded corners},
    left=5mm, right=5mm, top=7mm, bottom=5mm
]
\Large \textbf{SELAMAT! Kamu Sudah Lulus dan Siap Berlibur ke Jepang!}
\begin{itemize}
    \item \textbf{Antre dengan Benar:} Tunggu aba-aba \textit{Tsugi no kata douzo} sebelum maju ke meja kasir.
    \item \textbf{Siap Dipanaskan:} Bento/Nasi akan ditanya \textit{Atatamemasu ka?}. Jawab \textit{Hai, onegaishimasu!}
    \item \textbf{Jangan Lupa Tombol Layar:} Membeli alkohol = Sentuh layar \textit{Nenrei Kakunin} (Konfirmasi Umur 20+).
    \item \textbf{Ubah Pikiran:} Gunakan \textit{Yappari Irimasen} kalau kamu batal meminta kantong plastik.
\end{itemize}

Terima kasih sudah belajar dari Bab 1 hingga Bab 12 bersama buku **Andi - Japanese Handbook**. Perjalananmu mempelajari bahasa dan budaya Jepang sangat luar biasa. Bawalah buku ini bersamamu saat ke Jepang nanti, dan jangan ragu untuk berbicara dengan penduduk lokal! \textbf{気をつけて、行ってらっしゃい！ (Ki o tsukete, Itterasshai - Hati-hati di jalan dan selamat bersenang-senang!)}
\end{tcolorbox}

\newpage
\chapter{Revisit 1}

% ================= TOPIC 1 =================
\begin{lessonbox}{1. Simple sentence with です (Kalimat Dasar)}
    Halo Ayah, Bunda, dan Adik-adik! Kata \textbf{です (desu)} adalah kata ajaib yang dipakai di akhir kalimat agar terdengar sopan. Anggap saja artinya "Adalah". Kalau ditanya nama oleh orang Jepang, kamu tinggal sebut namamu dan tambah "desu"!
\end{lessonbox}
\begin{convobox}{1}
    \noindent
    \jpword{わたし}{私}{Watashi} \jpkana{は}{wa} 
    \jpkana{アンディ}{Andi} \jpkana{です。}{desu.}

    \vspace{1.5em}
    \noindent{\Large \color{articolor} \textbf{Arti:} Saya adalah Andi.}
\end{convobox}

% ================= TOPIC 2 =================
\begin{lessonbox}{2. The possessive の (Kepemilikan "Milik")}
    Saat kamu menemukan barangmu yang tertinggal di hotel, kamu bisa bilang "Ini milikku!". Partikel \textbf{の (no)} berfungsi menyambungkan pemilik dan barangnya.
\end{lessonbox}
\begin{convobox}{2}
    \noindent
    \jpword{わたし}{私}{Watashi} \jpkana{の}{no} 
    \jpkana{かばん}{kaban} \jpkana{です。}{desu.}

    \vspace{1.5em}
    \noindent{\Large \color{articolor} \textbf{Arti:} Ini adalah tas milik saya.}
\end{convobox}

% ================= TOPIC 3 =================
\begin{lessonbox}{3. Adjective basics (Kata Sifat Dasar)}
    Wah, makanannya enak! Pemandangannya indah! Untuk memuji sesuatu di Jepang, sebutkan bendanya lalu tambahkan kata sifat dan tutup dengan \textit{desu}.
\end{lessonbox}
\begin{convobox}{3}
    \noindent
    \jpkana{この}{Kono} \jpword{すし}{寿司}{sushi} \jpkana{は}{wa} 
    \jpkana{おいしい}{oishii} \jpkana{です。}{desu.}

    \vspace{1.5em}
    \noindent{\Large \color{articolor} \textbf{Arti:} Sushi ini lezat (enak).}
\end{convobox}

% ================= TOPIC 4 =================
\begin{lessonbox}{4. Adjectives - negative and linking (Kata Sifat: Tidak!)}
    Bagaimana kalau ternyata makanannya TIDAK pedas? Untuk kata sifat berakhiran "i", ubah huruf "i" terakhir menjadi \textbf{〜くない (〜kunai)}. 
\end{lessonbox}
\begin{convobox}{4}
    \noindent
    \jpkana{この}{Kono} \jpkana{カレー}{karee} \jpkana{は}{wa} 
    \jpword{から}{辛}{kara}\jpkana{くない}{kunai} \jpkana{です。}{desu.}

    \vspace{1.5em}
    \noindent{\Large \color{articolor} \textbf{Arti:} Kari ini tidak pedas.}
\end{convobox}

% ================= TOPIC 5 =================
\begin{lessonbox}{5. Adjectives: like, dislike, and want (Suka \& Ingin)}
    Kalau Adik-adik ingin es krim atau mainan baru, kalian bisa gunakan \textbf{好き (Suki - Suka)} atau \textbf{欲しい (Hoshii - Ingin)}. Jangan lupa pakai partikel \textit{ga} ya!
\end{lessonbox}
\begin{convobox}{5}
    \noindent
    \jpword{わたし}{私}{Watashi} \jpkana{は}{wa} 
    \jpkana{アイス}{aisu} \jpkana{が}{ga} 
    \jpword{す}{好}{su}\jpkana{き}{ki} \jpkana{です。}{desu.}

    \vspace{1.5em}
    \noindent{\Large \color{articolor} \textbf{Arti:} Saya suka es krim.}
\end{convobox}

% ================= TOPIC 6 =================
\begin{lessonbox}{6. Simple questions (Bertanya Darurat)}
    Saat jalan-jalan, Ayah dan Bunda pasti butuh bertanya arah. Tambahkan saja \textbf{か (ka)} di akhir kalimat, dan kalimatmu otomatis berubah menjadi pertanyaan!
\end{lessonbox}
\begin{convobox}{6}
    \noindent
    \jpkana{トイレ}{Toire} \jpkana{は}{wa} 
    \jpkana{どこ}{doko} \jpkana{ですか？}{desu ka?}

    \vspace{1.5em}
    \noindent{\Large \color{articolor} \textbf{Arti:} Toilet ada di mana?}
\end{convobox}

% ================= TOPIC 7 =================
\begin{lessonbox}{7. Basic Verbs: ます and ません (Kata Kerja Dasar)}
    Untuk bercerita tentang apa yang akan kamu lakukan, gunakan akhiran \textbf{〜ます (〜masu)}. Kalau TIDAK akan melakukannya, ubah menjadi \textbf{〜ません (〜masen)}.
\end{lessonbox}
\begin{convobox}{7}
    \noindent
    \jpkana{ラーメン}{Raamen} \jpkana{を}{o} 
    \jpword{た}{食}{ta}\jpkana{べます。}{bemasu.}

    \vspace{1.5em}
    \noindent{\Large \color{articolor} \textbf{Arti:} Saya (akan) makan ramen.}
\end{convobox}

% ================= TOPIC 8 =================
\begin{lessonbox}{8. Basic particles (を dan で)}
    Partikel adalah lem yang merekatkan kalimat! \textbf{を (o)} menempel pada objek yang kamu makan/beli. Sedangkan \textbf{で (de)} menempel pada TEMPAT kamu melakukan aktivitas itu.
\end{lessonbox}
\begin{convobox}{8}
    \noindent
    \jpkana{ホテル}{Hoteru} \jpkana{で}{de} 
    \jpword{あさ}{朝}{asa}\jpkana{ごはんを}{gohan o} 
    \jpword{た}{食}{ta}\jpkana{べます。}{bemasu.}

    \vspace{1.5em}
    \noindent{\Large \color{articolor} \textbf{Arti:} Saya makan sarapan di hotel.}
\end{convobox}

% ================= TOPIC 9 =================
\begin{lessonbox}{9. Basic particles (へ/に - Arah \& Tujuan)}
    Saat naik kereta, perhatikan partikel \textbf{へ (e)} atau \textbf{に (ni)}. Keduanya berarti "Ke". Ini sangat penting untuk memberi tahu supir taksi ke mana keluargamu mau pergi!
\end{lessonbox}
\begin{convobox}{9}
    \noindent
    \jpword{きょうと}{京都}{Kyouto} \jpkana{に}{ni} 
    \jpword{い}{行}{i}\jpkana{きます。}{kimasu.}

    \vspace{1.5em}
    \noindent{\Large \color{articolor} \textbf{Arti:} Saya pergi ke Kyoto.}
\end{convobox}

% ================= TOPIC 10 =================
\begin{lessonbox}{10. Basic particles (も/と - Juga \& Dan)}
    Saat berbelanja suvenir, kadang kita mau beli banyak hal. Gunakan \textbf{と (to)} untuk bilang "Dan", serta \textbf{も (mo)} untuk bilang "Juga".
\end{lessonbox}
\begin{convobox}{10}
    \noindent
    \jpkana{これ}{Kore} \jpkana{と}{to} 
    \jpkana{それ}{sore} \jpkana{を}{o} 
    \jpword{くだ}{下}{kuda}\jpkana{さい。}{sai.}

    \vspace{1.5em}
    \noindent{\Large \color{articolor} \textbf{Arti:} Tolong (berikan saya) ini dan itu.}
\end{convobox}

% ================= TOPIC 11 =================
\begin{lessonbox}{11. Existence Verbs: ある/いる (Ada!)}
    Di taman ada rusa? Di toko ada payung? Gunakan \textbf{います (Imasu)} jika yang ada adalah makhluk hidup (manusia/hewan). Gunakan \textbf{あります (Arimasu)} untuk benda mati (gedung/barang).
\end{lessonbox}
\begin{convobox}{11}
    \noindent
    \jpkana{あそこ}{Asoko} \jpkana{に}{ni} 
    \jpkana{パンダ}{panda} \jpkana{が}{ga} 
    \jpkana{います。}{imasu.}

    \vspace{1.5em}
    \noindent{\Large \color{articolor} \textbf{Arti:} Di sana ada Panda.}
\end{convobox}

% ================= TOPIC 12 =================
\begin{lessonbox}{12. Past tense: でした (Masa Lalu: Benda/Sifat)}
    Liburan kemarin hujan? Atau makanannya sangat enak? Jika peristiwanya SUDAH BERLALU, kata \textit{desu} diubah menjadi \textbf{でした (deshita)}.
\end{lessonbox}
\begin{convobox}{12}
    \noindent
    \jpword{きのう}{昨日}{Kinou} \jpkana{は}{wa} 
    \jpword{あめ}{雨}{ame} \jpkana{でした。}{deshita.}

    \vspace{1.5em}
    \noindent{\Large \color{articolor} \textbf{Arti:} Kemarin (cuacanya) hujan.}
\end{convobox}

% ================= TOPIC 13 =================
\begin{lessonbox}{13. Past tense: ました (Masa Lalu: Kata Kerja)}
    Sama seperti poin sebelumnya, kalau kegiatan (kata kerja) sudah selesai kamu lakukan hari itu, ubah akhiran \textit{~masu} menjadi \textbf{〜ました (〜mashita)}.
\end{lessonbox}
\begin{convobox}{13}
    \noindent
    \jpkana{たこやき}{Takoyaki} \jpkana{を}{o} 
    \jpword{か}{買}{ka}\jpkana{いました。}{imashita.}

    \vspace{1.5em}
    \noindent{\Large \color{articolor} \textbf{Arti:} Saya sudah membeli Takoyaki.}
\end{convobox}

% ================= TOPIC 14 =================
\begin{lessonbox}{14. Invitations (Mari Mengajak!)}
    Ayah dan Bunda mau mengajak anak-anak pergi ke Disneyland? Ubah kata kerjanya menjadi \textbf{〜ましょう (〜mashou)} yang artinya "Mari kita... / Ayo kita...".
\end{lessonbox}
\begin{convobox}{14}
    \noindent
    \jpword{いっしょ}{一緒}{Issho} \jpkana{に}{ni} 
    \jpword{い}{行}{i}\jpkana{きませんか？}{kimasen ka?}

    \vspace{1.5em}
    \noindent{\Large \color{articolor} \textbf{Arti:} Maukah pergi bersama-sama?}
\end{convobox}

% ================= TOPIC 15 =================
\begin{lessonbox}{15. Verb form: て (Bentuk Te untuk Menyambung)}
    Bentuk \textit{Te} sangat ajaib! Bisa dipakai untuk menyambung kalimat ("Saya makan, \textit{lalu} pergi") atau meminta tolong secara kasual.
\end{lessonbox}
\begin{convobox}{15}
    \noindent
    \jpkana{ちょっと}{Chotto} \jpword{ま}{待}{ma}\jpkana{って}{tte} 
    \jpword{くだ}{下}{kuda}\jpkana{さい。}{sai.}

    \vspace{1.5em}
    \noindent{\Large \color{articolor} \textbf{Arti:} Tolong tunggu sebentar.}
\end{convobox}

% ================= TOPIC 16 =================
\begin{lessonbox}{16. Verb form: ている (Sedang Berlangsung)}
    Jika saat ini kamu SEDANG melakukan sesuatu (misalnya sedang duduk di kereta), gunakan bentuk \textbf{〜ています (〜te imasu)}.
\end{lessonbox}
\begin{convobox}{16}
    \noindent
    \jpword{いま}{今}{Ima}\jpkana{、}{,} 
    \jpword{でんしゃ}{電車}{densha} \jpkana{に}{ni} 
    \jpword{の}{乗}{no}\jpkana{っています。}{tte imasu.}

    \vspace{1.5em}
    \noindent{\Large \color{articolor} \textbf{Arti:} Sekarang, saya sedang naik kereta.}
\end{convobox}

% ================= TOPIC 17 =================
\begin{lessonbox}{17. Requests and permission (Meminta Izin Terdengar Sopan)}
    Saat melihat maskot lucu atau pemandangan indah di kuil, kadang kita tidak yakin boleh memfotonya atau tidak. Minta izinlah dengan sopan pakai \textbf{〜てもいいですか (〜te mo ii desu ka?)}.
\end{lessonbox}
\begin{convobox}{17}
    \noindent
    \jpword{しゃしん}{写真}{Shashin} \jpkana{を}{o} 
    \jpword{と}{撮}{to}\jpkana{ってもいいですか？}{tte mo ii desu ka?}

    \vspace{1.5em}
    \noindent{\Large \color{articolor} \textbf{Arti:} Bolehkah saya mengambil foto?}
\end{convobox}

% --- SUMMARY BOX CLOSING ---
\vspace{2em}
\begin{tcolorbox}[
    colback=green!5!white,
    colframe=green!60!black,
    boxrule=3pt,
    arc=5mm,
    title={\huge\bfseries \sffamily 🌸 Selesai! Selamat Berlibur! 🌸},
    fonttitle=\color{white},
    attach boxed title to top center={yshift=-4mm},
    boxed title style={colback=green!60!black, rounded corners},
    left=5mm, right=5mm, top=7mm, bottom=5mm
]
\Large Keluarga Hebat! Kalian telah menyelesaikan seluruh pelajaran penting dalam buku \textbf{Andi - Japanese Handbook} ini. Dengan bekal tata bahasa dasar dan kumpulan kosakata penting ini, pengalaman berlibur kalian di Jepang dijamin akan menjadi jauh lebih seru, aman, dan berkesan. Jangan takut mencoba berbicara dengan penduduk lokal ya. \textbf{日本の旅行を楽しんでください！ (Nihon no ryokou o tanoshinde kudasai - Selamat menikmati liburan di Jepang!)}
\end{tcolorbox}

\newpage
\chapter{Revisit 2}

\begin{center}
    {\Huge \bfseries Andi - Japanese Handbook}\\[0.5em]
    {\LARGE \color{boxframe} \textbf{Bab Spesial 2: Kilas Balik Tata Bahasa (Lanjutan)}}
\end{center}
\vspace{1em}

% ================= TOPIC 1 =================
\begin{lessonbox}{1. Frequency words (Seberapa Sering?)}
    Halo lagi, petualang cilik! Kalau mau cerita seberapa sering kamu pergi ke konbini, gunakan kata frekuensi. \textbf{よく (Yoku)} = Sering, \textbf{ときどき (Tokidoki)} = Kadang-kadang.
\end{lessonbox}
\begin{convobox}{1}
    \noindent
    \jpword{わたし}{私}{Watashi} \jpkana{は}{wa} 
    \jpkana{よく}{yoku} 
    \jpkana{コンビニ}{konbini} \jpkana{に}{ni} 
    \jpword{い}{行}{i}\jpkana{きます。}{kimasu.}

    \vspace{1.5em}
    \noindent{\Large \color{articolor} \textbf{Arti:} Saya sering pergi ke minimarket.}
\end{convobox}

% ================= TOPIC 2 =================
\begin{lessonbox}{2. Wanting to do something (Ingin Melakukan Sesuatu)}
    Bunda, kalau si kecil punya keinginan, ajarkan mereka mengganti akhiran \textit{〜masu} menjadi \textbf{〜たい (〜tai)}. \textit{Tabemasu} (makan) menjadi \textit{Tabetai} (ingin makan)!
\end{lessonbox}
\begin{convobox}{2}
    \noindent
    \jpword{すし}{寿司}{Sushi} \jpkana{が}{ga} 
    \jpword{た}{食}{ta}\jpkana{べたいです。}{betai desu.}

    \vspace{1.5em}
    \noindent{\Large \color{articolor} \textbf{Arti:} Saya ingin makan sushi.}
\end{convobox}

% ================= TOPIC 3 =================
\begin{lessonbox}{3. で: Time, Method, Material (Partikel Serbaguna "De")}
    Partikel \textbf{で (De)} sangat sakti! Bisa berarti "dengan alat" (dengan bus), "dari bahan" (dari kayu), atau "batas waktu" (dalam 1 jam).
\end{lessonbox}

\begin{convobox}{3}
    \noindent
    \jpkana{バス}{Basu} \jpkana{で}{de} 
    \jpword{えき}{駅}{eki} \jpkana{に}{ni} 
    \jpword{い}{行}{i}\jpkana{きます。}{kimasu.}

    \vspace{1.5em}
    \noindent{\Large \color{articolor} \textbf{Arti:} Saya pergi ke stasiun DENGAN (menggunakan) bus.}
\end{convobox}

% ================= TOPIC 4 =================
\begin{lessonbox}{4. Casual speech (Bahasa Kasual)}
    Saat ngobrol santai dengan teman baru, buang kata \textit{Desu} dan \textit{Masu}! Ini bikin kamu terdengar ramah dan akrab.
\end{lessonbox}
\begin{convobox}{4}
    \noindent
    \jpkana{これ、}{Kore,} \jpkana{とても}{totemo} \jpkana{おいしい！}{oishii!}

    \vspace{1.5em}
    \noindent{\Large \color{articolor} \textbf{Arti:} Ini, enak banget!}
\end{convobox}

% ================= TOPIC 5 =================
\begin{lessonbox}{5. Casual Speech \#2 (Bertanya Kasual)}
    Untuk bertanya dalam bahasa santai, jangan pakai \textit{〜ka?}. Cukup naikkan nada suaramu di akhir kata seperti orang kaget!
\end{lessonbox}
\begin{convobox}{5}
    \noindent
    \jpword{あした}{明日}{Ashita}\jpkana{、}{,} 
    \jpword{い}{行}{i}\jpkana{く？}{ku?}

    \vspace{1.5em}
    \noindent{\Large \color{articolor} \textbf{Arti:} Besok, kamu pergi (berangkat)?}
\end{convobox}

% ================= TOPIC 6 =================
\begin{lessonbox}{6. Time Expressions (Hari \& Tanggal!)}
    \textbf{Info Penting Turis!} Di Jepang, cara membaca tanggal 1 sampai 10 sangat spesial:
    1 (Tsuitachi), 2 (Futsuka), 3 (Mikka), 4 (Yokka), 5 (Itsuka). 
    Hari ini = \textbf{Kyou}, Besok = \textbf{Ashita}, Kemarin = \textbf{Kinou}.
\end{lessonbox}
\begin{convobox}{6}
    \noindent
    \jpword{きょう}{今日}{Kyou} \jpkana{は}{wa} 
    \jpword{ついたち}{一日}{tsuitachi} \jpkana{です。}{desu.}

    \vspace{1.5em}
    \noindent{\Large \color{articolor} \textbf{Arti:} Hari ini adalah tanggal 1.}
\end{convobox}

% ================= TOPIC 7 =================
\begin{lessonbox}{7. Making Requests (Meminta Tolong)}
    Saat mau difotokan atau minta diambilkan barang, gunakan bentuk \textit{Te} lalu tambah \textbf{ください (Kudasai)}. Artinya "Tolong lakukan...".
\end{lessonbox}
\begin{convobox}{7}
    \noindent
    \jpkana{これ}{Kore} \jpkana{を}{o} 
    \jpword{み}{見}{mi}\jpkana{せて}{sete} 
    \jpword{くだ}{下}{kuda}\jpkana{さい。}{sai.}

    \vspace{1.5em}
    \noindent{\Large \color{articolor} \textbf{Arti:} Tolong perlihatkan (barang) ini.}
\end{convobox}

% ================= TOPIC 8 =================
\begin{lessonbox}{8. Finished/Not Yet (Sudah \& Belum)}
    "Apakah kereta Shinkansen-nya sudah datang?" "Belum!". Gunakan \textbf{もう (Mou)} untuk "Sudah" dan \textbf{まだ (Mada)} untuk "Belum".
\end{lessonbox}
\begin{convobox}{8}
    \noindent
    \jpkana{もう}{Mou} \jpword{た}{食}{ta}\jpkana{べました。}{bemashita.} 
    \jpkana{まだ}{Mada} \jpkana{です。}{desu.}

    \vspace{1.5em}
    \noindent{\Large \color{articolor} \textbf{Arti:} Saya SUDAH makan. / Saya BELUM.}
\end{convobox}

% ================= TOPIC 9 =================
\begin{lessonbox}{9. Change - なる/する (Menjadi...)}
    Jika cuaca berubah dari panas menjadi dingin, pakai kata \textbf{なる (Naru)}. Akhiran \textit{i} berubah jadi \textit{ku}. \textit{Samui} (Dingin) -> \textit{Samuku naru}.
\end{lessonbox}
\begin{convobox}{9}
    \noindent
    \jpword{さむ}{寒}{Samu}\jpkana{くなりました。}{ku narimashita.}

    \vspace{1.5em}
    \noindent{\Large \color{articolor} \textbf{Arti:} (Cuacanya) telah menjadi dingin.}
\end{convobox}

% ================= TOPIC 10 =================
\begin{lessonbox}{10. も: No(thing), even (Tidak Ada Apa-apa!)}
    Kata tanya + \textbf{も (mo)} + Negatif = Kosong/Tidak ada! \\
    Nani (Apa) + Mo + Arimasen = Nani mo arimasen (Tidak ada apa-apa!).
\end{lessonbox}
\begin{convobox}{10}
    \noindent
    \jpword{なに}{何}{Nani} \jpkana{も}{mo} 
    \jpkana{ありません。}{arimasen.}

    \vspace{1.5em}
    \noindent{\Large \color{articolor} \textbf{Arti:} Tidak ada apa-apa.}
\end{convobox}

% ================= TOPIC 11 =================
\begin{lessonbox}{11. Particles: か and を (Atau \& Objek)}
    Kamu ingin jus apel ATAU susu? Gunakan partikel \textbf{か (ka)} di antara pilihanmu. Partikel \textbf{を (o)} menandakan barang yang kamu pilih.
\end{lessonbox}
\begin{convobox}{11}
    \noindent
    \jpkana{ジュース}{Juusu} \jpkana{か}{ka} 
    \jpkana{ミルク}{miruku} \jpkana{を}{o} 
    \jpword{の}{飲}{no}\jpkana{みます。}{mimasu.}

    \vspace{1.5em}
    \noindent{\Large \color{articolor} \textbf{Arti:} Saya minum jus ATAU susu.}
\end{convobox}

% ================= TOPIC 12 =================
\begin{lessonbox}{12. Particles: と and が (Dan \& Subjek Pribadi)}
    Kalau perginya bersama-sama, sambung dengan \textbf{と (to)} yang berarti "Dan". \textbf{が (Ga)} adalah penunjuk siapa yang melakukan kegiatan itu.
\end{lessonbox}
\begin{convobox}{12}
    \noindent
    \jpword{ちち}{父}{Chichi} \jpkana{と}{to} 
    \jpword{はは}{母}{haha} \jpkana{が}{ga} 
    \jpword{い}{行}{i}\jpkana{きます。}{kimasu.}

    \vspace{1.5em}
    \noindent{\Large \color{articolor} \textbf{Arti:} Ayah DAN ibu akan pergi.}
\end{convobox}

% ================= TOPIC 13 =================
\begin{lessonbox}{13. Particles ね、よ、わ、なあ (Partikel Ekspresi)}
    Di akhir kalimat, \textbf{ね (ne)} berarti "ya kan?" (meminta setuju), dan \textbf{よ (yo)} berarti "lho!" (kasih info baru). Bikin bahasamu makin keren!
\end{lessonbox}
\begin{convobox}{13}
    \noindent
    \jpkana{これ、}{Kore,} \jpword{たか}{高}{taka}\jpkana{いですね！}{i desu ne!}

    \vspace{1.5em}
    \noindent{\Large \color{articolor} \textbf{Arti:} Ini, mahal ya!}
\end{convobox}

% ================= TOPIC 14 =================
\begin{lessonbox}{14. Comparisons (Arah \& Membandingkan)}
    \textbf{Info Penting Turis!} Hafalkan arah ini: \textbf{Migi} (Kanan), \textbf{Hidari} (Kiri), \textbf{Massugu} (Lurus). Untuk membandingkan, gunakan \textbf{〜のほうがいい (〜no hou ga ii)}.
\end{lessonbox}
\begin{convobox}{14}
    \noindent
    \jpword{みぎ}{右}{Migi} \jpkana{の}{no} 
    \jpword{ほう}{方}{hou} \jpkana{が}{ga} 
    \jpkana{いいです。}{ii desu.}

    \vspace{1.5em}
    \noindent{\Large \color{articolor} \textbf{Arti:} Sebelah KANAN lebih bagus.}
\end{convobox}

% ================= TOPIC 15 =================
\begin{lessonbox}{15. Amounts (Menyebut Jumlah)}
    Taman bermainnya ramai? Gunakan \textbf{たくさん (Takusan - Banyak)} atau \textbf{すこし (Sukoshi - Sedikit)}.
\end{lessonbox}
\begin{convobox}{15}
    \noindent
    \jpword{ひと}{人}{Hito} \jpkana{が}{ga} 
    \jpkana{たくさん}{takusan} 
    \jpkana{います。}{imasu.}

    \vspace{1.5em}
    \noindent{\Large \color{articolor} \textbf{Arti:} Ada BANYAK orang.}
\end{convobox}

% ================= TOPIC 16 =================
\begin{lessonbox}{16. Lists of things (Mendaftar Banyak Barang)}
    Saat belanja suvenir banyak macamnya, gunakan \textbf{や (ya)} untuk bilang "A, B, dan lain-lain". Berbeda dengan "to" yang menyebutkan semua benda.
\end{lessonbox}
\begin{convobox}{16}
    \noindent
    \jpkana{パン}{Pan} \jpkana{や}{ya} 
    \jpkana{ケーキを}{keeki o} 
    \jpword{か}{買}{ka}\jpkana{います。}{kaimasu.}

    \vspace{1.5em}
    \noindent{\Large \color{articolor} \textbf{Arti:} Saya membeli roti, kue (dan lain-lain).}
\end{convobox}

% ================= TOPIC 17 =================
\begin{lessonbox}{17. (Groups of) People (Kumpulan Orang)}
    Untuk membuat jamak (lebih dari satu), tambahkan \textbf{たち (tachi)}. Watashi (Saya) -> Watashi-tachi (Kami). Kodomo (Anak) -> Kodomo-tachi (Anak-anak).
\end{lessonbox}
\begin{convobox}{17}
    \noindent
    \jpword{わたし}{私}{Watashi}\jpkana{たちは}{tachi wa} 
    \jpword{かぞく}{家族}{kazoku} \jpkana{です。}{desu.}

    \vspace{1.5em}
    \noindent{\Large \color{articolor} \textbf{Arti:} Kami adalah keluarga.}
\end{convobox}

% ================= TOPIC 18 =================
\begin{lessonbox}{18. More て usages (Bentuk "Te" untuk Sifat)}
    Bagaimana bilang "Murah dan Enak"? Ubah kata sifat pertama (Yasui) menjadi \textbf{Yasukute}, lalu sambung!
\end{lessonbox}
\begin{convobox}{18}
    \noindent
    \jpword{やす}{安}{Yasu}\jpkana{くて、}{kute,} 
    \jpkana{おいしいです。}{oishii desu.}

    \vspace{1.5em}
    \noindent{\Large \color{articolor} \textbf{Arti:} Murah dan enak.}
\end{convobox}

% ================= TOPIC 19 =================
\begin{lessonbox}{19. Using から (Karena / Mulai Dari)}
    Partikel \textbf{から (kara)} punya 2 arti hebat: "Dari" (waktu/tempat), atau "Karena" (menjelaskan alasan).
\end{lessonbox}
\begin{convobox}{19}
    \noindent
    \jpword{す}{好}{Su}\jpkana{き}{ki} 
    \jpkana{だからです。}{da kara desu.}

    \vspace{1.5em}
    \noindent{\Large \color{articolor} \textbf{Arti:} Karena (saya) suka.}
\end{convobox}

% ================= TOPIC 20 =================
\begin{lessonbox}{20. Linking sentences (Kata Sambung)}
    Agar ceritamu panjang seperti penduduk lokal, pakai kata sambung! \textbf{でも (Demo)} = Tapi. \textbf{それから (Sorekara)} = Kemudian/Lalu.
\end{lessonbox}
\begin{convobox}{20}
    \noindent
    \jpword{たか}{高}{Taka}\jpkana{いです。}{i desu.} 
    \jpkana{でも、}{Demo,} 
    \jpword{か}{買}{ka}\jpkana{います。}{kaimasu.}

    \vspace{1.5em}
    \noindent{\Large \color{articolor} \textbf{Arti:} Mahal. Tapi, saya membelinya.}
\end{convobox}

% ================= TOPIC 21 =================
\begin{lessonbox}{21. Describing nouns (Menjelaskan Kata Benda)}
    Berbeda dengan bahasa Indonesia, di Jepang penjelasannya diletakkan di DEPAN! "Orang yang memakai baju merah" -> Baju merah dipakai + Orang.
\end{lessonbox}
\begin{convobox}{21}
    \noindent
    \jpword{あか}{赤}{Aka}\jpkana{い}{i} 
    \jpword{ふく}{服}{fuku} \jpkana{を}{o} 
    \jpword{き}{着}{ki}\jpkana{た}{ta} 
    \jpword{ひと}{人}{hito}\jpkana{。}{.}

    \vspace{1.5em}
    \noindent{\Large \color{articolor} \textbf{Arti:} Orang yang memakai pakaian merah.}
\end{convobox}

% ================= TOPIC 22 =================
\begin{lessonbox}{22. Describing actions (Bicara Cara Melakukan)}
    "Berjalan dengan CEPAT". Ubah \textit{Hayai} (Cepat) menjadi \textbf{Hayaku} agar bisa menempel ke kata kerja!
\end{lessonbox}
\begin{convobox}{22}
    \noindent
    \jpword{はや}{早}{Haya}\jpkana{く}{ku} 
    \jpword{ある}{歩}{aru}\jpkana{きます。}{kimasu.}

    \vspace{1.5em}
    \noindent{\Large \color{articolor} \textbf{Arti:} Berjalan dengan cepat.}
\end{convobox}

% ================= TOPIC 23 =================
\begin{lessonbox}{23. Using でしょう (Mungkin / Ya Kan?)}
    Kalau kamu menebak cuaca besok atau meminta persetujuan halus, gunakan \textbf{でしょう (Deshou)}. 
\end{lessonbox}
\begin{convobox}{23}
    \noindent
    \jpword{あした}{明日}{Ashita} \jpkana{は}{wa} 
    \jpword{あめ}{雨}{ame} \jpkana{でしょう。}{deshou.}

    \vspace{1.5em}
    \noindent{\Large \color{articolor} \textbf{Arti:} Besok mungkin akan hujan.}
\end{convobox}

% --- SUMMARY BOX CLOSING ---
\vspace{2em}
\begin{tcolorbox}[
    colback=purple!5!white,
    colframe=purple!60!black,
    boxrule=3pt,
    arc=5mm,
    title={\huge\bfseries \sffamily ⛩️ Penutup: Siap Jelajahi Jepang! ⛩️},
    fonttitle=\color{white},
    attach boxed title to top center={yshift=-4mm},
    boxed title style={colback=purple!60!black, rounded corners},
    left=5mm, right=5mm, top=7mm, bottom=5mm
]
\Large Sempurna! Buku Saku Super Cepat seri 2 ini sudah mencakup hampir seluruh kebutuhan komunikasi dasar untuk bertahan dan bersenang-senang di Jepang! Ingat selalu rumus menanyakan arah, menghitung tanggal, dan meminta tolong. Bawalah keceriaan dan keberanian kalian dalam mempraktikkan bahasa Jepang bersama buku ini. \textbf{また日本でお会いしましょう！ (Mata Nihon de o-ai shimashou - Sampai jumpa lagi di Jepang!)}
\end{tcolorbox}

\newpage
\chapter{Revisit 3}

\begin{center}
    {\Huge \bfseries Andi - Japanese Handbook}\\[0.5em]
    {\LARGE \color{boxframe} \textbf{Bab Spesial 3: Keahlian Bertahan Hidup (Advanced)}}
\end{center}
\vspace{1em}

% ================= THEME 1 =================
\begin{lessonbox}{1. Menerima, Memberi \& Hadiah (Giving \& Receiving 1-5)}
    \textbf{Topik: Giving and receiving \#1 - \#5} \\
    Di Jepang, budaya memberi oleh-oleh (\textit{Omiyage}) sangat kental! 
    \begin{itemize}
        \item \textbf{あげる (Ageru)} = Saya memberi ke orang lain.
        \item \textbf{もらう (Morau)} = Saya menerima dari orang lain.
        \item \textbf{くれる (Kureru)} = Orang lain memberi ke saya!
    \end{itemize}
    \textbf{Tips Turis:} Kalau ada orang Jepang yang sangat baik membantumu membawakan koper atau memberimu peta, kamu bisa bilang \textit{"Tetsudatte kurete, arigatou!"} (Terima kasih sudah memberiku bantuan!). Jika kamu ingin memberikan suvenir dari Indonesia, katakan \textit{"Kore, agemasu!"} (Ini, saya berikan untukmu!).
\end{lessonbox}

\begin{convobox}{1}
    \noindent
    \jpword{ともだち}{友達}{Tomodachi} \jpkana{に}{ni} 
    \jpword{おみやげ}{お土産}{o-miyage} \jpkana{を}{o} 
    \jpkana{あげます。}{agemasu.}

    \vspace{1.5em}
    \noindent{\Large \color{articolor} \textbf{Arti:} Saya memberikan oleh-oleh kepada teman.}
\end{convobox}

\begin{convobox}{2}
    \noindent
    \jpword{てつだ}{手伝}{Tetsuda}\jpkana{って}{tte} 
    \jpkana{くれて、}{kurete,} 
    \jpkana{ありがとう。}{arigatou.}

    \vspace{1.5em}
    \noindent{\Large \color{articolor} \textbf{Arti:} Terima kasih sudah (memberikan) bantuan kepadaku.}
\end{convobox}

% ================= THEME 2 =================
\begin{lessonbox}{2. Menjelaskan Benda secara Detail (Sifat \& Klausa)}
    \textbf{Topik: More の usages, Relative Clauses, Describing nouns/actions, こと+が} \\
    Bagaimana cara bilang "Kereta yang pergi ke Tokyo" atau "Sushi yang dimakan Ayah"? Di Jepang, \textbf{semua penjelasannya ditaruh di DEPAN kata benda}! 
    \begin{itemize}
        \item \textit{Toukyou ni iku} (Pergi ke Tokyo) + \textit{Densha} (Kereta) = \textbf{Toukyou ni iku densha} (Kereta yang pergi ke Tokyo).
    \end{itemize}
    \textbf{Tips Turis:} Kalau kamu kehilangan barang, kamu bisa menjelaskannya dengan detail. Misalnya mencari tas: \textit{"Watashi ga katta kaban"} (Tas yang saya beli). Oh ya, partikel \textbf{の (No)} juga bisa menggantikan benda agar kita tidak repot mengulang kata. "Saya suka yang merah" = \textit{"Akai \textbf{no} ga suki desu"}.
\end{lessonbox}

\begin{convobox}{3}
    \noindent
    \jpword{とうきょう}{東京}{Toukyou} \jpkana{に}{ni} 
    \jpword{い}{行}{i}\jpkana{く}{ku} 
    \jpword{でんしゃ}{電車}{densha} \jpkana{は}{wa} 
    \jpkana{どれ}{dore} \jpkana{ですか？}{desu ka?}

    \vspace{1.5em}
    \noindent{\Large \color{articolor} \textbf{Arti:} Kereta yang pergi ke Tokyo yang mana?}
\end{convobox}

\begin{convobox}{4}
    \noindent
    \jpword{あか}{赤}{Aka}\jpkana{い}{i} \jpkana{の}{no} \jpkana{を}{o} 
    \jpword{くだ}{下}{kuda}\jpkana{さい。}{sai.}

    \vspace{1.5em}
    \noindent{\Large \color{articolor} \textbf{Arti:} Tolong berikan (barang) YANG merah.}
\end{convobox}

% ================= THEME 3 =================
\begin{lessonbox}{3. Bisa, Boleh, dan Harus! (Aturan \& Kemampuan)}
    \textbf{Topik: Potential form, Must \& Must Not, Commands, Suggestions, Explanations (んです)} \\
    Saat jalan-jalan, sangat penting memahami papan peringatan dan aturan!
    \begin{itemize}
        \item \textbf{Bisa (Potential):} Ubah \textit{~masu} menjadi \textit{~emasu}. Contoh: \textit{Taberaremasu} (Bisa dimakan).
        \item \textbf{Harus (Must):} \textit{〜なければならない (〜nakereba naranai)}. "Saya harus pergi!"
        \item \textbf{Saran (Suggestions):} \textit{〜ほうがいい (〜hou ga ii)}. "Lebih baik bawa payung!"
    \end{itemize}
    \textbf{Tips Turis:} Kalau kamu tersesat dan butuh menjelaskan situasimu, gunakan \textbf{〜んです (〜n desu)} di akhir kalimatmu. Ini memberi kesan "Tolong dengarkan alasanku!". Contoh: \textit{"Michi ni mayotta n desu!"} (Aku tersesat!).
\end{lessonbox}

\begin{convobox}{5}
    \noindent
    \jpkana{ここで}{Koko de} 
    \jpword{しゃしん}{写真}{shashin} \jpkana{を}{o} 
    \jpword{と}{撮}{to}\jpkana{ってはいけません。}{tte wa ikemasen.}

    \vspace{1.5em}
    \noindent{\Large \color{articolor} \textbf{Arti:} Tidak boleh (Dilarang) mengambil foto di sini.}
\end{convobox}

\begin{convobox}{6}
    \noindent
    \jpword{きっぷ}{切符}{Kippu} \jpkana{を}{o} 
    \jpword{か}{買}{ka}\jpkana{わなければなりません。}{wanakereba narimasen.}

    \vspace{1.5em}
    \noindent{\Large \color{articolor} \textbf{Arti:} (Saya) harus membeli tiket kereta.}
\end{convobox}

% ================= THEME 4 =================
\begin{lessonbox}{4. Bahasa Dewa ala Jepang (Honorific \& Humble)}
    \textbf{Topik: Honorific speech, Humble speech, Passive voice, Causative verbs} \\
    Pernah dengar pengumuman stasiun atau kasir yang bahasanya terdengar asing dan panjang sekali? Itu adalah \textbf{敬語 (Keigo / Bahasa Hormat)}!
    \begin{itemize}
        \item \textbf{Honorific (Meninggikan tamu):} \textit{O-machi kudasai} (Tolong tunggu).
        \item \textbf{Humble (Merendahkan diri):} \textit{Itashimasu} (Saya akan melakukan).
    \end{itemize}
    \textbf{Tips Turis:} Kamu \textbf{TIDAK PERLU} bisa berbicara dengan Keigo! Cukup ucapkan bentuk \textit{〜masu} biasa. Namun, biasakan telingamu mendengarnya agar kamu paham maksud pelayan toko. Oh ya, kalau kakimu tidak sengaja terinjak di kereta, kamu bisa pakai \textbf{Pasif (Passive Voice)}: \textit{"Ashi o fUmaremashita!"} (Kakiku terinjak!).
\end{lessonbox}

\begin{convobox}{7}
    \noindent
    \jpkana{しょうしょう}{Shoushou} \jpkana{お}{o}\jpword{ま}{待}{ma}\jpkana{ち}{chi} 
    \jpword{くだ}{下}{kuda}\jpkana{さい。}{sai.}

    \vspace{1.5em}
    \noindent{\Large \color{articolor} \textbf{Arti:} Mohon tunggu sebentar. (Bahasa Honorific Kasir)}
\end{convobox}

\begin{convobox}{8}
    \noindent
    \jpword{わたし}{私}{Watashi} \jpkana{が}{ga} 
    \jpword{あんない}{案内}{annai} \jpkana{いたします。}{itashimasu.}

    \vspace{1.5em}
    \noindent{\Large \color{articolor} \textbf{Arti:} Saya yang akan memandu Anda. (Bahasa Humble Staf)}
\end{convobox}

% ================= THEME 5 =================
\begin{lessonbox}{5. Waktu, Pergerakan, dan Batas (Time \& Movement)}
    \textbf{Topik: Going \& Coming, Noun Suffixes (まで, Time, でも, Specification, Other)} \\
    Bagaimana cara bilang pergerakan dari A ke B? 
    \begin{itemize}
        \item \textbf{〜まで (〜made)} = Sampai. Contoh: \textit{Toukyou made} (Sampai Tokyo).
        \item \textbf{〜ていく / 〜てくる (〜te iku / 〜te kuru)} = Pergi / Datang. \textit{Katte kimasu!} (Aku beli dulu ya lalu kembali ke sini!).
        \item \textbf{〜でも (〜demo)} = "Atau semacamnya". \textit{Koohii demo nomimasu ka?} (Mau minum kopi atau semacamnya?).
    \end{itemize}
    \textbf{Tips Turis:} Gunakan \textit{〜made} saat bertanya pada supir taksi: \textit{"Kono hoteru \textbf{made} onegaishimasu"} (Tolong antarkan \textbf{sampai} hotel ini).
\end{lessonbox}

\begin{convobox}{9}
    \noindent
    \jpword{えき}{駅}{Eki} \jpkana{まで}{made} 
    \jpword{い}{行}{i}\jpkana{ってきます。}{tte kimasu.}

    \vspace{1.5em}
    \noindent{\Large \color{articolor} \textbf{Arti:} Saya akan pergi ke stasiun (lalu akan kembali lagi ke sini).}
\end{convobox}

% ================= THEME 6 =================
\begin{lessonbox}{6. Menyambung Kalimat \& Alasan (Conjunctions)}
    \textbf{Topik: ため/ために, Cause and Effect, のに/だけど, から, Linking sentences, Conjunctions 1-2} \\
    Agar kamu bisa bercerita layaknya penduduk lokal, gunakan kata sambung (Conjunctions)!
    \begin{itemize}
        \item \textbf{〜ために (〜tame ni)} = Demi / Untuk. \textit{Kazoku no tame ni} (Demi keluarga).
        \item \textbf{だから (Dakara) / から (Kara)} = Oleh karena itu / Karena.
        \item \textbf{なのに (Na noni) / だけど (Dakedo)} = Padahal / Tetapi.
    \end{itemize}
    \textbf{Tips Turis:} \textit{Dakedo} sangat sering dipakai orang Jepang di bahasa sehari-hari. "Sepatu ini mahal, \textit{dakedo} nyaman!"
\end{lessonbox}

\begin{convobox}{10}
    \noindent
    \jpword{あめ}{雨}{Ame} \jpkana{が}{ga} 
    \jpword{ふ}{降}{fu}\jpkana{っている}{tte iru} 
    \jpkana{から、}{kara,} 
    \jpkana{タクシーで}{takushii de} 
    \jpword{い}{行}{i}\jpkana{きましょう。}{kimashou.}

    \vspace{1.5em}
    \noindent{\Large \color{articolor} \textbf{Arti:} KARENA sedang turun hujan, mari kita pergi naik taksi.}
\end{convobox}

% ================= THEME 7 =================
\begin{lessonbox}{7. Kelihatannya \& Katanya (Senses \& Guesses)}
    \textbf{Topik: Seems らしい/よう/そう, Hearsay, Expectations, Intentions, Supposition} \\
    Melihat kue stroberi di etalase kafe? Kamu bisa bilang \textbf{「美味しそう！」 (Oishisou! - Kelihatannya enak!)}. 
    \begin{itemize}
        \item \textbf{〜そう (〜sou)} = Kelihatannya... (Dari apa yang kita lihat).
        \item \textbf{〜らしい (〜rashii) / 〜そうです (〜sou desu)} = Katanya / Kudengar... (Dari gosip atau berita).
        \item \textbf{〜つもり (〜tsumori)} = Rencana / Niat.
    \end{itemize}
    \textbf{Tips Turis:} Kalau besok ada rencana ke Gunung Fuji, kamu bisa bilang \textit{"Ashita, Fuji-san ni iku \textbf{tsumori} desu!"} (Besok, rencananya mau pergi ke Gunung Fuji!).
\end{lessonbox}

\begin{convobox}{11}
    \noindent
    \jpkana{この}{Kono} \jpkana{ケーキ}{keeki}\jpkana{、}{,} 
    \jpkana{おいし}{oishi}\jpkana{そうですね！}{sou desu ne!}

    \vspace{1.5em}
    \noindent{\Large \color{articolor} \textbf{Arti:} Kue ini, kelihatan enak ya!}
\end{convobox}

\begin{convobox}{12}
    \noindent
    \jpword{あした}{明日}{Ashita} \jpkana{は}{wa} 
    \jpword{ゆき}{雪}{yuki} \jpkana{が}{ga} 
    \jpword{ふ}{降}{fu}\jpkana{る}{ru} 
    \jpkana{そうです。}{sou desu.}

    \vspace{1.5em}
    \noindent{\Large \color{articolor} \textbf{Arti:} KATANYA besok akan turun salju (Dengar dari berita cuaca).}
\end{convobox}

% ================= THEME 8 =================
\begin{lessonbox}{8. Pengandaian \& Keputusan (Conditionals)}
    \textbf{Topik: Conditional たら/ば, ことにする/なる, ように} \\
    Bagaimana jika hujan? Bagaimana jika tiketnya habis? Gunakan \textbf{〜たら (〜tara)} atau \textbf{〜ば (〜ba)} untuk "Jika/Kalau".
    \begin{itemize}
        \item \textbf{〜たら (〜tara)} = Jika (Paling sering dipakai dan paling mudah!). \textit{Ame ga futtara...} (Kalau hujan...).
        \item \textbf{〜ことにする (〜koto ni suru)} = Saya memutuskan untuk... \textit{Kore o kau koto ni shimasu!} (Saya putuskan beli yang ini!).
    \end{itemize}
\end{lessonbox}

\begin{convobox}{13}
    \noindent
    \jpword{にほん}{日本}{Nihon} \jpkana{に}{ni} 
    \jpword{つ}{着}{tsu}\jpkana{いたら、}{itara,} 
    \jpkana{ラーメン}{raamen} \jpkana{を}{o} 
    \jpword{た}{食}{ta}\jpkana{べましょう！}{bemashou!}

    \vspace{1.5em}
    \noindent{\Large \color{articolor} \textbf{Arti:} JIKA (nanti) sudah tiba di Jepang, mari kita makan ramen!}
\end{convobox}

% ================= THEME 9 =================
\begin{lessonbox}{9. Memulai, Selesai, \& Kondisi (Actions in Progress)}
    \textbf{Topik: ところ, Completing (てしまう), Starting (はじめる), Verb States, Volitional Expanded, Advanced Questions, Comparisons 2} \\
    Terakhir, mari mengekspresikan tindakanmu dengan lebih tajam!
    \begin{itemize}
        \item \textbf{〜ところ (〜tokoro)} = Baru saja mau / Sedang / Baru saja selesai. \textit{Ima, iku tokoro desu!} (Sekarang, aku baru saja mau berangkat!).
        \item \textbf{〜てしまう (〜te shimau)} = Selesai (tuntas), atau "Waduh/Gawat!" (Tidak sengaja).
        \item \textbf{〜はじめる (〜hajimeru)} = Mulai melakukan sesuatu. \textit{Ame ga furi-hajimemashita} (Hujan mulai turun).
    \end{itemize}
    \textbf{Tips Turis Darurat:} Kalau kamu tidak sengaja menjatuhkan atau menghilangkan paspor/tiket, gunakan \textit{te shimaimashita} untuk menunjukkan rasa menyesal/gawat. \textit{"Kippu o nakushite shimaimashita!"} (Waduh, tiket saya hilang!). Petugas akan sangat mengerti kepanikanmu!
\end{lessonbox}

\begin{convobox}{14}
    \noindent
    \jpword{きっぷ}{切符}{Kippu} \jpkana{を}{o} 
    \jpword{な}{失}{na}\jpkana{くして}{kushite} 
    \jpkana{しまいました。}{shimaimashita.}

    \vspace{1.5em}
    \noindent{\Large \color{articolor} \textbf{Arti:} Waduh, saya (tidak sengaja) menghilangkan tiket saya.}
\end{convobox}

\begin{convobox}{15}
    \noindent
    \jpword{いま}{今}{Ima}\jpkana{、}{,} 
    \jpword{ひこうき}{飛行機}{hikouki} \jpkana{に}{ni} 
    \jpword{の}{乗}{no}\jpkana{る}{ru} 
    \jpkana{ところです！}{tokoro desu!}

    \vspace{1.5em}
    \noindent{\Large \color{articolor} \textbf{Arti:} Sekarang, saya baru saja mau naik ke pesawat!}
\end{convobox}

% --- SUMMARY BOX CLOSING ---
\vspace{2em}
\begin{tcolorbox}[
    colback=red!5!white,
    colframe=red!60!black,
    boxrule=3pt,
    arc=5mm,
    title={\huge\bfseries \sffamily 🏆 LULUS TINGKAT MASTER! 🏆},
    fonttitle=\color{white},
    attach boxed title to top center={yshift=-4mm},
    boxed title style={colback=red!60!black, rounded corners},
    left=5mm, right=5mm, top=7mm, bottom=5mm
]
\Large \textbf{Luar Biasa!} Ayah, Bunda, dan Adik-adik kini sudah memegang "Sabuk Hitam" untuk bertahan hidup di Jepang! Ke-50 topik tingkat lanjut ini memang menantang, tapi dengan mengingat tips-tips kuncinya:
\begin{itemize}
    \item Dengarkan kata kuncinya (Jangan panik jika mendengar \textit{Keigo} / bahasa dewa kasir).
    \item Gunakan \textbf{〜んです (n desu)} jika butuh menjelaskan alasan darurat.
    \item Gunakan \textbf{〜てしまいました (te shimaimashita)} jika tidak sengaja melakukan kesalahan.
\end{itemize}
Sekarang, liburan kalian di Jepang pasti akan penuh dengan petualangan bahasa yang luar biasa. \textbf{ご旅行、楽しんでくださいね！ (Go-ryokou, tanoshinde kudasai ne - Nikmati liburannya ya!)}
\end{tcolorbox}

\newpage
\chapter{Ensiklopedia Master Bahasa Jepang (Final)}

% ================= THEME 1 =================
\begin{lessonbox}{1. Pengalaman Pertama \& Mengubah Kata Kerja Jadi Benda}
    \textbf{Topik: Using 初めて, Nominalization, Noun Reference (の), Qualifiers, Noun Suffixes (\#1-\#3, に)} \\
    Saat liburan ke luar negeri, banyak sekali "Pengalaman Pertama" yang akan kalian hadapi. Kata \textbf{初めて (Hajimete - Pertama kali)} adalah senjata rahasia turis!
    \begin{itemize}
        \item \textbf{Hajimete desu:} Kalau pelayan restoran bertanya apakah kamu sudah pernah ke sini, jawab saja \textit{"Hajimete desu"}. Mereka pasti akan menjelaskan cara pesannya dengan sangat ramah!
        \item \textbf{Nominalization (Menjadikan Benda):} Kalau kamu mau bilang "Makan sushi itu menyenangkan", kamu harus mengubah kata kerja "Makan" jadi benda dengan menambah \textbf{の (no)} atau \textbf{こと (koto)}. Contoh: \textit{Sushi o taberu \textbf{no} wa tanoshii desu!}
        \item \textbf{Noun Suffixes に (ni):} Partikel \textit{ni} juga sangat hebat. Bisa dipakai untuk tujuan: \textit{Kaimono \textbf{ni} ikimasu} (Pergi \textbf{untuk} belanja).
    \end{itemize}
\end{lessonbox}

\begin{convobox}{1}
    \noindent
    \jpword{にほん}{日本}{Nihon} \jpkana{は}{wa} 
    \jpword{はじ}{初}{haji}\jpkana{めてです。}{mete desu.}

    \vspace{1.5em}
    \noindent{\Large \color{articolor} \textbf{Arti:} Ini pertama kalinya (saya datang) ke Jepang.}
\end{convobox}

\begin{convobox}{2}
    \noindent
    \jpword{ふじさん}{富士山}{Fuji-san} \jpkana{に}{ni} 
    \jpword{のぼ}{登}{nobo}\jpkana{る}{ru} 
    \jpkana{の}{no} \jpkana{は}{wa} 
    \jpword{たの}{楽}{tano}\jpkana{しいです。}{shii desu.}

    \vspace{1.5em}
    \noindent{\Large \color{articolor} \textbf{Arti:} Mendaki Gunung Fuji ITU menyenangkan.}
\end{convobox}

% ================= THEME 2 =================
\begin{lessonbox}{2. Menjelaskan Alasan dengan Sangat Lancar}
    \textbf{Topik: Sentence Links (Because 1-3), わけ (Wake), Explanations, Giving Information} \\
    Kenapa keretanya terlambat? Kenapa makanannya enak? Orang Jepang sangat suka menjelaskan alasan!
    \begin{itemize}
        \item \textbf{〜から (〜kara) / 〜ので (〜node):} Ini adalah kata "Karena" yang paling standar. \textit{Ame ga futta node...} (Karena turun hujan...).
        \item \textbf{わけ (Wake):} Pernahkah makan ramen yang sangat enak sampai kamu berpikir "Pantas saja enak, ternyata kedai legendaris!"? Dalam bahasa Jepang, itu adalah \textbf{美味しいわけです (Oishii wake desu - Pantas saja enak! / Masuk akal kalau ini enak!)}. Ini akan membuat teman Jepangmu kagum!
    \end{itemize}
\end{lessonbox}

\begin{convobox}{3}
    \noindent
    \jpword{ゆき}{雪}{Yuki} \jpkana{が}{ga} 
    \jpword{ふ}{降}{fu}\jpkana{っている}{tte iru} 
    \jpword{わけ}{訳}{wake} \jpkana{ですね。}{desu ne.}

    \vspace{1.5em}
    \noindent{\Large \color{articolor} \textbf{Arti:} Pantas saja (masuk akal jika) turun salju ya. (Pantas saja dingin banget!)}
\end{convobox}

% ================= THEME 3 =================
\begin{lessonbox}{3. Memberi Penekanan, Contoh, dan Penyangkalan Total!}
    \textbf{Topik: Emphasis (こと/Other), Negation (~ない), も～ない, Negative Verb Endings, など・なんて・なんか, Giving Examples} \\
    Bagaimana cara bilang "Tidak ada SATU PUN"? Atau menyebutkan contoh oleh-oleh?
    \begin{itemize}
        \item \textbf{も〜ない (Mo...nai) / Penyangkalan Total:} Jika kamu ditanya apakah ada barang yang tertinggal, dan kamu yakin tidak ada, katakan \textbf{一つもないです (Hitotsu mo nai desu - Satu pun tidak ada)}.
        \item \textbf{など (Nado - Dan lain-lain):} Saat membeli banyak hal. \textit{Kashi ya, o-cha \textbf{nado} o kaimashita} (Saya membeli snack, teh, dan lain-lain).
        \item \textbf{なんか (Nanka):} Versi sangat kasual dari \textit{Nado}. Sering dipakai anak muda! \textit{Sushi nanka tabetai na~} (Pengen makan sushi atau semacamnya nih~).
    \end{itemize}
\end{lessonbox}

\begin{convobox}{4}
    \noindent
    \jpword{おかね}{お金}{O-kane} \jpkana{が}{ga} 
    \jpword{いちえん}{1円}{ichi-en} \jpkana{も}{mo} 
    \jpkana{ないです。}{nai desu.}

    \vspace{1.5em}
    \noindent{\Large \color{articolor} \textbf{Arti:} Tidak ada uang 1 yen pun (Sama sekali tidak punya uang).}
\end{convobox}

% ================= THEME 4 =================
\begin{lessonbox}{4. Kuantitas, Perbandingan, dan Derajat (Level)}
    \textbf{Topik: ほど expanded, Degree/Amount, Comparisons, (Not) Only} \\
    Adik-adik pasti suka membandingkan mainan atau makanan. Mari kita pelajari levelnya!
    \begin{itemize}
        \item \textbf{〜ほど〜ない (〜hodo...nai):} Artinya "Tidak se-...". Contoh: \textit{Oosaka wa, Toukyou \textbf{hodo} samuku nai desu} (Osaka, TIDAK SE-dingin Tokyo).
        \item \textbf{〜だけ (〜dake) / 〜しか〜ない (〜shika...nai):} \textit{Dake} artinya "Hanya". Tapi kalau mau lebih dramatis, pakai \textit{Shika + Tidak}. Contoh: \textit{Hyaku-en \textbf{shika nai} desu} (Hanya ada 100 yen saja / Tidak ada selain 100 yen).
    \end{itemize}
\end{lessonbox}

\begin{convobox}{5}
    \noindent
    \jpword{きょう}{今日}{Kyou} \jpkana{は}{wa} 
    \jpword{きのう}{昨日}{kinou} \jpkana{ほど}{hodo} 
    \jpword{あつ}{暑}{atsu}\jpkana{くないです。}{kunai desu.}

    \vspace{1.5em}
    \noindent{\Large \color{articolor} \textbf{Arti:} Hari ini tidak SE-panas kemarin.}
\end{convobox}

% ================= THEME 5 =================
\begin{lessonbox}{5. Permintaan Tingkat Lanjut \& Kebutuhan Mendesak}
    \textbf{Topik: Commands/Requests 1-4, Necessity} \\
    Kadang di bandara atau hotel, kamu harus mengisi formulir atau diminta melakukan sesuatu oleh petugas.
    \begin{itemize}
        \item \textbf{〜てください (〜te kudasai):} Tolong lakukan. (Bentuk standar).
        \item \textbf{〜ないでください (〜nai de kudasai):} Tolong JANGAN lakukan! \textit{Koko ni hairanai de kudasai!} (Tolong jangan masuk ke sini!).
        \item \textbf{Necessity (Harus):} Jika petugas bilang \textit{Kakanakereba narimasen}, artinya kamu "Harus / Wajib" menulis sesuatu di formulir tersebut!
    \end{itemize}
\end{lessonbox}

\begin{convobox}{6}
    \noindent
    \jpkana{ここで}{Koko de} 
    \jpword{た}{食}{ta}\jpkana{べないで}{benai de} 
    \jpword{くだ}{下}{kuda}\jpkana{さい。}{sai.}

    \vspace{1.5em}
    \noindent{\Large \color{articolor} \textbf{Arti:} Tolong JANGAN makan di sini.}
\end{convobox}

% ================= THEME 6 =================
\begin{lessonbox}{6. Harapan, Doa, dan Berandai-andai (If / Wishes)}
    \textbf{Topik: Wishes \& Hopes 1-2, Using ように 1-2, Using もし} \\
    Saat kamu berdoa di Kuil Jepang (Jinja), bagaimana cara meminta permohonan dalam bahasa Jepang?
    \begin{itemize}
        \item \textbf{〜ように (〜you ni):} Digunakan untuk berdoa! Gabungkan dengan tepuk tangan 2 kali di kuil dan ucapkan: \textit{Nihongo ga jouzu ni narimasu \textbf{you ni}!} (Semoga bahasa Jepangku jadi pintar!).
        \item \textbf{もし (Moshi):} Artinya "Seandainya/Misalkan". \textit{\textbf{Moshi} jikan ga attara, asobimashou} (Seandainya ada waktu, ayo kita main!).
    \end{itemize}
\end{lessonbox}

\begin{convobox}{7}
    \noindent
    \jpword{あした}{明日}{Ashita} 
    \jpword{は}{晴}{ha}\jpkana{れます}{remasu} 
    \jpkana{ように。}{you ni.}

    \vspace{1.5em}
    \noindent{\Large \color{articolor} \textbf{Arti:} Semoga besok cuacanya cerah.}
\end{convobox}

% ================= THEME 7 =================
% \begin{lessonbox}{7. Memanfaatkan Waktu: "Mumpung" \& Penyelesaian}
%     \textbf{Topik: Actions \& Time (うち, Endings, #3), Continuing Actions, Completing Actions 1-2} \\
%     Waktu liburan di Jepang sangat berharga! Mari kita manfaatkan "Mumpung"!
%     \begin{itemize}
%         \item \textbf{〜うちに (〜uchi ni):} Artinya "Selagi / Mumpung". \textit{Nihon ni iru \textbf{uchi ni}, takusan tabetai!} (Mumpung/Selagi masih ada di Jepang, aku mau makan banyak!).
%         \item \textbf{Completing Actions (〜てしまう / 〜てある):} \textit{Te shimau} (Sudah selesai semua / Tidak sengaja). Sedangkan \textit{Te aru} artinya suatu tindakan sudah disiapkan. \textit{Hoteru wa yoyaku shi\textbf{te arimasu}} (Hotelnya sudah dalam keadaan selesai dipesan).
%     \end{itemize}
% \end{lessonbox}

\begin{convobox}{8}
    \noindent
    \jpword{あたた}{温}{Atata}\jpkana{かい}{kai} 
    \jpword{うち}{内}{uchi} \jpkana{に}{ni} 
    \jpword{た}{食}{ta}\jpkana{べて}{bete} 
    \jpword{くだ}{下}{kuda}\jpkana{さい。}{sai.}

    \vspace{1.5em}
    \noindent{\Large \color{articolor} \textbf{Arti:} Tolong makan SELAGI (mumpung) masih hangat.}
\end{convobox}

% ================= THEME 8 =================
\begin{lessonbox}{8. Meskipun, Padahal, dan Sebagai Gantinya (Konjungsi Ajaib)}
    \textbf{Topik: Conjunctions でも 1-2, Although, Even Though, Instead, Considering} \\
    Ini adalah tata bahasa untuk anak pintar yang suka berdebat dengan alasan yang logis!
    \begin{itemize}
        \item \textbf{〜のに (〜noni):} Artinya "Padahal". Mengandung perasaan kecewa. \textit{Takai \textbf{noni}, oishikunai!} (PADAHAL harganya mahal, tapi tidak enak!).
        \item \textbf{〜かわりに (〜kawari ni):} Artinya "Sebagai gantinya / Alih-alih". \textit{Kyou wa sushi no \textbf{kawari ni}, ramen o tabemasu} (Hari ini sebagai ganti sushi, kita makan ramen).
    \end{itemize}
\end{lessonbox}

\begin{convobox}{9}
    \noindent
    \jpword{あめ}{雨}{Ame} \jpkana{が}{ga} 
    \jpword{ふ}{降}{fu}\jpkana{っている}{tte iru} 
    \jpkana{のに、}{noni,} 
    \jpkana{かさを}{kasa o} 
    \jpword{も}{持}{mo}\jpkana{っていません。}{tte imasen.}

    \vspace{1.5em}
    \noindent{\Large \color{articolor} \textbf{Arti:} PADAHAL sedang turun hujan, (tapi) saya tidak membawa payung.}
\end{convobox}

% ================= THEME 9 =================
\begin{lessonbox}{9. Melihat Sifat, Gosip, dan Kesan Orang Lain}
    \textbf{Topik: Using っぽい, Using ～がる, Seeming/Pretending, Reporting Information 1-2, Expectations} \\
    Menebak-nebak sifat barang atau perasaan orang Jepang lain di sekitarmu:
    \begin{itemize}
        \item \textbf{〜っぽい (〜ppoi):} Kelihatan seperti / Mirip seperti. \textit{Kodomo-ppoi} (Sifatnya seperti anak-anak).
        \item \textbf{〜がる (〜garu):} Kita tidak bisa menebak perasaan orang 100\%. Kalau ada temanmu yang menggigil, jangan bilang "Dia dingin", tapi bilangnya \textit{Samu\textbf{gatte iru}} (Dia KELIHATANNYA merasa dingin).
        \item \textbf{〜はずです (〜hazu desu):} Ekspektasi / Seharusnya. \textit{Densha wa 10-ji ni kuru \textbf{hazu desu}} (SEHARUSNYA keretanya datang jam 10).
    \end{itemize}
\end{lessonbox}

\begin{convobox}{10}
    \noindent
    \jpkana{あの}{Ano} \jpword{ひと}{人}{hito} \jpkana{は}{wa} 
    \jpword{さむ}{寒}{samu}\jpkana{がっています。}{gatte imasu.}

    \vspace{1.5em}
    \noindent{\Large \color{articolor} \textbf{Arti:} Orang itu (kelihatannya) sedang merasa kedinginan.}
\end{convobox}

\begin{convobox}{11}
    \noindent
    \jpword{あした}{明日}{Ashita} \jpkana{の}{no} 
    \jpword{てんき}{天気}{tenki} \jpkana{は}{wa} 
    \jpword{は}{晴}{ha}\jpkana{れる}{reru} 
    \jpword{はず}{筈}{hazu} \jpkana{です。}{desu.}

    \vspace{1.5em}
    \noindent{\Large \color{articolor} \textbf{Arti:} Cuaca besok SEHARUSNYA (pasti) cerah.}
\end{convobox}

% ================= THEME 10 =================
\begin{lessonbox}{10. Bahasa Gaul Tingkat Akhir (Informal Endings \& Particles Expanded)}
    \textbf{Topik: Informal Sentence Endings, More Questions, Particles Expanded} \\
    Tingkat tertinggi belajar bahasa Jepang adalah berbicara seperti karakter anime atau teman akrab yang sedang bersantai!
    \begin{itemize}
        \item \textbf{〜かな (〜kana):} Diucapkan saat berbicara sendiri atau ragu. \textit{Doko ni ikou kana~?} (Kira-kira kita pergi ke mana ya~?). Sangat imut untuk anak-anak!
        \item \textbf{〜さ (〜sa) / 〜ぜ (〜ze) / 〜わ (〜wa):} Partikel di akhir kalimat untuk gaya bicara tertentu. \textit{Ze} (Kasar/Laki-laki banget), \textit{Wa} (Feminin/Anggun), \textit{Sa} (Santai).
        \item \textbf{Expanded Particles:} Menggabungkan partikel! \textit{Tomodachi \textbf{to no} shashin} (Foto BERSAMA teman). \textit{Nihon \textbf{kara no} omiyage} (Oleh-oleh DARI Jepang).
    \end{itemize}
\end{lessonbox}

\begin{convobox}{12}
    \noindent
    \jpword{きょう}{今日}{Kyou} \jpkana{は}{wa} 
    \jpword{なに}{何}{nani} \jpkana{を}{o} 
    \jpword{た}{食}{ta}\jpkana{べよう}{beyou} \jpkana{かな〜。}{kana~.}

    \vspace{1.5em}
    \noindent{\Large \color{articolor} \textbf{Arti:} Hari ini, kira-kira enaknya makan apa ya~.}
\end{convobox}

% --- SUMMARY BOX CLOSING ---
\vspace{2em}
\begin{tcolorbox}[
    colback=yellow!5!white,
    colframe=orange!80!black,
    boxrule=4pt,
    arc=6mm,
    title={\huge\bfseries \sffamily 🎆 T A M A T - KELULUSAN MASTER BAHASA JEPANG! 🎆},
    fonttitle=\color{white},
    attach boxed title to top center={yshift=-5mm},
    boxed title style={colback=orange!80!black, rounded corners},
    left=5mm, right=5mm, top=7mm, bottom=5mm
]
\Large \textbf{Air Mata Kebahagiaan! Kalian Berhasil Menyelesaikannya!}\\[0.5em]
Ayah, Bunda, dan Adik-adik yang super hebat, dengan selesainya Bab Puncak (Revisit 4) ini, \textbf{seluruh materi di dalam Andi - Japanese Handbook telah tamat 100\%!} 

Dari tidak bisa mengucapkan sepatah kata pun, hingga bisa memesan Ramen di mesin tiket, mengobrol dengan kasir Konbini, menebak perasaan orang dengan \textit{\textasciitilde garu}, hingga mendoakan cuaca di kuil dengan \textit{\textasciitilde you ni}. Ini adalah pencapaian yang luar biasa!

\textbf{Pesan Terakhir untuk Turis Hebat:}
\begin{itemize}
    \item Jangan takut berbuat salah saat berbicara di Jepang. Orang Jepang sangat menghargai senyuman dan usahamu!
    \item Bawalah buku ini kemanapun kalian pergi di Jepang. Jadikan buku ini saksi petualangan keluarga kalian yang luar biasa.
\end{itemize}

\begin{center}
    \Huge \textbf{いってらっしゃい！}\\[0.2em]
    (Itterasshai!)\\[0.2em]
    \Large Selamat Berpetualang di Jepang!
\end{center}
\end{tcolorbox}

\newpage



\end{document}